\subsection{The explicit threefold with identity monodromy}

Throughout this part we orient complex spaces by their complex structures.
All homology groups in the example have coefficients in $\Z$. A meridian is
counterclockwise in its indicated holomorphic coordinate; at infinity this
coordinate is $1/t$. We write
\[
 E=\C/(\Z+i\Z)=\C^*/q^\Z,\qquad q=e^{-2\pi},\qquad
 A=E^2,\qquad \Lambda=\Z e_1\oplus\Z e_2\oplus\Z e_3\oplus\Z e_4,
\]
where $(e_1,e_2,e_3,e_4)=((1,0),(i,0),(0,1),(0,i))$. The comparison
$E\longrightarrow\C^*/q^\Z$ is $z\mapsto e^{2\pi iz}$.

\begin{statement}{P4.1}{Proved}
Let $V=\operatorname{Bl}_{(0,0)}(B\times\Delta)$, where $B=E$, let $R$ be
the exceptional curve, and put $L=\OO_V(R)$. The holomorphic quotient
$U=L^\times/q^\Z$ is a smooth Hausdorff threefold, proper over $\Delta$.
Its noncentral fibres are $E\times B$. Its central fibre is the reduced
union of a primary Kodaira surface and the scalar Hopf surface
$(\C^2\setminus\{0\})/q^\Z$, meeting transversally in one elliptic curve.
\end{statement}
\begin{proof}
We use coordinates $(z,t)$ at the blown-up point. The blowup has charts
\[
 (z,t)=(u,uv),\qquad (z,t)=(rs,s),\qquad r=v^{-1},\quad s=uv.
\]
The frames $u^{-1}$ and $s^{-1}$ for $L$ give fibre coordinates related by
$\lambda_B=v\lambda_A$. Each nonzero fibre coordinate, modulo multiplication
by $q$, gives a product chart with fibre $E$. The transitions are holomorphic
translations on $E$. The scalar action is free and properly discontinuous:
the logarithm of the norm of a nonzero fibre coordinate changes by the
nonzero constant $\log q$, and compact subsets meet only finitely many
translates. Thus the quotient is Hausdorff and smooth. These product charts
have compact fibre $E$, so $U\to V$ is proper. The blowup and $B\times\Delta
\to\Delta$ are proper, proving properness of their composite.

The canonical section of $\OO(R)$ has coefficients $u$ and $s$. Over
$t\ne0$ it never vanishes. The ratios $\lambda_A/u=\lambda_B/s$ therefore
give the specified holomorphic product trivialization over $\Delta^*$.

The central divisor in $V$ is $B'\cup R$, with $R^2=-1$ and $R\cdot B'=1$.
Thus $L|_R=\OO_{\Pp^1}(-1)$ and $L|_{B'}$ has degree one. Over $B'$ the
quotient is a principal elliptic bundle with primitive nonzero circle Chern
class. Its canonical bundle is trivial: the base differential wedged with
$d\lambda/\lambda$ is invariant under the line-bundle transition, because
the additional term is pulled back from the one-dimensional base. The
circle-bundle calculation in P4.3 gives $b_1=3$. The surface classification
BHPV, Theorem VI.1.1 and Table 10, consequently identifies this surface as
a primary Kodaira surface.

The tautological evaluation morphism
$\OO_{\Pp^1}(-1)^\times\to\C^2\setminus\{0\}$ is a biholomorphism: a
nonzero vector uniquely determines its line. It intertwines both scalar
actions, and hence identifies the other component with the stated Hopf
surface. Over $B'\cap R$ the components meet in $E$. In local product
coordinates the map is $(u,v,c)\mapsto uv$; this proves transversality,
multiplicity one, and the absence of other critical points.
\end{proof}

\begin{statement}{P4.2}{Proved}
The three markings and the principal-bundle clutch in the prompt define a
smooth compact complex threefold $f:Y\to\Pp^1$ with connected fibres,
exactly three singular fibres of the form in P4.1, and identity monodromy
on $H_1$ of the regular fibre.
\end{statement}
\begin{proof}
The map $\ell(t)=(\log t)/(2\pi i)$ with values in $E$ is single-valued and
holomorphic on an annulus: changing the logarithm adds an integer.
Consequently $a_{\rm out}=a_{\rm in}+(\ell(t),0)$ is a holomorphic transition
for a principal $A$-bundle on $\Pp^1$. Product charts with compact fibre
prove that this bundle is Hausdorff and compact.

We choose disjoint discs about $0,1,\infty$ and replace their product models
by P4.1 using the holomorphic group isomorphisms
\[
 C_i\times B_i\longrightarrow A:\qquad
 (z,w)\mapsto(w,z),\quad(iz,iz+w),\quad(z,iz+w),
\]
respectively. Each lattice matrix $S_i$ displayed below has determinant
one. These maps, $\lambda_B=v\lambda_A$, the canonical punctured-disc
trivialization in P4.1, and the clutch $\ell$ specify all transitions.
Different base points separate in the base. Points over the same base
point lie in one Hausdorff local model, which separates them. The gluing is
therefore Hausdorff. Properness holds over each base chart and hence over
the base; compactness follows.

The two central components have connected elliptic fibres over connected
curves and meet each other. All remaining fibres equal $A$. The equations
$uv=t$ give exactly the indicated critical values. The punctured-disc
trivializations identify all fibre lattices with $\Lambda$, and the remaining
transition is a translation on $A$. Translation induces the identity on
$H_1(A;\Z)$, giving the asserted monodromy representation.
\end{proof}

\begin{statement}{P4.3}{Proved}
Order the local fibre lattice as $(a,c,b,d)$, where $(a,c)$ are the cycles
$(1,i)$ in $C_i$. The regular-fibre inclusions induce the quotient maps
\[
 H_1(A)=\Lambda\longrightarrow\Lambda/\Z a=H_1(U),\qquad
 H_2(A)=\bigwedge^2\Lambda\longrightarrow
 \bigwedge^2\Lambda/\Z(a\wedge c)=H_2(U).
\]
Both target groups are free. The meridian and every product of the
meridian with a fibre one-cycle map to zero in these degrees.
\end{statement}
\begin{proof}
A Hermitian metric gives a diffeomorphism
$L^\times/q^\Z\cong S(L)\times S^1_c$, with factors the circle bundle and
the logarithm of the norm modulo $\log q$. We first specify a homotopy
model for $V$. Blowing up a four-ball replaces it by the disc bundle of
$\OO(-1)$ along its Hopf boundary. The complement of the ball in
$B\times\Delta$ has homotopy type $B\vee S^3$; the boundary sphere is
the $S^3$ summand, since its map to $B$ is nullhomotopic. The resulting
homotopy pushout is $B\vee S^2$. Equivalently, on simply connected
covers the replacement adds one sphere at each lift of the centre and
no other homology; the comparison is an isomorphism on $\pi_1$ and on
the homology of those covers, and hence a homotopy equivalence. We may
choose the $B$ summand to be a regular $B_t$ and the sphere to be $R$,
joining their base points by a path.

The circle Chern class is zero on $B_t$ and $-1$ on $R$. Therefore
$S(L)$ is represented by the homotopy pushout
\[
 T^3_{(a,b,d)}\longleftarrow S^1_a\longrightarrow S^3.
\]
The second map is a Hopf fibre and is nullhomotopic. Thus this pushout
has the homotopy type $(T^3\cup_aD^2)\vee S^3$. We multiply by $S^1_c$.
The resulting integral cellular complex starts with the exterior-cell
complex of $T^4$, ordered by $(a,c,b,d)$, and adds cells $r$ in degree
two, $rc,h$ in degree three, and $hc$ in degree four. Its only nonzero
differentials are
\[
 \partial r=a,\qquad \partial(rc)=a\wedge c.
\]
We choose the orientation of $r$ to give the first equality and use the
product orientation for $rc$. The complete differential matrices are
printed below. The complex has free homology of ranks $(1,3,5,5,2)$;
the inclusion of the torus subcomplex gives exactly the two quotient maps.

We represent a base point, and both base generators $b,d$, away from
$0\in B$. Their products with the disc miss $R$, so the canonical section
trivializes $L$ over the resulting chains. The meridian and the products
with $b,d$ consequently bound. The products with $a,c$ bound over one
such base point. This proves all the remaining assertions on the maps.
\end{proof}

The lattice matrices, with columns $(a_i,c_i,b_i,d_i)$ in the common basis,
are
\[
S_1=\begin{pmatrix}0&0&1&0\\0&0&0&1\\1&0&0&0\\0&1&0&0\end{pmatrix},\quad
S_2=\begin{pmatrix}0&-1&0&0\\1&0&0&0\\0&-1&1&0\\1&0&0&1\end{pmatrix},\quad
S_3=\begin{pmatrix}1&0&0&0\\0&1&0&0\\0&-1&1&0\\1&0&0&1\end{pmatrix}.
\]
Put $Q_i=S_i^{-1}$. Deleting the first row of $Q_i$ gives $R_{1,i}$.
The row indices of $\bigwedge^2Q_i$ are $(ac,ab,ad,cb,cd,bd)$ and its
column indices are $(12,13,14,23,24,34)$; deleting its first row gives
$R_{2,i}$. Thus these are the integer matrices of the local quotient maps,
with target bases $(c_i,b_i,d_i)$ and $(a_ib_i,a_id_i,c_ib_i,c_id_i,b_id_i)$.

\begin{statement}{P4.4}{Proved}
The threefold $Y$ has $H_1(Y;\Z)=H_2(Y;\Z)=0$ and $\pi_1(Y)=1$.
\end{statement}
\begin{proof}
Let $P$ be the complement of three small open base discs. The remaining
bundle $W=f^{-1}(P)$ is trivial. We use $W=A\times P$, let $s_1,s_2$
generate $H_1(P)$, and set $s_3=-s_1-s_2$. The overlaps are
$N_i=A\times S^1$. Positive continuation of $\ell$ around zero adds $e_1$.
With
\[
 g_1=e_1,\quad g_2=g_3=0,\qquad
 v_1=\binom10,\quad v_2=\binom01,\quad v_3=\binom{-1}{-1},
\]
the overlap maps into $W$ are
\[
 x\mapsto x,\quad s_i\mapsto g_i+v_i,\qquad
 x\wedge y\mapsto x\wedge y,\quad
 x\wedge s_i\mapsto x\wedge g_i+x\otimes v_i.
\]
These are the maps represented by the products of the cycles, or by the
induced maps into the aspherical space $A\times P$; in particular they
retain the degree-two contribution of the clutch.

We use the ordered basis $(e_1,e_2,e_3,e_4,s_1,s_2)$ for $H_1(W)$ and
append the three bases $(c_i,b_i,d_i)$ for the local groups. On the source
we concatenate the three bases $(e_1,e_2,e_3,e_4,s_i)$. The resulting
$15\times15$ Mayer--Vietoris matrix $F_1$ has, in its $i$th five-column
block, upper block
\[
 \begin{pmatrix}I_4&g_i\\0_{2\times4}&v_i\end{pmatrix}
\]
and local block $(-R_{1,i},0)$ in the $i$th summand; its other local
blocks vanish. The sign is the usual difference of the two inclusions.

For $H_2(W)$ we use
\[
 (12,13,14,23,24,34,e_1s_1,\ldots,e_4s_1,e_1s_2,\ldots,e_4s_2)
\]
and append the three local five-element bases. The source is the
concatenation of $(12,13,14,23,24,34,e_1s_i,\ldots,e_4s_i)$.
The $29\times30$ matrix $F_2$ has top six rows
$[I_6,(e_j\wedge g_i)_{j=1}^4]$ in its $i$th block; its next eight rows
have the coefficients of $e_j\otimes v_i$ in the displayed order, and
its $i$th local block is $[-R_{2,i},0]$. All other blocks vanish.
Finally
\[
 F_0=\begin{pmatrix}1&1&1\\-1&0&0\\0&-1&0\\0&0&-1\end{pmatrix}
\]
has zero kernel. Thus the required exact sequence is
\[
 \Z^{30}\xrightarrow{F_2}\Z^{29}\longrightarrow H_2(Y)
 \longrightarrow\Z^{15}\xrightarrow{F_1}\Z^{15}\longrightarrow H_1(Y)
 \longrightarrow\Z^3\xrightarrow{F_0}\Z^4.
\]

We verify integral surjectivity by hand. Eliminating the local $H_1$
groups and two independent base meridians leaves $\Lambda$ modulo
$e_1,a_1,a_2,a_3$. The relation matrix is
\[
 R_{H_1}=\begin{pmatrix}1&0&0&1\\0&0&1&0\\0&1&0&0\\0&0&1&1\end{pmatrix},
 \qquad\det R_{H_1}=-1.
\]
Hence $F_1$ is an integral isomorphism. In degree two the fibre
$\bigwedge^2\Lambda$ columns eliminate the fifteen local generators,
leaving the three relations $\eta_i=a_i\wedge c_i$. The second meridian
columns kill $\Lambda\otimes s_2$, the third then kill
$\Lambda\otimes s_1$, and the first leave $e_1\wedge\Lambda$. The six
relations, in the order $e_{12},e_{13},e_{14},\eta_1,\eta_2,\eta_3$, have
matrix
\[
R_{H_2}=\begin{pmatrix}
1&0&0&0&1&1\\0&1&0&0&0&-1\\0&0&1&0&1&0\\
0&0&0&0&-1&0\\0&0&0&0&0&-1\\0&0&0&1&1&1
\end{pmatrix},\qquad \det R_{H_2}=1.
\]
This proves surjectivity of $F_2$ over $\Z$ and excludes torsion. The
executed independent Smith calculation gives fifteen unit factors for
$F_1$ and twenty-nine unit factors for $F_2$, and computes $\det F_1=-1$.

For the fundamental group the local pushout gives
$\pi_1(U_i)=\Lambda/\Z a_i$, with trivial meridian. Van Kampen starts
from $\pi_1(W)=\Lambda\times F(s_1,s_2)$ and imposes
\[
 a_1=a_2=a_3=0,\qquad s_1+e_1=0,\quad s_2=0,\quad s_1+s_2=0.
\]
Both base generators and $e_1$ vanish. The same unimodular matrix
$R_{H_1}$ then kills the fibre lattice, proving $\pi_1(Y)=1$.
\end{proof}

The program \texttt{code/part4/example\_mayer\_vietoris.py} recomputes all
entries, differentials, determinants and Smith forms exactly. Its output
and machine-readable matrices are supplied alongside it. For completeness
we print every integer matrix, including the zero differentials:
% Generated by code/part4/matrices_to_tex.py from example_matrices.json.
\begingroup
\setlength{\arraycolsep}{1.7pt}
\renewcommand{\arraystretch}{0.88}
\par\smallskip\noindent
$S_1$ has size $4\times 4$; its full integer matrix is
\[
\resizebox{!}{11.2mm}{$\displaystyle S_1=\begin{pmatrix}
0 & 0 & 1 & 0\\
0 & 0 & 0 & 1\\
1 & 0 & 0 & 0\\
0 & 1 & 0 & 0\\
\end{pmatrix}$}
\]
\par\smallskip\noindent
$S_2$ has size $4\times 4$; its full integer matrix is
\[
\resizebox{!}{11.2mm}{$\displaystyle S_2=\begin{pmatrix}
0 & -1 & 0 & 0\\
1 & 0 & 0 & 0\\
0 & -1 & 1 & 0\\
1 & 0 & 0 & 1\\
\end{pmatrix}$}
\]
\par\smallskip\noindent
$S_3$ has size $4\times 4$; its full integer matrix is
\[
\resizebox{!}{11.2mm}{$\displaystyle S_3=\begin{pmatrix}
1 & 0 & 0 & 0\\
0 & 1 & 0 & 0\\
0 & -1 & 1 & 0\\
1 & 0 & 0 & 1\\
\end{pmatrix}$}
\]
\par\smallskip\noindent
$R_{1,1}$ has size $3\times 4$; its full integer matrix is
\[
\resizebox{!}{8.399999999999999mm}{$\displaystyle R_{1,1}=\begin{pmatrix}
0 & 0 & 0 & 1\\
1 & 0 & 0 & 0\\
0 & 1 & 0 & 0\\
\end{pmatrix}$}
\]
\par\smallskip\noindent
$R_{1,2}$ has size $3\times 4$; its full integer matrix is
\[
\resizebox{!}{8.399999999999999mm}{$\displaystyle R_{1,2}=\begin{pmatrix}
-1 & 0 & 0 & 0\\
-1 & 0 & 1 & 0\\
0 & -1 & 0 & 1\\
\end{pmatrix}$}
\]
\par\smallskip\noindent
$R_{1,3}$ has size $3\times 4$; its full integer matrix is
\[
\resizebox{!}{8.399999999999999mm}{$\displaystyle R_{1,3}=\begin{pmatrix}
0 & 1 & 0 & 0\\
0 & 1 & 1 & 0\\
-1 & 0 & 0 & 1\\
\end{pmatrix}$}
\]
\par\smallskip\noindent
$R_{2,1}$ has size $5\times 6$; its full integer matrix is
\[
\resizebox{!}{14.0mm}{$\displaystyle R_{2,1}=\begin{pmatrix}
0 & -1 & 0 & 0 & 0 & 0\\
0 & 0 & 0 & -1 & 0 & 0\\
0 & 0 & -1 & 0 & 0 & 0\\
0 & 0 & 0 & 0 & -1 & 0\\
1 & 0 & 0 & 0 & 0 & 0\\
\end{pmatrix}$}
\]
\par\smallskip\noindent
$R_{2,2}$ has size $5\times 6$; its full integer matrix is
\[
\resizebox{!}{14.0mm}{$\displaystyle R_{2,2}=\begin{pmatrix}
1 & 0 & 0 & 1 & 0 & 0\\
0 & 0 & 0 & 0 & 1 & 0\\
0 & -1 & 0 & 0 & 0 & 0\\
1 & 0 & -1 & 0 & 0 & 0\\
1 & 0 & -1 & 1 & 0 & 1\\
\end{pmatrix}$}
\]
\par\smallskip\noindent
$R_{2,3}$ has size $5\times 6$; its full integer matrix is
\[
\resizebox{!}{14.0mm}{$\displaystyle R_{2,3}=\begin{pmatrix}
1 & 1 & 0 & 0 & 0 & 0\\
0 & 0 & 1 & 0 & 0 & 0\\
0 & 0 & 0 & 1 & 0 & 0\\
1 & 0 & 0 & 0 & 1 & 0\\
1 & 1 & 0 & 0 & 1 & 1\\
\end{pmatrix}$}
\]
\par\smallskip\noindent
$\partial_1$ has size $1\times 4$; its full integer matrix is
\[
\resizebox{!}{2.8mm}{$\displaystyle \partial_1=\begin{pmatrix}
0 & 0 & 0 & 0\\
\end{pmatrix}$}
\]
\par\smallskip\noindent
$\partial_2$ has size $4\times 7$; its full integer matrix is
\[
\resizebox{!}{11.2mm}{$\displaystyle \partial_2=\begin{pmatrix}
0 & 0 & 0 & 0 & 0 & 0 & 1\\
0 & 0 & 0 & 0 & 0 & 0 & 0\\
0 & 0 & 0 & 0 & 0 & 0 & 0\\
0 & 0 & 0 & 0 & 0 & 0 & 0\\
\end{pmatrix}$}
\]
\par\smallskip\noindent
$\partial_3$ has size $7\times 6$; its full integer matrix is
\[
\resizebox{!}{19.599999999999998mm}{$\displaystyle \partial_3=\begin{pmatrix}
0 & 0 & 0 & 0 & 1 & 0\\
0 & 0 & 0 & 0 & 0 & 0\\
0 & 0 & 0 & 0 & 0 & 0\\
0 & 0 & 0 & 0 & 0 & 0\\
0 & 0 & 0 & 0 & 0 & 0\\
0 & 0 & 0 & 0 & 0 & 0\\
0 & 0 & 0 & 0 & 0 & 0\\
\end{pmatrix}$}
\]
\par\smallskip\noindent
$\partial_4$ has size $6\times 2$; its full integer matrix is
\[
\resizebox{!}{16.799999999999997mm}{$\displaystyle \partial_4=\begin{pmatrix}
0 & 0\\
0 & 0\\
0 & 0\\
0 & 0\\
0 & 0\\
0 & 0\\
\end{pmatrix}$}
\]
\par\smallskip\noindent
$F_0$ has size $4\times 3$; its full integer matrix is
\[
\resizebox{!}{11.2mm}{$\displaystyle F_0=\begin{pmatrix}
1 & 1 & 1\\
-1 & 0 & 0\\
0 & -1 & 0\\
0 & 0 & -1\\
\end{pmatrix}$}
\]
\par\smallskip\noindent
$F_1$ has size $15\times 15$; its full integer matrix is
\[
\resizebox{0.94\linewidth}{!}{$\displaystyle F_1=\begin{pmatrix}
1 & 0 & 0 & 0 & 1 & 1 & 0 & 0 & 0 & 0 & 1 & 0 & 0 & 0 & 0\\
0 & 1 & 0 & 0 & 0 & 0 & 1 & 0 & 0 & 0 & 0 & 1 & 0 & 0 & 0\\
0 & 0 & 1 & 0 & 0 & 0 & 0 & 1 & 0 & 0 & 0 & 0 & 1 & 0 & 0\\
0 & 0 & 0 & 1 & 0 & 0 & 0 & 0 & 1 & 0 & 0 & 0 & 0 & 1 & 0\\
0 & 0 & 0 & 0 & 1 & 0 & 0 & 0 & 0 & 0 & 0 & 0 & 0 & 0 & -1\\
0 & 0 & 0 & 0 & 0 & 0 & 0 & 0 & 0 & 1 & 0 & 0 & 0 & 0 & -1\\
0 & 0 & 0 & -1 & 0 & 0 & 0 & 0 & 0 & 0 & 0 & 0 & 0 & 0 & 0\\
-1 & 0 & 0 & 0 & 0 & 0 & 0 & 0 & 0 & 0 & 0 & 0 & 0 & 0 & 0\\
0 & -1 & 0 & 0 & 0 & 0 & 0 & 0 & 0 & 0 & 0 & 0 & 0 & 0 & 0\\
0 & 0 & 0 & 0 & 0 & 1 & 0 & 0 & 0 & 0 & 0 & 0 & 0 & 0 & 0\\
0 & 0 & 0 & 0 & 0 & 1 & 0 & -1 & 0 & 0 & 0 & 0 & 0 & 0 & 0\\
0 & 0 & 0 & 0 & 0 & 0 & 1 & 0 & -1 & 0 & 0 & 0 & 0 & 0 & 0\\
0 & 0 & 0 & 0 & 0 & 0 & 0 & 0 & 0 & 0 & 0 & -1 & 0 & 0 & 0\\
0 & 0 & 0 & 0 & 0 & 0 & 0 & 0 & 0 & 0 & 0 & -1 & -1 & 0 & 0\\
0 & 0 & 0 & 0 & 0 & 0 & 0 & 0 & 0 & 0 & 1 & 0 & 0 & -1 & 0\\
\end{pmatrix}$}
\]
\par\smallskip\noindent
$F_2$ has size $29\times 30$; its full integer matrix is
\[
\resizebox{0.94\linewidth}{!}{$\displaystyle F_2=\begin{pmatrix}
1 & 0 & 0 & 0 & 0 & 0 & 0 & -1 & 0 & 0 & 1 & 0 & 0 & 0 & 0 & 0 & 0 & 0 & 0 & 0 & 1 & 0 & 0 & 0 & 0 & 0 & 0 & 0 & 0 & 0\\
0 & 1 & 0 & 0 & 0 & 0 & 0 & 0 & -1 & 0 & 0 & 1 & 0 & 0 & 0 & 0 & 0 & 0 & 0 & 0 & 0 & 1 & 0 & 0 & 0 & 0 & 0 & 0 & 0 & 0\\
0 & 0 & 1 & 0 & 0 & 0 & 0 & 0 & 0 & -1 & 0 & 0 & 1 & 0 & 0 & 0 & 0 & 0 & 0 & 0 & 0 & 0 & 1 & 0 & 0 & 0 & 0 & 0 & 0 & 0\\
0 & 0 & 0 & 1 & 0 & 0 & 0 & 0 & 0 & 0 & 0 & 0 & 0 & 1 & 0 & 0 & 0 & 0 & 0 & 0 & 0 & 0 & 0 & 1 & 0 & 0 & 0 & 0 & 0 & 0\\
0 & 0 & 0 & 0 & 1 & 0 & 0 & 0 & 0 & 0 & 0 & 0 & 0 & 0 & 1 & 0 & 0 & 0 & 0 & 0 & 0 & 0 & 0 & 0 & 1 & 0 & 0 & 0 & 0 & 0\\
0 & 0 & 0 & 0 & 0 & 1 & 0 & 0 & 0 & 0 & 0 & 0 & 0 & 0 & 0 & 1 & 0 & 0 & 0 & 0 & 0 & 0 & 0 & 0 & 0 & 1 & 0 & 0 & 0 & 0\\
0 & 0 & 0 & 0 & 0 & 0 & 1 & 0 & 0 & 0 & 0 & 0 & 0 & 0 & 0 & 0 & 0 & 0 & 0 & 0 & 0 & 0 & 0 & 0 & 0 & 0 & -1 & 0 & 0 & 0\\
0 & 0 & 0 & 0 & 0 & 0 & 0 & 1 & 0 & 0 & 0 & 0 & 0 & 0 & 0 & 0 & 0 & 0 & 0 & 0 & 0 & 0 & 0 & 0 & 0 & 0 & 0 & -1 & 0 & 0\\
0 & 0 & 0 & 0 & 0 & 0 & 0 & 0 & 1 & 0 & 0 & 0 & 0 & 0 & 0 & 0 & 0 & 0 & 0 & 0 & 0 & 0 & 0 & 0 & 0 & 0 & 0 & 0 & -1 & 0\\
0 & 0 & 0 & 0 & 0 & 0 & 0 & 0 & 0 & 1 & 0 & 0 & 0 & 0 & 0 & 0 & 0 & 0 & 0 & 0 & 0 & 0 & 0 & 0 & 0 & 0 & 0 & 0 & 0 & -1\\
0 & 0 & 0 & 0 & 0 & 0 & 0 & 0 & 0 & 0 & 0 & 0 & 0 & 0 & 0 & 0 & 1 & 0 & 0 & 0 & 0 & 0 & 0 & 0 & 0 & 0 & -1 & 0 & 0 & 0\\
0 & 0 & 0 & 0 & 0 & 0 & 0 & 0 & 0 & 0 & 0 & 0 & 0 & 0 & 0 & 0 & 0 & 1 & 0 & 0 & 0 & 0 & 0 & 0 & 0 & 0 & 0 & -1 & 0 & 0\\
0 & 0 & 0 & 0 & 0 & 0 & 0 & 0 & 0 & 0 & 0 & 0 & 0 & 0 & 0 & 0 & 0 & 0 & 1 & 0 & 0 & 0 & 0 & 0 & 0 & 0 & 0 & 0 & -1 & 0\\
0 & 0 & 0 & 0 & 0 & 0 & 0 & 0 & 0 & 0 & 0 & 0 & 0 & 0 & 0 & 0 & 0 & 0 & 0 & 1 & 0 & 0 & 0 & 0 & 0 & 0 & 0 & 0 & 0 & -1\\
0 & 1 & 0 & 0 & 0 & 0 & 0 & 0 & 0 & 0 & 0 & 0 & 0 & 0 & 0 & 0 & 0 & 0 & 0 & 0 & 0 & 0 & 0 & 0 & 0 & 0 & 0 & 0 & 0 & 0\\
0 & 0 & 0 & 1 & 0 & 0 & 0 & 0 & 0 & 0 & 0 & 0 & 0 & 0 & 0 & 0 & 0 & 0 & 0 & 0 & 0 & 0 & 0 & 0 & 0 & 0 & 0 & 0 & 0 & 0\\
0 & 0 & 1 & 0 & 0 & 0 & 0 & 0 & 0 & 0 & 0 & 0 & 0 & 0 & 0 & 0 & 0 & 0 & 0 & 0 & 0 & 0 & 0 & 0 & 0 & 0 & 0 & 0 & 0 & 0\\
0 & 0 & 0 & 0 & 1 & 0 & 0 & 0 & 0 & 0 & 0 & 0 & 0 & 0 & 0 & 0 & 0 & 0 & 0 & 0 & 0 & 0 & 0 & 0 & 0 & 0 & 0 & 0 & 0 & 0\\
-1 & 0 & 0 & 0 & 0 & 0 & 0 & 0 & 0 & 0 & 0 & 0 & 0 & 0 & 0 & 0 & 0 & 0 & 0 & 0 & 0 & 0 & 0 & 0 & 0 & 0 & 0 & 0 & 0 & 0\\
0 & 0 & 0 & 0 & 0 & 0 & 0 & 0 & 0 & 0 & -1 & 0 & 0 & -1 & 0 & 0 & 0 & 0 & 0 & 0 & 0 & 0 & 0 & 0 & 0 & 0 & 0 & 0 & 0 & 0\\
0 & 0 & 0 & 0 & 0 & 0 & 0 & 0 & 0 & 0 & 0 & 0 & 0 & 0 & -1 & 0 & 0 & 0 & 0 & 0 & 0 & 0 & 0 & 0 & 0 & 0 & 0 & 0 & 0 & 0\\
0 & 0 & 0 & 0 & 0 & 0 & 0 & 0 & 0 & 0 & 0 & 1 & 0 & 0 & 0 & 0 & 0 & 0 & 0 & 0 & 0 & 0 & 0 & 0 & 0 & 0 & 0 & 0 & 0 & 0\\
0 & 0 & 0 & 0 & 0 & 0 & 0 & 0 & 0 & 0 & -1 & 0 & 1 & 0 & 0 & 0 & 0 & 0 & 0 & 0 & 0 & 0 & 0 & 0 & 0 & 0 & 0 & 0 & 0 & 0\\
0 & 0 & 0 & 0 & 0 & 0 & 0 & 0 & 0 & 0 & -1 & 0 & 1 & -1 & 0 & -1 & 0 & 0 & 0 & 0 & 0 & 0 & 0 & 0 & 0 & 0 & 0 & 0 & 0 & 0\\
0 & 0 & 0 & 0 & 0 & 0 & 0 & 0 & 0 & 0 & 0 & 0 & 0 & 0 & 0 & 0 & 0 & 0 & 0 & 0 & -1 & -1 & 0 & 0 & 0 & 0 & 0 & 0 & 0 & 0\\
0 & 0 & 0 & 0 & 0 & 0 & 0 & 0 & 0 & 0 & 0 & 0 & 0 & 0 & 0 & 0 & 0 & 0 & 0 & 0 & 0 & 0 & -1 & 0 & 0 & 0 & 0 & 0 & 0 & 0\\
0 & 0 & 0 & 0 & 0 & 0 & 0 & 0 & 0 & 0 & 0 & 0 & 0 & 0 & 0 & 0 & 0 & 0 & 0 & 0 & 0 & 0 & 0 & -1 & 0 & 0 & 0 & 0 & 0 & 0\\
0 & 0 & 0 & 0 & 0 & 0 & 0 & 0 & 0 & 0 & 0 & 0 & 0 & 0 & 0 & 0 & 0 & 0 & 0 & 0 & -1 & 0 & 0 & 0 & -1 & 0 & 0 & 0 & 0 & 0\\
0 & 0 & 0 & 0 & 0 & 0 & 0 & 0 & 0 & 0 & 0 & 0 & 0 & 0 & 0 & 0 & 0 & 0 & 0 & 0 & -1 & -1 & 0 & 0 & -1 & -1 & 0 & 0 & 0 & 0\\
\end{pmatrix}$}
\]
\par\smallskip\noindent
$R_{H_1}$ has size $4\times 4$; its full integer matrix is
\[
\resizebox{!}{11.2mm}{$\displaystyle R_{H_1}=\begin{pmatrix}
1 & 0 & 0 & 1\\
0 & 0 & 1 & 0\\
0 & 1 & 0 & 0\\
0 & 0 & 1 & 1\\
\end{pmatrix}$}
\]
\par\smallskip\noindent
$R_{H_2}$ has size $6\times 6$; its full integer matrix is
\[
\resizebox{!}{16.799999999999997mm}{$\displaystyle R_{H_2}=\begin{pmatrix}
1 & 0 & 0 & 0 & 1 & 1\\
0 & 1 & 0 & 0 & 0 & -1\\
0 & 0 & 1 & 0 & 1 & 0\\
0 & 0 & 0 & 0 & -1 & 0\\
0 & 0 & 0 & 0 & 0 & -1\\
0 & 0 & 0 & 1 & 1 & 1\\
\end{pmatrix}$}
\]
\endgroup


\begin{statement}{P4.5}{Proved}
For this example,
\[
 a(Y)=1,\qquad c_3[Y]=0,\qquad b_3(Y)=2,
\]
and the integral homology is free of ranks $(1,0,0,2,0,0,1)$.
\end{statement}
\begin{proof}
The regular fibre has Euler number zero. Both components of each special
fibre are elliptic bundles over compact curves, and their intersection is
elliptic, so additivity gives zero for each special support. Stratifying
the base and applying multiplicativity to the smooth bundle gives
$e(Y)=0$. The Euler class of the real bundle underlying a complex
rank-three tangent bundle is its top Chern class, with the complex
orientation. Hence $c_3[Y]=0$. Duality and P4.4 give $b_4=b_5=0$ and
$e(Y)=2-b_3$, so $b_3=2$. Integral duality and the universal-coefficient
theorem identify the torsion of $H_3$ with that of $H_2$, which is zero.
All other groups have already been determined.

We also have $H^2(Y;\Z)=0$. A horizontal effective divisor would restrict
to a nonzero effective divisor on the Kähler torus $A$ with zero first
Chern class. Its pairing with a Kähler class would be positive by
integration over the divisor and zero by the Chern-class computation.
Thus every divisor on $Y$ is vertical. The polar divisor of a meromorphic
function is consequently vertical; on a general fibre the function is
holomorphic and therefore constant. Local holomorphic sections give a
descended meromorphic function on the regular part of the base. Near
each exceptional base value a sufficiently high power of the base
coordinate removes the vertical poles. Connected fibres then give
holomorphic descent there, and dividing by that power gives meromorphic
descent. Therefore $\C(Y)=f^*\C(\Pp^1)$ and $a(Y)=1$.
\end{proof}

\subsection{The normal-crossings sign theorem}

\begin{statement}{P4.6}{Proved}
Let $X$ be a compact complex threefold with $b_2(X)=0$, and let $F$ be a
reduced compact normal-crossings divisor. We allow different local
branches of one global irreducible component to meet. If $T(F)$ is the
number of ordinary triple points, then
\[
 e(F)=T(F),\qquad T(F)\text{ is even}.
\]
This assertion concerns the reduced support, regardless of the
multiplicities with which it occurs in a fibre.
\end{statement}
\begin{proof}
Write $\nu:S=\coprod S_i\to X$ for the normalization immersion, and $D_i$
for the conductor divisor on $S_i$. The surfaces $S_i$ are smooth and
the curves $D_i$ are reduced nodal divisors. The quotient
$N_i=\nu_i^*T_X/T_{S_i}$ is a line bundle. A local equation of $F$ gives
\[
 N_i\simeq \nu_i^*\OO_X(F)\otimes\OO_{S_i}(-D_i).
\]
Indeed, for the coordinate equation $x_1\cdots x_k=0$ the normal
derivative on $x_1=0$ is $x_2\cdots x_k$, whose zero divisor is $D_i$.
These are morphisms of holomorphic vector bundles, including when
$S_i$ is not embedded globally.

Since $b_2(X)=0$, both $c_1(T_X)$ and $c_1(\OO_X(F))$ vanish in real
cohomology. Duality gives $b_4(X)=0$, so $c_2(T_X)$ also vanishes. The
tangent exact sequence yields
\[
 K_{S_i}\equiv-D_i,\qquad c_2(S_i)=K_{S_i}^2.
\]
If $a=\sum_i\chi(\OO_{S_i})$, the surface Riemann--Roch formula
$12\chi(\OO_{S_i})=K_{S_i}^2+e(S_i)$ consequently gives
$\sum_i e(S_i)=6a$.

Let $\Gamma$ be the normalization of the double-curve locus, with one
sheet for each unordered pair of local branches. The normalization of
$\coprod D_i$ is an unramified degree-two cover of $\Gamma$, recording
the choice of branch. Each triple point contributes three nodes to
$\coprod D_i$. Adjunction gives
$\chi(\OO_{D_i})=-D_i\cdot(D_i+K_{S_i})/2=0$. The normalization
sequence for these nodal curves therefore gives
\[
 0=2\chi(\OO_\Gamma)-3T(F)=e(\Gamma)-3T(F).
\]
The normalization complex of $F$ is the exact alternating restriction
complex of the local coordinate hyperplanes. Permuting the branches
globally replaces its double-curve term by $\OO_\Gamma$ tensored with
the sign line of the two-sheeted branch cover. That line has order at
most two, and hence degree zero on every component; it has the same
coherent Euler characteristic as $\OO_\Gamma$. The triple term has
length one at each triple point. It follows that
\[
 \chi(\OO_F)=a-\chi(\OO_\Gamma)+T(F)=a-T(F)/2.
\]
The exact sequence for the Cartier divisor $F$ and analytic
Hirzebruch--Riemann--Roch give $\chi(\OO_F)=0$, since
$\operatorname{ch}(\OO_X(-F))=\operatorname{ch}(\OO_X)$ in rational
cohomology. Thus $T(F)=2a$. Finally ordinary Euler additivity gives
\[
 e(F)=\sum_i e(S_i)-e(\Gamma)+T(F)=6a-2T(F)=T(F).
\]
The alternating count is also verified directly at a point with one,
two or three branches. This proves the identity and parity, without
using fibre multiplicities or a base change.
\end{proof}

\begin{statement}{P4.7}{Proved}
Let $f:X\to B$ be a proper surjective holomorphic map with connected
fibres from a compact complex threefold satisfying $b_1(X)=b_2(X)=0$
to a smooth compact curve. Suppose all reduced fibre supports have
normal crossings. Then
\[
 B=\Pp^1,\qquad
 c_3[X]=\sum_{b\in B}T((f^{-1}(b))_{\mathrm{red}})\in\{0,2\},\qquad
 b_3(X)=2-c_3[X].
\]
The numerical conclusion does not require the general fibre to be a
torus.
\end{statement}
\begin{proof}
Connected fibres give $f_*\R=\R$. The beginning of the constant-sheaf
Leray sequence injects $H^1(B;\R)$ into $H^1(X;\R)$, so $B$ has genus
zero. A smooth fibre $G$ has trivial normal line bundle. Its tangent
sequence gives $c_2(G)=c_2(X)|_G=0$, hence $e(G)=0$ for any such fibre.
The critical values are finite. Ehresmann's theorem gives a smooth
bundle over their complement. Euler additivity and P4.6 give the
displayed sum. Duality gives $e(X)=2-b_3(X)$, and the nondegenerate
alternating cup product on $H^3(X;\R)$ makes $b_3$ even. The
nonnegative sum is therefore zero or two.
\end{proof}

\begin{statement}{P4.8}{Proved in the normal-crossings class}
For a torus fibration in P4.7, an irreducible singular support with
rational normalization has Euler number two; every other irreducible
support has Euler number zero. In the zero case the supports have the
forms stated in the prompt: a smooth support has a finite unramified
torus cover; a singular irreducible support has smooth minimal
elliptic-ruled or Hopf normalization; a reducible support has smooth
minimal Hopf or non-Kähler elliptic normalization components, with
smooth elliptic double curves and finite unramified branch maps.
Thus $c_3[X]=0$ precisely when every fibre belongs to these zero cases.
\end{statement}
\begin{proof}
We apply BHPV, Theorem VI.1.1 and Table 10, to the smooth compact
normalization components, using the identity $e(S_i)=K_{S_i}^2$ proved
in P4.6. The relevant non-Kähler classification is stated in
Ornea--Vuletescu--Verbitsky, Theorems 4.2 and 5.1: a minimal non-Kähler
surface is elliptic or of class VII, and the reduced fibres of a
non-Kähler elliptic surface are smooth elliptic curves.

For a component with nonempty conductor, a Kähler form $\omega$ would
give $K_{S_i}\cdot\omega=-D_i\cdot\omega<0$. No positive power of
the canonical bundle can then have a nonzero section. The classification
makes a Kähler component rational or ruled. If it is obtained by $k$
blowups of a ruled surface over a genus-$g$ curve, then
\[
 e=4(1-g)+k,\qquad K^2=8(1-g)-k.
\]
Equality gives $k=2(1-g)$, so $g\le1$ and the genus-one case is minimal.
A rational component has $\chi(\OO)=1$, hence $e=6$ here, including
the case with minimal model $\Pp^2$.

A minimal non-Kähler elliptic surface has $e=K^2=0$, by its smooth
elliptic reduced fibres. A minimal class-VII surface has $b_1=1$ and
$p_g=0$, so $\chi(\OO)=0$ and $K^2=-e$, with $e=b_2\ge0$.
Blowups preserve $\chi(\OO)$ and change $(e,K^2)$ by $(1,-1)$.
Consequently $e=K^2$ forces minimality and $e=0$ in all these
non-Kähler cases. The class-VII case has $b_2=0$. Theorem 1 of
Buonerba--Bogomolov--Kurnosov supplies a finite unramified Hopf or
Inoue cover. The conductor excludes an Inoue cover: Inoue surfaces
have no compact curves, as stated in Khaled--Teleman, Theorem 3.1
(p.~303), also proved for the original covering families in Inoue
1972, Propositions 1(i) and 2(i). The third family has a type-II
unramified double cover (Khaled--Teleman, proof of Remark 3.4,
p.~305). A finite unramified quotient of a Hopf surface has universal
cover $\C^2\setminus\{0\}$ and is therefore a Hopf surface.

For an irreducible singular support there is one normalization
component. P4.6 gives $T=2\chi(\OO_{S_1})$, which is two in the
rational case and zero in the other cases. P4.12 proves that the
dualizing sheaf of an irreducible support is torsion. Adjunction on
the normalization therefore makes $K_{S_1}+D_1$ torsion. Since
$D_1\ne0$, every positive plurigenus of $S_1$ vanishes: multiplying
a pluricanonical section by an appropriate positive power of the
conductor section would give a section of a torsion line bundle with
a zero, whose sufficiently high power is a nonzero constant. Thus a
non-Kähler normalization in this irreducible case is of class VII,
and hence is the Hopf case already obtained.

For a reducible connected fibre, each normalization component meets
another global irreducible component. Pulling back the divisor of
that component gives a nonzero effective divisor on $S_i$ with zero
real first Chern class. A Kähler form would pair positively with it,
contradicting this zero class. Hence every component is non-Kähler,
has $\chi(\OO)=0$, and $T=0$. The conductor components are disjoint
and smooth, because all their nodes would come from triple points.
For any conductor component $C$, adjunction and $K+D\equiv0$ give
$2g(C)-2=(K+C)\cdot C=0$, since the other conductor components
are disjoint. Thus the double curves are smooth elliptic, and their
branch maps are the unramified maps already described in P4.6.

Finally let the support $G$ of a multiple fibre $mG$ be smooth. The
normalized cyclic base change $t=u^m$ is locally a disjoint union of
smooth charts: the original equation is $t=z^m$ times a holomorphic
unit, whose root can be extracted locally. These charts show that
the new central fibre maps finitely and unramifiedly to $G$, and
that the new map to the $u$-disc is a proper submersion. Its nearby
fibres are the original tori. Catanese, \emph{Deformation in the
large of some complex manifolds, I}, Theorem 2.1, states that every
deformation of a complex torus is a complex torus; its definition
uses precisely a proper holomorphic submersion over a connected
disc, with no Kähler hypothesis on the total space. We have verified
all these hypotheses. The new central fibre is therefore a torus;
Ehresmann's theorem also shows that it is connected. This gives
the required finite unramified cover of $G$.
\end{proof}

\begin{statement}{P4.9}{Proved}
In the requested irreducible and normal-crossing fibre classes,
the complete sign, exact Euler values and normalization equality
classification hold as stated in P4.58. In particular every
irreducible singular support with rationally resolved normalization
has Euler number two, and every other irreducible support has
Euler number zero.
\end{statement}
\begin{proof}
P4.26 proves the unrestricted irreducible numerical sign;
P4.35--P4.37 give the rational and class-VII exact values.
P4.32 proves the zero-Euler classification. The positive ruled
alternative is excluded by the complete sequence P4.54--P4.57
and P4R.18--P4R.20, synthesized in P4.58.
\end{proof}
We checked CDP 2020, Lemma~4.6, Propositions~4.7,
4.11 and~7.3, deriving the sign and classification independently
from its local conductor identities. P4.14 records the precise
incorrect curve-classification inference in Proposition~7.3.
Earlier proof obligations and unsuccessful reductions are retained
in the working notes; the final classification is P4.58.
The original stated-value row also claimed a sign for arbitrary
reducible supports without normal crossings. Such supports are
outside the two classes explicitly requested in Part~4.2; the
proved global assertion retains the stated irreducible-or-NC
support hypothesis.

\subsection{Explicit checks on the hypotheses}

\begin{statement}{P4.10}{Proved}
There is a smooth projective threefold $Z$ with $b_1(Z)=0$ and a
connected torus fibration $Z\to\Pp^1$ whose reduced special fibres
are irreducible normal-crossings surfaces, but $c_3[Z]=24$.
Consequently the hypothesis $b_2=0$ cannot be omitted, even in the
intersection of the two classes of fibres in the question.
\end{statement}
\begin{proof}
On the two affine charts of $\Pp^1$, with $s=t^{-1}$, we take
\[
 y^2z=x^3-3xz^2+t^6z^3,\qquad
 y_\infty^2z=x_\infty^3-3s^4x_\infty z^2+z^3,
\]
glued by $x_\infty=s^2x$ and $y_\infty=s^3y$. These are the
Weierstrass cubics in $\Pp(\OO\oplus\OO(2)\oplus\OO(3))$ for
the line bundle $\OO_{\Pp^1}(1)$. They define an elliptic surface
$\pi:S\to\Pp^1$ with zero section $[0:1:0]$. The discriminant
is a nonzero constant times $t^{12}-4$, or $1-4s^{12}$ on the
other chart. There are twelve simple zeros and a smooth fibre at
infinity. At a discriminant zero the cubic has a double root
$x=1$ or $x=-1$ and a distinct third root, so the fibre has an
ordinary node. Since $6t^5\ne0$ there, the total surface is smooth.
At the zero section the derivative in $z$ is nonzero. The map near
each node is therefore $xy=t$ after holomorphic coordinate changes
and translation of the base.

The relative dualizing differential $dx/y$ has transitions
$\pi^*\OO(1)$ and generates the dualizing sheaf also at the nodes.
Thus $\omega_S=\pi^*\OO(-1)$. Relative duality gives
$R^1\pi_*\OO_S=\OO(-1)$, and connected fibres give
$\pi_*\OO_S=\OO$. Leray yields $h^1(\OO_S)=0$, while all
positive plurigenera vanish. BHPV, Theorem VI.1.1, makes $S$
rational; in the ruled case the base genus equals $h^1(\OO_S)=0$.
In particular $\pi_1(S)=1$, since the minimal rational models
$\Pp^2$ and the Hirzebruch surfaces are simply connected and a
point blowup preserves the fundamental group.

The projective fibre product $W=S\times_{\Pp^1}S$ has exactly
twelve singularities, at the pair of nodes in each nodal fibre.
Their equations are $xy=uv$, with base function $xy=uv$. We blow
up the coherent ideal of the diagonal $\Delta_S\subset W$. This
divisor is Cartier away from the twelve points, so the blowup is
the identity there. Put $r=x-u$, $h=y-v$ at a singular point.
The equation becomes $ry+uh=0$ and the diagonal ideal is $(r,h)$.
The blowup has charts
\[
 h=ar,\quad y=-au,\quad x=u+r,\quad f=-aux;
 \qquad
 r=bh,\quad u=-by,\quad x=-bv,\quad f=-byv.
\]
The independent coordinates are $(a,u,x)$ and $(b,y,v)$.
Both charts are smooth, with $b=a^{-1}$ on the overlap, and the
exceptional fibre is $\Pp^1$. Thus the resulting $Z$ is smooth
and projective. The displayed equations make every special fibre
reduced and normal crossings. Its support remains irreducible:
the original product of two irreducible nodal cubics is irreducible,
and the exceptional locus has dimension one, so introduces no
surface component. Its normalization is rational, by the
birational map from $\Pp^1\times\Pp^1$.

A nodal cubic has Euler number $2-1=1$, so its self-product has
Euler number one. The resolution replaces one point by $\Pp^1$;
each resolved special support therefore has Euler number two.
All other fibres are $E_t\times E_t$ with Euler number zero.
Hence $e(Z)=12\cdot2=24=c_3[Z]$.

For completeness we compute $\pi_1(Z)$. Every loop can be moved
off the fibre divisors, so the fundamental group of their complement
surjects onto $\pi_1(Z)$. The complement is a bundle with fibre
$E\times E$ and a section over a twelve-punctured sphere. The
section extends over $\Pp^1$, killing the base generators. The
two maps
\[
 S\longrightarrow Z:\quad p\mapsto(p,O(\pi(p))),\qquad
 p\mapsto(O(\pi(p)),p)
\]
are defined because their images miss every conifold point, and
extend the two factor inclusions of $E$. Since $S$ is simply
connected they kill both fibre fundamental groups. Thus
$\pi_1(Z)=1$ and $b_1(Z)=0$. A projective Kähler class is nonzero
in $H^2(Z;\R)$, because its cube has positive integral. This
identifies the omitted hypothesis: $b_2(Z)>0$.
\end{proof}

\begin{statement}{P4.11}{Proved}
There is a smooth compact complex threefold $C$ with
$b_1(C)=b_2(C)=0$, $e(C)=0$, and a holomorphic submersion
$C\to\Pp^1$ whose fibres are primary Hopf surfaces. No fibre has
a finite unramified torus cover. Thus the torus hypothesis is
necessary for the equality classification as stated in the prompt.
P4.7 proves it unnecessary for the numerical normal-crossings
sign theorem.
\end{statement}
\begin{proof}
The additive group $\C$ acts on
$V=(\C^2\setminus\{0\})\times(\C^2\setminus\{0\})$ by
\[
 \tau\cdot(z,w)=(e^\tau z,e^{i\tau}w).
\]
For $\tau=a+ib$, the logarithms of the two Euclidean norms change
by $a$ and $-b$. The action is free. It is proper: when both a
point and its translate lie in prescribed compact subsets of $V$,
the two norm ratios bound $a$ and $b$. Each orbit has a unique
representative with both norms one, obtained with
$a=-\log\|z\|$ and $b=\log\|w\|$. This gives an explicit
diffeomorphism $C=V/\C\cong S^3\times S^3$ and proves compactness.
Holomorphic transverse slices, for example $z_j=1$ after choosing
a local logarithm, give quotient charts with holomorphic
transitions. Hence $C$ is a smooth compact complex threefold.
The product cell complex gives free integral homology of ranks
$(1,0,0,2,0,0,1)$ and therefore the stated Betti and Euler numbers.

The invariant map $[z,w]\mapsto[z_1:z_2]$ is a holomorphic
submersion to $\Pp^1$. In an affine base chart write
$z=\lambda(1,\xi)$ and set $\lambda=1$ using the action.
The choices of logarithm differ by $2\pi ik$, changing $w$ by
$e^{-2\pi k}$. Thus we obtain the explicit local fibre bundle
with fibre
$H=(\C^2\setminus\{0\})/\langle w\mapsto e^{-2\pi}w\rangle$.
Its universal cover retracts onto the simply connected $S^3$, so
$\pi_1(H)=\Z$. A connected finite unramified torus cover would
inject $\Z^4$ as a finite-index subgroup of $\Z$. Every subgroup
of $\Z$ is cyclic, whereas $\Z^4$ has rank four, which is
impossible. This proves the failure of the stated smooth-support
classification using the exact covering morphism.
\end{proof}

The executed program \texttt{code/part4/necessity\_examples.py}
checks the degree and squarefreeness of the discriminant, both
resolution substitutions, the Euler count $24$, and the product
cell ranks used above. Its complete output is supplied.

\subsection{Irreducible supports: canonical torsion and the normal case}

\begin{statement}{P4.12}{Proved}
In the torus-fibration setting of Part 4.2, let $f^{-1}(b)=mS$
have irreducible reduced support. Then
\[
 \omega_S^{\otimes m}\simeq\OO_S,\qquad
 \chi(\OO_S)=0,\qquad h^1(\OO_S)\in\{1,2\}.
\]
Every normal such support is smooth. For a singular, necessarily
nonnormal support, a minimal resolved model of its normalization
has Kodaira dimension $-\infty$.
\end{statement}
\begin{proof}
\emph{Canonical torsion.}
Choose a disc $\Delta$ with no other critical value, with coordinate
$t$ vanishing at $b$, and write $U=f^{-1}(\Delta)$. The coherent
direct image $f_*\omega_U$ is torsion-free of rank one: nonzero
base functions act injectively on sections, and a nearby torus has
one-dimensional canonical space. It is therefore a line bundle
on the disc. A generator gives a section $v$ of $\omega_U$,
nonzero and nowhere vanishing on every noncentral fibre. Its
divisor is $aS$ for an integer $a\ge0$. If $a\ge m$, then $v/t$
would be holomorphic, contradicting primitivity of the generator
at zero. We have $0\le a<m$ and
\[
 \omega_U\simeq\OO_U(aS),\qquad\OO_U(mS)\simeq\OO_U.
\]
Adjunction gives $\omega_S\simeq\OO_U((a+1)S)|_S$, proving
$\omega_S^{\otimes m}\simeq\OO_S$. The ambient
Riemann--Roch computation of P4.6, applied to $S$, gives
$\chi(\OO_S)=0$. Reduced irreducibility gives $h^0(\OO_S)=1$.
A nonzero section of a torsion line bundle has a power equal to
a nonzero constant, and hence is nowhere zero; its space of
sections consequently has dimension at most one. Serre duality
on the Gorenstein surface $S$ gives
$h^2(\OO_S)=h^0(\omega_S)\in\{0,1\}$, proving the assertion
about $h^1$.

\emph{Exceptional curves and discrepancies.}
For a minimal resolution $\pi:T\to S$ of a normal Gorenstein
surface write $K_T=\pi^*K_S+\sum a_jE_j$. Its exceptional
intersection matrix $M$ has nonnegative off-diagonal entries
and is negative definite. We recall the local analytic proof.
Pulling back a generic holomorphic function vanishing at a
contracted point gives $H+\sum m_jE_j$, with $m_j>0$ and $H$
the strict transform. Intersecting gives $Mm\le0$, strictly in
at least one entry of each connected block. After scaling by
$\operatorname{diag}(m_j)$, the negative quadratic form is a
sum of nonnegative terms $c_{ij}(z_i-z_j)^2$ and $d_i z_i^2$,
with connected support and some $d_i>0$ in every block. It is
positive definite. This is the elementary argument in Mumford
1961, Section I, p.~6 (an unnumbered assertion and proof), and
applies directly to these analytic germs.

Adjunction gives
$K_T\cdot E_j=2p_a(E_j)-2-E_j^2\ge0$: the sole negative
possibility is a smooth rational $(-1)$-curve, excluded by
minimality of the resolution. The inverse of a negative definite
matrix with nonnegative off-diagonal entries has nonpositive
entries. Indeed, put $Q=-M$ and choose $c$ larger than all its
eigenvalues and diagonal entries; then
$Q^{-1}=c^{-1}\sum_{k\ge0}(I-Q/c)^k$ is entrywise nonnegative.
The discrepancy equations $Ma=(K_T\cdot E_j)_j$ imply
$a_j\le0$. Hence $K_T=\pi^*K_S-E$ with $E$ an effective
exceptional integral divisor.

Put $p=\dim H^0(S,R^1\pi_*\OO_T)$. Leray, with this sheaf
supported at finitely many points, gives
\begin{equation}\label{eq:P4chi-resolution}
 \chi(\OO_T)=\chi(\OO_S)-p=-p.
\end{equation}
For any effective exceptional cycle $Z$, the higher direct image
$R^2\pi_*I_Z$ vanishes because all fibres of $\pi$ have dimension
at most one. The exact sequence of $Z$ gives a surjection
$R^1\pi_*\OO_T\to R^1\pi_*\OO_Z$, and therefore
\begin{equation}\label{eq:P4exceptional-genus}
 p\ge h^1(Z,\OO_Z).
\end{equation}

\emph{Exclusion of normal singular supports.}
Suppose first that $E=0$. Then $K_T$ is torsion and the
classification gives Kodaira dimension zero and nonnegative
holomorphic Euler characteristic. Equation
\eqref{eq:P4chi-resolution} forces $p=0$ and $\chi(\OO_T)=0$.
The remaining surfaces are tori, bielliptic surfaces and Kodaira
surfaces. They have no rational curves: a map from $\Pp^1$ to
a torus is constant, bielliptic surfaces have finite unramified
torus covers, and a primary Kodaira surface has an elliptic
fibration with smooth elliptic reduced fibres over an elliptic
curve. A map from $\Pp^1$ to a secondary Kodaira surface lifts
to its finite unramified primary Kodaira cover. On a torsion
canonical surface a negative curve is a smooth rational
$(-2)$-curve by adjunction. Thus there are no exceptional
curves, and $S$ is smooth.

Suppose next that $E\ne0$. Since $K_T+E$ is torsion, all
positive plurigenera vanish, so $T$ is rational, ruled, or a
blowup of a class-VII surface. Rationality contradicts
\eqref{eq:P4chi-resolution}, since $\chi(\OO_T)=1$. In class
VII that equation gives $p=0$, whereas
\[
 \chi(\OO_E)=-E\cdot(E+K_T)/2=0.
\]
The nonempty effective compact divisor $E$ has $h^0(\OO_E)\ge1$,
so $h^1(\OO_E)\ge1$, contradicting
\eqref{eq:P4exceptional-genus}.

Finally suppose $T$ is ruled, after point blowups, over a curve
$C$ of genus $g\ge1$. Equation \eqref{eq:P4chi-resolution}
gives $p=g-1$. A general ruling fibre $F$ has $K_T\cdot F=-2$,
so $E\cdot F=2$. An exceptional component $E_j$ therefore
maps onto $C$. Its normalization has genus at least $g$ by
Riemann--Hurwitz, so $h^1(\OO_{E_j})\ge g$. Taking $Z=E_j$
in \eqref{eq:P4exceptional-genus} gives $p\ge g$, a contradiction.
This exhausts the normal case.

\emph{The resolved normalization of a nonnormal support.}
Write $\nu:N\to S$ for the finite normalization and $\pi:T\to N$
for its minimal resolution. The residue comparison of dualizing
sheaves gives a pole divisor along the conductor; its strict
transform $D$ on $T$ is nonzero and effective. For completeness,
pull back a local generator of $\omega_S$ as a meromorphic
two-form. Its canonical divisor has the form $-D+\sum a_jE_j$.
The equations for the exceptional coefficients have right side
$(K_T+D)\cdot E_j\ge0$, because both $K_T\cdot E_j$ and
$D\cdot E_j$ are nonnegative. The same inverse-matrix argument
gives $a_j\le0$. Taking the $m$th power to trivialize the pulled
back $\omega_S$ consequently gives
\[
 mK_T+D'\sim0
\]
for a nonzero effective integral divisor $D'$. This is also the
conductor formula of CDP 2020, Lemma 4.6; its local proof does
not use the disputed global restriction conclusion. Any
positive pluricanonical section on $T$, multiplied by the
appropriate effective section and raised to a power, would give
a nonzero holomorphic function on compact $T$ with a zero.
It follows that $\kappa(T)=-\infty$. BHPV, Theorem VI.1.1,
gives precisely the stated alternatives for the minimal model.
\end{proof}

The ruled-genus and nonnormal-conductor arguments needed after
this reduction are proved below and completed in P4.58.

\begin{statement}{P4.13}{Proved}
Let a proper holomorphic degeneration of two-dimensional complex
tori have smooth total threefold and central fibre $mS$ with
irreducible reduced support. If $A$ is the monodromy on $H_1$ of
a nearby torus, then
\[
 m e(S_{\mathrm{reg}})=\det(I-A^m),\qquad
 0\le e(S_{\mathrm{reg}})\le16/m.
\]
This assertion requires no Betti-number hypothesis on a compactification.
\end{statement}
\begin{proof}
We resolve the pair by a proper holomorphic modification
$\pi:U'\to U$ which is an isomorphism away from $S$ and over
$S_{\mathrm{reg}}$. The total central divisor has normal crossings.
The strict transform has multiplicity $m$, and each exceptional
component has multiplicity at least $2m$. Indeed, its first centre
lies over a singular point of the reduced divisor, whose order is
at least two; subsequent centres either retain such an exceptional
multiplicity or add further positive multiplicities.

We spell out the trace comparison. At a point on $k$ branches,
with local equation $z_1^{n_1}\cdots z_k^{n_k}=t$, the nearby
fibre retracts onto a disjoint union of real $(k-1)$-tori.
Monodromy permutes these components and translates any component
that it preserves. For $k\ge2$ the alternating cohomological
trace is zero. On a single branch of multiplicity $n$ there are
$n$ contractible nearby components in a cyclic permutation;
the $d$th iterate has trace $n$ when $n\mid d$, and zero otherwise.
Alternating traces in the nearby-cycle spectral sequence, with a
finite stratification of the compact central divisor, therefore give
\[
 L(A^d;H^*(T^4;\C))=\sum_{n_j\mid d}n_j e(E_j^\circ).
\]
This is the global form of A'Campo 1975, Theorem 1; his proper
analytic construction in Section 2 and spectral-sequence proof in
Section 3 give precisely this comparison. The explicit local
calculation verifies the hypotheses without a Kähler assumption.
For $d=m$ the strict transform is the sole contribution, and its
open stratum is $S_{\mathrm{reg}}$.

These same charts prove quasi-unipotence. A common multiple $N$
of the $n_j$ makes monodromy act trivially on all local
nearby-cohomology sheaves, hence on the $E_2$ page. Its action
on the filtered abutment is consequently unipotent. Every
eigenvalue of $A$ is a root of unity. Since torus cohomology is
the exterior algebra on $H_1$, its Lefschetz number is
$\det(I-A^m)$. Complex conjugate eigenvalues contribute
$|1-\lambda^m|^2$; real eigenvalues are $1$ or $-1$ and
contribute zero or two. Thus this determinant is nonnegative
and at most $2^4$, as required.
\end{proof}

\begin{statement}{P4.14}{Proved; a specific intermediate inference is false}
The reducible-curve inference in CDP 2020, Proposition 7.3,
printed p.~694, omits the two plane cubics
\[
 Z(Y^2-XZ)=0,\qquad XY(X-Y)=0.
\]
They are reduced Gorenstein curves of arithmetic genus one with
trivial dualizing sheaf and rational normalization, but have
respectively two tangent components and three concurrent
components. Their conductor order is two on every normalization
branch. These curves refute the asserted nodal-cycle
classification at that step; they are not counterexamples to
the threefold sign assertion.
\end{statement}
\begin{proof}
Plane adjunction gives $\omega_C=\OO_C(3-3)=\OO_C$ and
$p_a(C)=1$. The first cubic has a line tangent to a smooth
conic at $[1:0:0]$. Its local equation is $y(y-x^2)=0$.
The normalization restrictions are $(p(t,0),p(t,t^2))$.
Their image is exactly the pairs $(a,b)$ with $b-a$ divisible
by $t^2$: the inverse expression is
$p(x,y)=a(x)+y(b(x)-a(x))/x^2$. Thus the normalization
quotient has length two and the conductor exponent is two on
each branch.

For the second cubic the branches are $(t,0),(0,t),(t,t)$.
All constant coefficients must agree, and the third linear
coefficient is the sum of the first two. These three conditions
are sufficient. In every degree $k\ge2$, the monomials
$x^k-x^{k-1}y$, $y^k-xy^{k-1}$ and $x^{k-1}y$ independently
realize the coefficients on the three branches. Hence the
normalization quotient has length three and each conductor
exponent is again two. The normalization Euler characteristics
are two and three; subtracting these quotient lengths gives
$\chi(\OO_C)=0$ in both cases.

The program \texttt{code/part4/genus\_one\_conductor.py}
executes the exact jet restriction matrices through orders
$2,3,4,6,8$, computing the ranks and conductor exponents.
When these branches vary in a surface, their permutation may
assemble them into a single normalization curve covering the
image conductor with degree two or three. Consequently a
unique horizontal conductor component with multiplicity two
does not alone force a degree-one map or generically reduced
image conductor. An additional argument is needed at that
precise point of the general irreducible case.
\end{proof}

The branch permutations in P4.14 have explicit local realizations.
Take an elliptic curve $C$ with a free translation of order two or
three. The map $(x,y)\mapsto(x,x^2-y)$ exchanges the two tacnode
branches and preserves $y(y-x^2)$. The order-three map
$(x,y)\mapsto(-y,x-y)$ cycles the three branches and preserves
$xy(x-y)$. The diagonal action on $C$ times a small invariant
neighbourhood of the origin is free. Its quotient is a smooth
threefold; the divisor is irreducible because the action permutes
its branches transitively. Its normalization conductor is one
elliptic curve mapping to $C/\langle\tau\rangle$ with degree two
or three. These specify the finite morphisms in P4.14 explicitly;
they make no assertion of a compact torus smoothing.

\begin{statement}{P4.15}{Proved}
Let $mS$ be an irreducible singular fibre in the setting of 4.2.
Suppose its normalization $T$ is smooth and $d\nu:T_T\to\nu^*T_X$
has rank two away from finitely many points. Let $\ell$ be the
length of its maximal-minor scheme. Then
\[
 e(T)=K_T^2-\ell.
\]
The alternatives for $T$ are a rational surface with
$e(T)=6-\ell/2$ and $0\leq\ell\leq6$, a minimal ruled surface
over an elliptic curve with $\ell=0$, or a Hopf surface with
$\ell=0$. In either of the last two alternatives $e(S)=0$.
The reduced conductor and singular locus are disjoint unions of
smooth elliptic curves, and their normalization maps are finite
and unramified.
\end{statement}
\begin{proof}
The differential is injective as a sheaf map: its kernel vanishes
on a dense open set, and a locally free sheaf is torsion-free.
Its rank-one cokernel $Q$ is locally free away from finitely many
points. The depth lemma applied to the length-one locally free
resolution gives depth at least one at those points, excluding
a finite-length subsheaf. All torsion would be supported there,
so $Q$ is torsion-free. On the smooth surface we obtain
\[
 Q=\mathcal I_Z\otimes L,\qquad
 L=K_T\otimes\nu^*(\det T_X),\qquad \operatorname{length}(Z)=\ell.
\]
The two-column presentation by $d\nu$ identifies the local ideal
of $Z$ with its three maximal minors. The classes $c_1(T_X)$ and
$c_2(T_X)$ vanish in real cohomology, since $b_2(X)=b_4(X)=0$.
The second Chern class of the exact sequence therefore gives
\[
 0=e(T)-K_T^2+\ell.
\]
P4.12 gives $\kappa(T)=-\infty$. For a ruled surface over a
genus-$g$ curve, blown up $k$ times, we have
\[
 e(T)=4(1-g)+k,\qquad K_T^2=8(1-g)-k.
\]
Consequently $4(g-1)+2k+\ell=0$. Thus $g\leq1$, and $g=1$
forces $k=\ell=0$. In the class-VII case, including blowups,
$K_T^2=-e(T)$ and $e(T)=b_2(T)\geq0$. Hence
$2e(T)+\ell=0$, and the surface is minimal with $b_2=0$.
Its nonempty conductor, the finite-cover classification, and the
absence of curves on Inoue surfaces proved in P4.8 force it to
be Hopf. For rational $T$, Noether's formula gives
$2e(T)+\ell=12$. The inequality $e(T)\geq3$ gives the asserted
range. When $\ell=0$, the cokernel is a line bundle and $d\nu$
has rank two everywhere.

We now compute the image Euler number in the two latter cases.
The conductor divisor $D$ satisfies $K_T+D\equiv0$. First let
$T$ be a minimal ruled surface over an elliptic curve $E$;
it is a $\mathbb P^1$-bundle. For a ruling fibre $F$, $D\cdot F=2$.
If a horizontal component $H$ occurs with multiplicity one,
adjunction gives
\[
 2p_a(H)-2=(K_T+H)H=-(D-H)H\leq0.
\]
The map to $E$ gives geometric genus at least one. Equality
therefore holds, $H$ is smooth elliptic, and $H$ is disjoint
from every other component of $D$. There is no vertical
component because every ruling fibre meets $H$. The possibilities
are one elliptic double-section or two disjoint elliptic sections.

The other possibility is $D=2H+\sum_j a_jF_j$, where $H$ is a
section, $F_j$ are distinct ruling fibres and $a_j>0$.
The reduced conductor has Euler number equal to the number $k$
of these fibres: $e(H)=0$, $e(F_j)=2$, and each meets $H$ once.
The immersive normalization identifies $T\setminus D_{\mathrm{red}}$
with $S_{\mathrm{reg}}$. Hence
\[
 e(S_{\mathrm{reg}})=e(T)-e(D_{\mathrm{red}})=-k.
\]
P4.13 forces $k=0$. Thus in every case the reduced conductor is
a disjoint union of smooth elliptic curves.

On a Hopf surface every irreducible curve $H$ satisfies
$H^2=K_TH=0$, so $p_a(H)=1$. A rational curve would lift from
$\mathbb P^1$ to the universal cover $\mathbb C^2\setminus\{0\}$,
where both coordinate functions must be constant. Thus $H$
is smooth elliptic. Distinct curves have intersection zero and
are disjoint, proving the same conductor description.

To check the image, at a singular point the finite immersive
normalization presents the germ as distinct smooth surface
branches $S_i$. The reduced conductor on each branch is a single
smooth curve germ $C_i$. Any two branches intersect in a curve
in the smooth ambient threefold, and its support lies in both
$C_i$ and $C_j$. Their irreducibility forces those curve germs
to agree. Thus the image singular locus is one smooth curve
germ, and the conductor maps locally isomorphically to it.
Globally this map is finite and unramified. Riemann--Hurwitz
then makes every image component elliptic. Finally
\[
 e(S)=e(T)-e(D_{\mathrm{red}})+e(\operatorname{Sing}S)=0.
\]
\end{proof}

\begin{statement}{P4.16}{Proved}
For every integer $n\geq1$, the reduced hypersurface germ
$x^2-zy^n=0$ cannot occur in a fibre of a proper holomorphic
degeneration of complex two-tori with smooth total threefold,
even with any positive common multiplicity on that germ.
In particular the Whitney umbrella ($n=2$) is excluded.
No hypothesis on $b_2$ is needed.
\end{statement}
\begin{proof}
Put $F_n=\{x^2-zy^n=1\}\subset\mathbb C^3$. We exhibit its
compact spherical homology class and compute its square.
The open subset $U=\{y\ne0\}$ is
$\mathbb C^*\times\mathbb C$ in coordinates $(y,x)$.
For $0<\delta<1$, take the two disjoint graph neighbourhoods
\[
 N_\pm:\quad x=\pm\sqrt{1+zy^n},\quad |zy^n|<\delta.
\]
Their $(y,z)$ domains contract by simultaneous real scaling.
Each $N_\pm\cap U$ has coordinates $(y,zy^n)$ in
$\mathbb C^*\times\{|w|<\delta\}$. Its map to $U$ takes
the generating $y$-loop to the generating $y$-loop. These three
sets cover $F_n$. Van Kampen gives $\pi_1(F_n)=0$.
Mayer--Vietoris identifies $H_2(F_n;\mathbb Z)$ with the
kernel of $(a,b)\mapsto a+b$ on the two intersection
$H_1$ groups, hence with $\mathbb Z$. Hurewicz represents its
generator by a map $S^2\to F_n$.

We compute the intersection form by compactification. Blow up
$\mathbb P^1_x\times\mathbb P^1_y$ successively $n$ times above
each point $(1,0)$ and $(-1,0)$. Near either point let
$a=x^2-1$, $b=y$. The successive charts at the centres are
\[
 a=b^j a_j,\qquad z=\frac{a_j}{b^{n-j}}.
\]
For $j<n$, blow up the next centre $a_j=b=0$.
At $j=n$, $z=a_n$ is regular. The final exceptional curve
minus its point on the preceding boundary is precisely the
line $y=0,x=\pm1$ in $F_n$. Every earlier exceptional curve
is a pole divisor of $z$. Consequently $F_n=T_n\setminus B$,
where $T_n$ is this smooth projective blown-up surface and
$B$ consists of the two infinity divisors, the strict transform
of $y=0$, and the first $n-1$ exceptional curves on each chain.
Both infinity divisors are removed, so no extension of the
rational function at their intersection is asserted or needed.

In the total-transform basis $H_x,H_y,E_{\pm,j}$ the nonzero
products are $H_xH_y=1$ and $E_{\pm,j}^2=-1$. The boundary
classes are
\[
 H_x,\ H_y,\ H_y-E_{+,1}-E_{-,1},\quad
 E_{+,j}-E_{+,j+1},\ E_{-,j}-E_{-,j+1}\quad(1\leq j<n).
\]
They form a connected tree of rational curves; thus $H_1(B)=0$.
Compact-support duality and the sequence of $(T_n,B)$ identify
the image of $H_2(F_n)$ in $H_2(T_n)$ with the kernel of
intersection against these boundary classes. Orthogonality to
$H_x,H_y$ removes both ruling coefficients. Orthogonality to
the exceptional differences makes the coefficients constant
along each chain. The strict transform of $y=0$ makes the two
constants opposite. Thus the kernel is generated by the
primitive integral vector
\[
 v=\sum_{j=1}^n E_{+,j}-\sum_{j=1}^n E_{-,j},\qquad v^2=-2n.
\]
The sequence introduces no multiple: $H_1(B)=0$ injects
$H^2(T_n,B)$ into $H^2(T_n)$ with precisely the stated kernel.
Therefore the compact spherical generator has square $-2n$.
The exact matrices, primitive kernels, boundary graphs and
blowup identities for $n=1,\ldots,12$ are independently checked
by \texttt{code/part4/milnor\_intersection.py}; its output and
JSON data are included. The proof is uniform in $n$.

The defining polynomial is weighted homogeneous of degree
$d=2n+4$ with positive weights $(n+2,2,4)$. Scaling by
\[
 (x,y,z)\longmapsto
 (\varepsilon^{n+2}x,\varepsilon^2y,\varepsilon^4z)
\]
identifies $F_n$ with the fibre of value $\varepsilon^{2n+4}$.
It carries any compact representative of $v$, together with a
compact transverse representative for its square, into an
arbitrarily small neighbourhood of the origin.

We also check the unit and multiplicity in an actual defining
map. Take an analytic $m$-th root of its unit. For a unit $h$,
the local analytic isomorphism
\[
 \Phi(x,y,z)=
 \bigl(h^{(n+2)/d}x,h^{2/d}y,h^{4/d}z\bigr)
\]
satisfies $g\circ\Phi=hg$, where $g=x^2-zy^n$; its derivative
at the origin is invertible. Thus the actual local map to the
base is $(g\circ\Phi)^m$. The small local fibre of $g$ embeds
in the corresponding smooth compact torus fibre.

The spherical representative is null-homotopic in that torus,
whose universal cover is the contractible space $\mathbb C^2$.
It is therefore zero in $H_2$, and its self-intersection there
is zero. The two compact representatives already chosen compute
the same intersection in the open local fibre and in the torus.
Their intersection is $-2n$, which is a contradiction.
\end{proof}

P4.16 excludes the displayed pinching and ramification germs.
The subsequent arguments retain singular normalizations and
nonimmersive normalization maps and complete their classification
in P4.58.

\begin{statement}{P4.17}{Proved}
For all $a\geq2$ and $n\geq1$, the germ $x^a-zy^n=0$ is
excluded in the setting of P4.16. Its affine nearby fibre has
$\pi_1=0$ and compact second-homology intersection lattice
$A_{a-1}(-n)$. This includes the order-three cyclic-ramification
germ $x^3-zy^3$ and examples with singular normalization.
\end{statement}
\begin{proof}
Replace the two roots of $x^2=1$ in P4.16 by the $a$ distinct
roots of $x^a=1$. The open set $y\ne0$ is still
$\mathbb C^*\times\mathbb C$. There are $a$ contractible graph
neighbourhoods, each meeting it in $\mathbb C^*\times\Delta$
with the same generating $y$-loop. Thus van Kampen gives
$\pi_1=0$, and Mayer--Vietoris gives
\[
 H_2(F_{a,n};\mathbb Z)=\ker\bigl(\mathbb Z^a
   \xrightarrow{(c_i)\mapsto\sum_i c_i}\mathbb Z\bigr).
\]
At each root, $x^a-1$ is a local coordinate. The same $n$
successive blowup charts give a compactification with $a$
exceptional chains and boundary classes
\[
 H_x,\ H_y,\ H_y-\sum_i E_{i,1},\qquad
 E_{i,j}-E_{i,j+1}\quad(1\leq i\leq a,\ 1\leq j<n).
\]
This boundary is a tree of rational curves. The exact sequence
of the pair identifies compact $H_2$ integrally with the kernel
of intersection against these classes. Its primitive basis is
\[
 v_i=\sum_{j=1}^n E_{i,j}-\sum_{j=1}^n E_{i+1,j}
       \quad(1\leq i<a).
\]
The products are $v_i^2=-2n$, $v_iv_{i+1}=n$, with all other
off-diagonal products zero. This is $A_{a-1}(-n)$.
Hurewicz makes every class spherical. Positive weights
$(n+1,a,a)$, of degree $a(n+1)$, scale compact representatives
into arbitrarily small nearby fibres. The explicit unit removal
and torus argument of P4.16 now exclude $v_1$, whose square
is $-2n$. The program independently checks the complete
matrices and primitive bases for $2\leq a\leq6$, $1\leq n\leq6$.

We specify the normalization morphism. Let $\mu_a$ act on
$\mathbb C^2_{u,v}$ by $(u,v)\mapsto(\zeta u,\zeta^{-n}v)$.
Then
\[
 x=u^nv,\qquad y=u^a,\qquad z=v^a
\]
gives a finite map to $x^a=zy^n$: both $u$ and $v$ satisfy
monic equations over its coordinate ring. The function-field
degree is $a$, with the displayed faithful $\mu_a$ action.
Consequently the induced quotient map is finite and birational.
The invariant ring is integrally closed, so this is the
normalization. For $(a,n)=(2,3)$ it has coordinates
$U=u^2,V=uv,W=v^2$ with $UW=V^2$, and its map is
$(x,y,z)=(UV,U,W)$. Thus it has an $A_1$ singularity.
For $n=a$ the quotient is smooth with coordinates $(y,v)$,
and its conductor map is $v\mapsto v^a$. This verifies
explicitly the cyclic branch-permutation description.
\end{proof}

\begin{statement}{P4.18}{Proved}
For $a\geq2$ and $b,c\geq1$, the reduced germ
$x^a-y^bz^c=0$ cannot occur with a common positive multiplicity
in a torus degeneration with smooth total threefold. Its affine
nearby fibre contains a spherical class of nonzero square.
In particular this excludes $x^2-z^3y^4$, whose smooth
normalization has a divisorial differential rank drop.
\end{statement}
\begin{proof}
Write $d=\gcd(b,c)$, $b=dB$, $c=dC$, so $\gcd(B,C)=1$.
The nearby fibre $F=\{x^a-y^bz^c=1\}$ maps to the smooth
affine curve $A=\{x^a-w^d=1\}$ by $w=y^Bz^C$.
The function $w$ on $A$ has $a$ simple zeros. Choose an
embedded arc between two zeros avoiding the others, and a
small open disc $D$ around it. In a holomorphic coordinate
$\xi$ on $D$, write $w=u(\xi)(\xi-\alpha)(\xi-\beta)$
with $u$ a unit. Taking a $B$-th root of $u$, rescaling $y$,
and making an affine coordinate change identifies the preimage
of $D$ with
\[
 W=\{y^Bz^C=X^2-1,\ X\in D'\},
\]
where $D'$ is a disc containing both $-1$ and $1$.

Choose $Bp+Cq=1$. With $h=X^2-1$, the explicit trivialization
away from the two special fibres is
\[
 y=h^pt^C,\quad z=h^qt^{-B},\quad t=y^qz^{-p}.
\]
Its fundamental group is $F(\beta_-,\beta_+)\times\mathbb Z\tau$.
The neighbourhood of each special fibre is the graph
$X=\pm\sqrt{1+y^Bz^C}$ on $|y^Bz^C|<\delta$, a contractible
domain by real scaling. Its punctured part has two generating
loops in $y,z$, whose images in $\langle\beta_\pm,\tau\rangle$
are the columns of
\[
 \begin{pmatrix}B&C\\q&-p\end{pmatrix},\qquad \det=-1.
\]
Thus adjoining each neighbourhood kills $\beta_\pm$ and $\tau$.
Van Kampen gives $\pi_1(W)=0$.

The map
\[
 \pi:W\longrightarrow W'=\{y^BZ=X^2-1,\ X\in D'\},
 \qquad (X,y,z)\longmapsto(X,y,z^C)
\]
is proper finite of degree $C$. The compact class $v$ in
P4.16, of square $-2B$, is already supported over $D'$.
Indeed, the cover just given with $C=1$ has the identical
Mayer--Vietoris maps before and after enlarging $D'$ to the
whole plane: both punctured bases have the same two meridians,
and the same two graph neighbourhoods are attached. The
inclusion is therefore an isomorphism on $H_2$ and preserves
the primitive class and its square.
For its compact-support dual $\alpha\in H_c^2(W';\mathbb Z)$,
properness and degree give
\[
 \int_W(\pi^*\alpha)^2=C\int_{W'}\alpha^2=-2BC.
\]
One may verify this by smooth compactly supported forms:
outside the measure-zero branch divisor the map is an oriented
$C$-sheeted covering. The dual of $\pi^*\alpha$ is spherical
by $\pi_1(W)=0$ and Hurewicz. Its compact representatives
embed in $F$ with the same nonzero square.

The polynomial has positive weights $(b+c,a,a)$. Weighted
scaling places these representatives in arbitrarily small
nearby fibres. P4.16's explicit unit removal handles the unit
and common multiplicity. The spherical class becomes zero
in a torus, contradicting its square $-2BC$.

For the mixed example the normalization is the finite
birational map
\[
 (u,v)\longmapsto(u^3v^2,v,u^2).
\]
Finiteness follows from $u^2=z$, and $u=x/(zy^2)$ gives
birationality. The image ring contains all even powers of $u$
and the odd monomials $u^iv^j$ with $i\geq3,j\geq2$.
Its conductor is $u^2v^2\mathbb C\{u,v\}$, giving divisor
$2\{u=0\}+2\{v=0\}$. Its differential minors are
$3u^2v^2,-4u^4v,-2u$, with common divisorial factor $u$.
Thus this example explicitly lies outside P4.15's finite
rank-drop assumption. The 64 exponent matrices for
$1\leq b,c\leq8$, transfer squares and normalization
identities are executed in
\texttt{code/part4/binomial\_milnor.py}; its output and
JSON data are saved alongside it.
\end{proof}

\begin{statement}{P4.19}{Proved}
The germ $y(x^2-zy^2)=0$ has ramified conductor, of multiplicity
two on each normalization component, but its affine nearby
fibre is $\mathbb C^*\times\mathbb C$. Ramification alone
therefore does not give the spherical intersection obstruction
of P4.16--P4.18. This is a local statement, not a compact
counterexample to 4.2.
\end{statement}
\begin{proof}
One normalization component is the plane $y=0$; the other is
the Whitney normalization $(u,v)\mapsto(uv,v,u^2)$.
The singular locus is the smooth $z$-axis. The plane conductor
is $x^2=0$, obtained by restricting the other component's
equation. On the Whitney component the intrinsic conductor
is $v\mathbb C\{u,v\}$; multiplication by the added plane
equation $y=v$ gives $v^2\mathbb C\{u,v\}$. Directly, the
union ring is the fibre product of the component rings over
$\mathbb C\{x,z\}/(x^2)$, giving exactly these conductor ideals.
The reduced conductor cover is the identity sheet from the
plane together with $u\mapsto u^2$ from the Whitney component.
The latter is ramified.

On $y(x^2-zy^2)=1$ we have $y\ne0$, and the inverse maps
\[
 (x,y,z)\longmapsto(y,x),\qquad
 (y,x)\longmapsto\left(x,y,\frac{yx^2-1}{y^3}\right)
\]
identify this fibre with $\mathbb C^*\times\mathbb C$.
Its $\pi_2$ and $H_2$ vanish. Hence every compact spherical
class in an open subset has zero square after inclusion,
and therefore has zero square already in that open subset.
\end{proof}

\begin{statement}{P4.20}{Proved}
Let $mS$ be an irreducible fibre of a proper holomorphic
degeneration of complex two-tori with smooth total threefold.
Suppose the normalization of $S$ is smooth and, outside
finitely many points, $S$ is smooth or has the ordinary-double
germ $xy=0$. Then $S$ has normal crossings everywhere.
In the global setting of 4.2, P4.6--P4.8 apply to this subclass.
\end{statement}
\begin{proof}
We give the boundary and transport arguments explicitly.
Take a small ball $B$ around an exceptional point, with
boundary transverse to the smooth and double strata away
from that point. Let $g$ be a reduced equation. An analytic
$m$-th root of the unit in the actual defining function lets
us write that function as $g^m$.

Consider a small holomorphic deformation $G_s$ of $g$ whose
zero set near $\partial B$ retains its smoothly moving smooth
and ordinary-double strata, transverse to $\partial B$.
There is a uniform $\delta>0$ such that
$G_s|_{\partial B}$ is submersive over $0<|t|<\delta$.
At a boundary point on a double curve, choose coordinates
$a,b,c$ with $G_s=ab$, absorbing a unit in $a$.
Transversality of the double curve to $\partial B$ allows one
real $c$-coordinate to be solved from the boundary equation.
The complex coordinates $a,b$ remain free there, so $ab$ is
submersive when nonzero. At a smooth point use $G_s=a$.
Compactness gives uniform bounds; points outside these
neighbourhoods stay away from zero.

For each fixed small $s$ there are finitely many nonzero
critical values in this disc. After shrinking, the analytic
critical set in a neighbourhood of the compact closed ball
has finitely many relevant local components. The function is constant on
each analytic critical component. A component at a nonzero
small value cannot meet $\partial B$, so it is compact;
the maximum principle excludes a positive-dimensional
compact analytic subset of a ball. Thus the relevant critical
points are isolated and finite. Paths avoiding these values
join a small regular value to one fixed regular value $t_0$.
Proper Ehresmann transport on the closed ball, with the
verified boundary submersion, identifies the fibres along
such paths. For sufficiently small deformation the fibres
$G_s=t_0$ also form a smooth proper family to $s=0$, by
compactness. Hence every compact spherical class of nonzero
square in a local nearby fibre of $G_s$ transports into
$g=t_0$. Its inclusion in the actual torus fibre is impossible.

At a normalization preimage write the map germ as
$\phi:(\mathbb C^2,0)\to(\mathbb C^3,0)$. Away from the
exceptional point it is immersive with transverse multiple
images. If $d\phi$ drops rank, its maximal minors define a
nonempty finite scheme. Small generic polynomial
perturbations have transverse rank-one points near it.
We check both genericity and conservation of their positive
number. Varying linear coefficients makes the universal
derivative map submersive onto the $3$-by-$2$ matrix space.
The rank-one locus is smooth of codimension two; the rank-zero
locus has codimension six. Generic perturbations are transverse
to the first and avoid the second. The height-two maximal-minor
ideal has its Hilbert--Burch resolution. Its one-parameter
zero scheme is Cohen--Macaulay of dimension one; the finite
special fibre excludes a vertical component. Thus it is flat
and preserves the positive length. Further generic polynomial
coefficients keep the finitely many rank-one images off other
sheets: incidence of a zero-dimensional critical set with a
two-dimensional image in a three-dimensional target has
negative expected dimension.

Each transverse rank-one point is a Whitney crosscap, as we
now verify. Choose its first target coordinate as a source
coordinate $u$, and write $\phi=(u,F(u,v),H(u,v))$.
Transversality of the equations $F_v=H_v=0$ gives a linear
combination with nonzero second $v$-derivative. The
holomorphic Morse lemma with parameter $u$, followed by
subtracting the critical value, makes the second coordinate
$v^2$. Subtract the even part of the third coordinate to
write it as $vA(u,v^2)$. Transversality gives $A_u(0,0)\ne0$.
The target change $u\mapsto A(u,Y)$ and the matching source
change give exactly $(u,v)\mapsto(u,v^2,uv)$.
Its image is the crosscap of P4.16, with a compact spherical
class of square $-4$ in its nearby fibre.

The deformed images have holomorphic equations $G_s$ with
$G_0=g$. To see this, a generic projection of the original
finite normalization to two target coordinates is finite;
Weierstrass preparation preserves this after shrinking the
deformation. The total image is an analytic hypersurface
in $\mathbb C^3\times\Delta$. Its central equation has
multiplicity one, as seen at a smooth point of $S$; a unit
adjustment gives $G_0=g$. Outside a smaller ball the original
map is a transverse immersion, and this compact boundary
data persist under small deformation. We therefore have
the boundary property already proved. Transporting the
crosscap class contradicts its torus inclusion. The
normalization is consequently immersive everywhere.

Write its image germ as a unit times $g_1\cdots g_r$, where
each branch is smooth and $v_i=dg_i(0)\ne0$. Absorb the
unit in one factor. Independent gradients give $r\leq3$
and normal-crossing coordinates by the inverse function
theorem. Suppose instead $\sum_i c_iv_i=0$ is a nonzero
relation, and let $I$ be its nonzero support. Set
$a_i=\varepsilon/c_i$ for $i\in I$ and $a_j=\eta$ for
$j\notin I$, where $0<|\varepsilon|\ll|\eta|\ll1$.
Perturb constants so that $h_i(0)=a_i$. Keep all gradients
except one $i_0\in I$, and change that gradient by
\[
 \Delta v_{i_0}=-a_{i_0}\sum_{j\notin I}\frac{v_j}{a_j}.
\]
This is arbitrarily small and makes $\sum_i dh_i(0)/a_i=0$.
Thus $H=\prod_i h_i$ is critical at the origin and has
nonzero critical value $\prod_i a_i$. Its Hessian divided
by that value is
\[
 \sum_i\frac{\operatorname{Hess}(h_i)}{a_i}
 -\sum_i\frac{dh_i\otimes dh_i}{a_i^2}.
\]
An arbitrarily small generic quadratic change to $h_{i_0}$
makes this matrix invertible: its determinant is a nonzero
polynomial in the freely variable symmetric Hessian.
The branch functions remain smooth and retain their
ordinary-double/smooth boundary strata. We have constructed
a holomorphic Morse point off the divisor.

Its nearby fibre contains the real vanishing sphere in
$w_1^2+w_2^2+w_3^2=\epsilon$. Its normal bundle is
$iTS^2$, with the opposite orientation to $TS^2$ under the
complex surface orientation. Its self-intersection is
therefore $-e(S^2)=-2$. Boundary transport again puts this
class in the original torus fibre, a contradiction. Hence
the original branch gradients are independent everywhere,
which proves the normal-crossing assertion.
\end{proof}

The rank-one count is also stated in Gim\'enez Conejero--Pint\'er,
arXiv:2401.14951v2, Lemma 2.3; Section 5.3, p.~33,
discusses the stable-image/nearby-fibre comparison.
The proof above verifies the boundary transport required here.
It does not extend that argument to cuspidal, tacnodal or
triple-curve transverse types.

\paragraph{Verified source locations for this part.}
The classification table used above is the publisher preview of
BHPV, Chapter VI, Theorem VI.1.1 and Table 10, pp.~243--244:
\url{https://page-one.springer.com/pdf/preview/10.1007/978-3-642-57739-0_7}.
The additional exact theorem texts are Ornea--Vuletescu--Verbitsky,
Theorems 4.2 and 5.1, \url{https://arxiv.org/pdf/1810.05768};
Buonerba--Bogomolov--Kurnosov, Theorem 1,
\url{https://arxiv.org/pdf/1709.00062}; Catanese, Theorem 2.1,
\url{https://www.mathe8.uni-bayreuth.de/de/team/prof-fabrizio-catanese-old/pdf/085.pdf};
Inoue 1972, Propositions 1 and 2,
\url{https://www.jstage.jst.go.jp/article/pjab1945/48/7/48_7_445/_pdf};
and Khaled--Teleman, Theorem 3.1 and the proof of Remark 3.4,
\url{https://www.ms.u-tokyo.ac.jp/journal/33-3-1.pdf}.
Mumford's local intersection argument is at
\url{https://www.dam.brown.edu/people/mumford/alg_geom/papers/1961a--TopSingSurface-NumDam.pdf},
Section I, p.~6. All were opened and checked in this run; the
chronological reasons and the claims taken are recorded in the
source log. The proofs of P4.1--P4.7 and the two explicit
counterexamples above were derived here from their charts and
exact maps. No external result for the full arbitrary-fibre sign
assertion is being claimed.


\subsection{The remaining smooth ruled normalizations}

\begin{statement}{P4.21}{Proved}
In the setting of Part~4.2, an irreducible singular support $S$
with smooth rational normalization $T$ has $e(S)>0$.
\end{statement}
\begin{proof}
Let $D$ be the conductor divisor and
$C=\operatorname{Sing}(S)_{\mathrm{red}}$. By P4.12,
$K_T+D\equiv0$. Contract $T$ to a minimal rational surface
$T_0$ and push $D$ forward to $D_0$. At every reversed
blowup the exceptional coefficient is $\mu-1$, where
$\mu$ is the multiplicity of the current conductor at the
centre. In particular every centre lies on the conductor.

We give the complete minimal-surface case split. On
$\mathbb P^2$, the divisor $D_0$ is a plane cubic. An
irreducible cubic is smooth elliptic, rational nodal, or
rational cuspidal. A reducible cubic is a smooth conic
and a line, three lines, a double line and another line,
or a triple line. The conic and line meet at two
transverse points or at one tangency; three distinct
lines form a triangle or are concurrent. Thus either
$D_0$ is reduced nodal and has fundamental group
$\mathbb Z$ or $\mathbb Z^2$, or its reduced support is
simply connected and all normalized components are
rational.

On a Hirzebruch surface, write $F$ for a ruling fibre.
Since $D_0\cdot F=2$, its horizontal part is a double
section, two distinct reduced sections, or a reduced
irreducible bisection. A double section together with
vertical fibres is a rational tree. For the other cases
write
\[
 D_0=H_1+H_2+\sum_j a_jF_j
 \quad\hbox{or}\quad
 D_0=H+\sum_j a_jF_j.
\]
In the first case adjunction on either section gives
\[
 H_1\cdot H_2+\sum_j a_j=2.
\]
The support is therefore a rational tree or a single
rational cycle. A cycle occurs only with all coefficients
one and ordinary double intersection points. This
includes the cases of two section-intersection points,
one section-intersection point and one vertical fibre,
or two vertical fibres. Coincident intersections or a
vertical fibre through the section-intersection point
instead give trees.

For a bisection, adjunction gives
\[
 p_a(H)=1-\sum_j a_j.
\]
There is either no vertical component and $H$ has
arithmetic genus one, or one coefficient-one vertical
fibre and a smooth rational bisection. In the former
case $H$ is smooth elliptic, rational nodal, or rational
cuspidal. In the latter it meets the fibre at two distinct
points or at one tangency. Again the only nontrivial
fundamental groups are $\mathbb Z$ and $\mathbb Z^2$,
and those cases have reduced nodal conductor.

We follow these alternatives through the blowups. At a
point of conductor multiplicity at least two, the new
conductor contains the exceptional $\mathbb P^1$.
Its contraction is the quotient which collapses that
connected simply connected subcomplex to the old
point. Van Kampen therefore makes the induced map on
fundamental groups an isomorphism. At a point of
multiplicity one the conductor is smooth and reduced,
and its strict transform is isomorphic to it. Thus
$\pi_1(D_{\mathrm{red}})$ is unchanged throughout.
If the initial group was nontrivial, the conductor was
reduced nodal. A smooth-point blowup adds no conductor
component, and a node blowup inserts one component of
coefficient one. It remains reduced nodal.

In this latter case the generic transverse conductor
exponent of every normalization branch is one. The
exponent is its intrinsic conductor exponent plus the
intersection multiplicities with the other branches.
A singular branch has intrinsic exponent at least two.
Exponent one therefore means two smooth branches
meeting transversely. Consequently $S$ has ordinary
double curves outside finitely many points. P4.20 makes
$S$ normal crossing everywhere, and P4.6--P4.8 give
$e(S)=2$ for this rational normalization.

In the remaining case $D_{\mathrm{red}}$ is connected
and simply connected, and all its normalized components
are rational. The conductor pushout verified in P4R.3
gives
\[
 \pi_1(S)=\pi_1(T)*_{\pi_1(D_{\mathrm{red}})}\pi_1(C)
          =\pi_1(C).
\]
Every component of $C$ has rational normalization,
because it is the finite nonconstant image of a
rational component of $D_{\mathrm{red}}$. P4R.1 makes
$\pi_1(C)$ virtually abelian. For a connected curve
with rational normalized components, this is the free
group $F_b$ of its point-identification graph, so
$b\le1$. If $c\ge1$ is its number of components, then
\[
 e(C)=c+1-b\ge c\ge1.
\]
P4.13 gives $e(S_{\mathrm{reg}})\ge0$. Additivity proves
$e(S)=e(S_{\mathrm{reg}})+e(C)>0$.
\end{proof}

\begin{statement}{P4.22}{Proved}
In the setting of Part~4.2, an irreducible singular support
$S$ whose smooth normalization $T$ is birationally ruled
over an elliptic curve has $e(S)\ge0$. Equality holds
exactly when $T$ is minimal.
\end{statement}
\begin{proof}
Push the effective conductor to the minimal ruled
surface $T_0\to B$, where $B$ is elliptic. We have
$K_{T_0}+D_0\equiv0$. If a horizontal component $H$
has coefficient one, adjunction gives $p_a(H)\le1$,
whereas its nonconstant map to $B$ gives geometric
genus at least one. Thus $H$ is smooth elliptic and
$(D_0-H)\cdot H=0$. There is no vertical component.
The horizontal part is one elliptic bisection or two
disjoint elliptic sections, all with coefficient one.
All subsequent blowups have smooth coefficient-one
conductor centres; their exceptional curves are absent
from the conductor. The normalization differential
has no divisorial rank drop. P4.15 makes $T$ minimal,
makes the differential immersive everywhere, and
gives $e(S)=0$.

The other possibility is
\[
 D_0=2H+\sum_j a_jF_j,
\]
where $H$ is a smooth elliptic section. As in P4R.3,
$D_{\mathrm{red}}$ is $H$ with rational trees attached,
and its inclusion in $T$ induces an isomorphism on
fundamental groups. The conductor pushout and P4R.1
make $\pi_1(C)$ virtually abelian. The reduced-curve
calculation in P4R.3 gives $e(C)\ge0$, and P4.13 then
gives $e(S)\ge0$.

If $e(S)=0$, the curve calculation forces $C$ to have
one component, genus-one smooth normalization, and
no point-identification graph cycles. Each conductor
component maps finitely and nonconstantly onto $C$.
Riemann--Hurwitz excludes rational conductor
components. Thus $D_{\mathrm{red}}=H$ and no blowup
occurred: the first blowup on $2H$ would create a
vertical conductor component of coefficient one.
Therefore $T$ is minimal.

Conversely suppose $T$ is minimal. In the double-section
case let $v$ be the number of distinct vertical fibres
in the conductor. Since $e(T)=e(H)=0$ and each vertical
$\mathbb P^1$ meets $H$ once,
\[
 e(S_{\mathrm{reg}})=e(T\setminus D_{\mathrm{red}})=-v.
\]
P4.13 forces $v=0$, hence $D=2H$. At a general point
of its image curve, every normalization branch has
conductor exponent two. The branch exponent equals
its intrinsic conductor exponent plus the sum of its
intersections with the other branches. A singular
branch together with another branch would have
exponent greater than two. Thus either there is one
ordinary cusp branch, or all branches are smooth:
two with intersection multiplicity two, or three
with pairwise intersection multiplicity one. Four
smooth branches would have exponent at least three,
so the list is exhaustive.

In the cusp case $H\to C$ has generic degree one.
The normalization of $C$ is therefore $H$, an elliptic
curve. Virtual abelianness of $\pi_1(C)$ excludes any
additional point-identification graph cycle, so
$e(C)=0$. Additivity gives $e(S)=0$. In the other two
cases the normalization differential is immersive at
the generic point of $H$. There is no other conductor
component, so its rank-drop set is finite. P4.15
applies and again gives $e(S)=0$.
\end{proof}


\subsection{Class-VII minimal models with positive second Betti number}

\begin{statement}{P4.23}{Proved}
In the setting of Part~4.2, let an irreducible singular support
$S$ have smooth normalization $T$. If its minimal model $T_0$
belongs to class VII and $b_2(T_0)>0$, then $e(S)>0$.
\end{statement}
\begin{proof}
Let $D$ be the conductor. P4.12 makes $K_T+D$ torsion;
pushing down the point blowups gives $K_{T_0}+D_0$ torsion,
with $D_0$ effective and nonzero. This is a numerically
anticanonical divisor in the sense of Dloussky,
Definition~2.20, arXiv:math/0406387v3. His Theorem~0.2
therefore gives a global spherical shell. Proposition~2.27
and the square-zero-divisor paragraph after Lemma~2.22
leave the parabolic Inoue or Inoue--Hirzebruch cases,
or a connected rational cycle with attached rational
trees. The latter configuration and $\pi_1(T_0)=\mathbb Z$
are stated explicitly in Dloussky--Oeljeklaus--Toma,
\S1, pp.~284--285. We verify the consequences needed
for our conductor.

For an Inoue--Hirzebruch surface, the reduced cycle or
union of two cycles is numerically anticanonical. The
curve intersection matrix is negative definite. Subtracting
that divisor from $D_0$ and intersecting with every curve
therefore makes every coefficient difference zero.
Thus $D_0$ is reduced nodal. In the parabolic Inoue case
there is an elliptic curve $E$ and a disjoint rational
cycle $\Gamma$, with
\[
 K_{T_0}+E+\Gamma\sim0,\qquad \Gamma\equiv0.
\]
These relations are stated after Dloussky's Lemma~2.22.
Intersections with $E$ and the rational curves give
$D_0=E+k\Gamma$: the cycle matrix has kernel generated
by the sum of its components. Torsion of $K_{T_0}+D_0$
then makes $\mathcal O((k-1)\Gamma)$ torsion. A nonzero
effective divisor cannot have torsion associated line
bundle on a compact connected reduced complex space:
a power of its defining section would be a nonzero
constant and could have no zero. Hence $k=1$ and
$D_0$ is again reduced nodal.

The blowup coefficient $\mu-1$ preserves the reduced
nodal property, as in P4.21. Thus $S$ has ordinary
double curves outside finitely many points. P4.20
makes it normal crossing everywhere. Its normalization
is then immersive and P4.6 gives $e(T)=K_T^2$.
For a class-VII surface, however,
\[
 e(T)=b_2(T)>0,\qquad K_T^2=-b_2(T),
\]
a contradiction. These two exceptional minimal-model
cases are excluded.

In the remaining case the support of the conductor is
a connected rational cycle with attached rational trees.
It is the whole connected curve configuration: Dloussky's
Lemma~2.22(2) puts every curve meeting the anticanonical
support in that support, and Proposition~2.27 gives its
connectedness. Its fundamental group is $\mathbb Z$.
The blowups preserve this group by the quotient argument
of P4.21, and preserve the surface fundamental group.
Consequently
\[
 \pi_1(D_{\mathrm{red}})=\mathbb Z,\qquad
 \pi_1(T)=\mathbb Z.
\]
Every normalized component of
$C=\operatorname{Sing}(S)_{\mathrm{red}}$ is rational,
being the finite nonconstant image of a rational conductor
component. The curve is connected. Write $\pi_1(C)=F_b$
for the free group of its point-identification graph.
The conductor pushout gives
\[
 \pi_1(S)=
 \langle t,c_1,\ldots,c_b\mid t^k=w(c_1,\ldots,c_b)\rangle.
\]
Here $k$ records the conductor-cycle map to $\pi_1(T)$,
and $w$ records its map to $\pi_1(C)$. Neither map is
assumed to be an isomorphism.

Virtual abelianness from P4R.1 forces $b\le1$.
If $k\ne0$ and $w\ne1$, both cyclic amalgam maps are
injective; alternating nonidentity coset representatives
give the reduced normal form and show that $F_b$ injects
in the amalgam. If $w=1$, the group is
$(\mathbb Z/k\mathbb Z)*F_b$, with
$\mathbb Z/0\mathbb Z=\mathbb Z$, and $F_b$ again
injects. A nonabelian free subgroup is impossible in a
virtually abelian group. Finally, if $k=0$ and $w\ne1$,
the group is
\[
 \mathbb Z*\bigl(F_b/\langle\!\langle w\rangle\!\rangle\bigr).
\]
A free product of an infinite cyclic group and a nontrivial
group contains a nonabelian free subgroup: for a nontrivial
$g$ in the second factor, $t$ and $gtg^{-1}$ have no
nonempty reduced relation. Thus the second factor must
be trivial. Its abelianization has rank at least $b-1$,
which again gives $b\le1$.

For the number $c\ge1$ of components of $C$, we obtain
\[
 e(C)=c+1-b\ge c\ge1.
\]
By P4.13, $e(S_{\mathrm{reg}})\ge0$, so
$e(S)=e(S_{\mathrm{reg}})+e(C)>0$.
\end{proof}

\subsection{Minimality of a smooth Hopf normalization}

\begin{statement}{P4.24}{Proved}
In the setting of Part~4.2, let $S$ be an irreducible
singular fibre support with smooth normalization $T$.
If the minimal model of $T$ is a Hopf surface, then
$T$ is already minimal. The reduced conductor on $T$
and its reduced image in $S$ are disjoint unions of
smooth elliptic curves. Each map between their
components is a finite unramified covering, and
$e(S)=0$.
\end{statement}
\begin{proof}
Write $\nu:T\to S\hookrightarrow Y$ for the finite
holomorphic normalization map followed by the closed
embedding. We use the effective integral conductor
divisor $D$ with $K_T+D$ torsion. Let $T\to T_0$
be the sequence of $k$ point blowups giving the
minimal Hopf model. At each step the conductor
coefficient on the exceptional curve is $\mu-1$,
where $\mu$ is the multiplicity of the conductor
at the centre. Effectiveness therefore puts every
centre on the current conductor.

There is a finite unramified primary Hopf cover
$T'_0\to T_0$. Pulling back every blowup gives a
finite unramified cover $T'\to T$ and a blowup
presentation of $T'$ over $T'_0$. A centre has
exactly the degree of the covering many lifts,
so the new number of blowups is that degree times
$k$. The composite $T'\to Y$ is finite and is
immersive away from the pulled-back conductor.
The rational Chern classes of its pulled-back
ambient tangent bundle vanish. Indeed $b_2(Y)=0$
gives $c_1(T_Y)=c_2(T_Y)=0$ in rational cohomology,
using Poincar\'e duality in degree four. The
argument below applies to this finite map, without
requiring it to be a normalization. Proving
minimality of $T'$ proves $k=0$. We therefore use
$T,T_0,\nu,k$ for the covering objects during the
Chern calculation.

We first construct on $T_0$ a nonsingular
holomorphic foliation tangent to every compact
curve. Khaled, \emph{Real structures on primary
Hopf surfaces}, Theorem~2.1(1) and formula~(1),
gives
\[
 T_0=(\mathbb C^2\setminus\{0\})/\langle F\rangle,
 \qquad F(z,w)=(\alpha z+\lambda w^n,\beta w),
\]
where $0<|\alpha|\le|\beta|<1$ and
$\lambda(\alpha-\beta^n)=0$. We verify the needed
curve statement from these coordinates. Every
compact irreducible curve is smooth elliptic:
adjunction and $b_2(T_0)=0$ give arithmetic genus
one, and a rational normalization would lift to
$\mathbb C^2\setminus\{0\}$ and have constant
coordinate functions.

The map from the fundamental group of such an
elliptic curve to $\pi_1(T_0)=\mathbb Z$ has
nonzero image, since zero image would give the
same impossible compact lift. Its full inverse
image has finitely many connected components.
The Remmert--Stein extension across the origin
is an analytic curve germ; a power $F^h$ fixes
one of its branches. On the normalization of
that branch the induced contraction can be
written $t\mapsto\eta t$, $0<|\eta|<1$.
For completeness, if the initial coordinate
has return map $g(t)=\eta t+O(t^2)$, choose
$|\eta|<\rho<\sqrt{|\eta|}$. On a sufficiently
small disc, $|g(t)|\le\rho|t|$, and the
successive differences in $g^j(t)/\eta^j$
are bounded by a constant times
$(\rho^2/|\eta|)^j|t|^2$. Their convergent sum
gives the asserted linearizing coordinate.

When $\lambda=0$, comparison of coefficients in
$z(\eta t)=\alpha^h z(t)$ and
$w(\eta t)=\beta^h w(t)$ shows that each
nonzero coordinate is a single monomial.
Thus every branch away from the coordinate axes
has parametrization $(a t^p,b t^q)$. Its ratio
$p/q=\log|\alpha|/\log|\beta|$ is independent
of the branch. If such a branch exists, take
coprime positive integers $p,q$ in this ratio
and the vector field
$p z\partial_z+q w\partial_w$.
If no such branch exists, use
$z\partial_z+w\partial_w$.
Both are invariant under $F$, have no zero on
$\mathbb C^2\setminus\{0\}$, and are tangent
to every compact curve.

When $\lambda\ne0$, we have $\alpha=\beta^n$ and
\[
 F^h(z,w)=
 (\alpha^h z+h\lambda\alpha^{h-1}w^n,\beta^h w).
\]
If a normalized invariant branch had
$w(t)=b t^q\ne0$, the coefficient of $t^{nq}$
in the first coordinate identity would give
$0=h\lambda\alpha^{h-1}b^n$, a contradiction.
Thus the only curve is the image of $w=0$.
The vector field $n z\partial_z+w\partial_w$
is invariant, nowhere zero, and tangent to it.
This proves the foliation assertion.

Pull this foliation through the point blowups
and saturate its tangent sheaf. It is a line
bundle $L\subset T_T$. Every strict transform
of a conductor component and every exceptional
curve is invariant. The only singularities of
the foliation are simple saddles. To verify
this assertion and its multiplicities, at a
regular point use coordinates $(x,y)$ with
leaves $y=\text{constant}$. In the chart
$x=u,y=uv$ the saturated generator is
$u\partial_u-v\partial_v$; there is one
saddle, and the other blowup chart is regular
away from it. At an existing saddle use
$b x\partial_x-a y\partial_y$, with positive
integers $a,b$. The two charts have generators
\[
 b u\partial_u-(a+b)v\partial_v,\qquad
 (a+b)u\partial_u-a v\partial_v.
\]
One saddle is replaced by two. Each has
zero ideal $(u,v)$ and length one.
Consequently there are exactly $k$ saddles.

Let $V$ be the divisor measuring the
divisorial zero of the pulled-back transverse
one-form before saturation. Numerically
\[
 c_1(L)=-K_T+V.
\]
The omitted pullback line bundle is topologically
trivial because it comes from $T_0$.
The preceding charts give, at every blowup $\mu$,
\[
 K'=\mu^*K+E,\qquad
 V'=\mu^*V+\varepsilon E,\qquad
 \varepsilon=
 \begin{cases}0&\text{at a regular point},\\
 1&\text{at a saddle}.
 \end{cases}
\]
Starting with $V=0$ and $K_{T_0}^2=0$, we obtain
\[
 K_T^2=-k,\qquad K_TV-V^2=0.
\]
The second identity is preserved because its
change is $\varepsilon^2-\varepsilon=0$.

Put $B=\nu^*T_Y$ and compose the holomorphic
maps $L\to T_T\to B$. This map has no
divisorial zero. Away from the conductor this
follows from immersivity. At the generic point
of a conductor component the foliation is
tangent to that curve, and finiteness of
$\nu$ makes its differential along the curve
nonzero. A nonconstant holomorphic map of
curves has nonzero derivative at the generic
point. Thus the cokernel $Q=B/L$ is a
torsionfree sheaf of rank two. Its double
dual $Q^{**}$ is locally free, and
\[
 z_0=\operatorname{length}(Q^{**}/Q)
\]
is finite. At a point, write the map $L\to B$
as a column $(f_1,f_2,f_3)$ in trivializations.
If this column vanishes, the ideal
$(f_1,f_2,f_3)$ has finite colength. Dualizing
$0\to\mathcal O\to\mathcal O^3\to Q\to0$
gives $\mathcal E xt^1(Q,\mathcal O)
=\mathcal O/(f_1,f_2,f_3)$.
Dualizing $0\to Q\to Q^{**}\to Q^{**}/Q\to0$
and using the two-dimensional regular local
ring identifies this with
$\mathcal E xt^2(Q^{**}/Q,\mathcal O)$,
which has the same length. At each foliation
saddle all three coefficients vanish.
It follows that
\[
 z_0\ge k.
\]

Let $E_B\subset B$ be the saturation of the
generic rank-two image of $d\nu:T_T\to B$.
Its determinant differs from that of $T_T$
by an effective integral divisor $R$, the
divisorial rank-drop divisor:
\[
 c_1(E_B)=-K_T+R.
\]
The support of $R$ is contained in the
conductor. Saturate $E_B/L$ in $Q^{**}$.
The resulting rank-one reflexive subsheaf
is a line bundle $J$, with
$c_1(J)=R-V$. Its torsionfree rank-one
quotient has the form
$\mathcal I_{Z_1}\otimes J_1$ with
$c_1(J_1)=K_T-R$ and
$z_1=\operatorname{length}(Z_1)\ge0$.
All these assertions concern sheaves on the
smooth surface $T$; determinants are computed
off finite sets and therefore determine their
line-bundle classes.

The identities $c_1(B)=c_2(B)=0$ give
\[
 c_2(Q)=(-K_T+V)^2,\qquad
 c_2(Q^{**})=(-K_T+V)^2-z_0.
\]
The exact sequence with $J$ gives another
expression for the same number:
\[
 (-K_T+V)^2-z_0=(R-V)(K_T-R)+z_1.
\]
Using $K_T^2=-k$ and $K_TV-V^2=0$, this becomes
\[
 G(R):=R^2-K_TR-VR=k+z_0+z_1\ge2k.
\]

We compare this expression with the actual
blowup sequence. Push $R$ and $V$ down at
each step; $R$ has integral coefficients.
Writing $R'=\mu^*R+cE$ gives $c\in\mathbb Z$,
without any positivity requirement on $c$.
The intersection calculation is
\[
 G(R')-G(R)=-c^2+(1+\varepsilon)c.
\]
For $\varepsilon=0$ this is at most zero
for every integer $c$; for $\varepsilon=1$
it is $1-(c-1)^2\le1$. On $T_0$ every
divisor is numerically zero, so $G(R_0)=0$.
Consequently
\[
 G(R)\le\#\{\text{blowups at saddles}\}\le k.
\]
The lower bound $G(R)\ge2k$ forces
$k=z_0=z_1=0$. Returning through the finite
unramified cover proves that the original
normalization $T$ is minimal.

It remains to identify the conductor maps,
rather than only their Euler number.
The reduced conductor on a minimal Hopf
surface consists of disjoint smooth elliptic
curves. Since $z_0=0$ on the primary cover,
the differential of $\nu$ along each such
curve is nowhere zero, also before that
cover. Thus every branch of its image is
smooth. At an image point, two different
curve germs would have to come from two
different normalization points. Their
surface branches in the smooth threefold
intersect in a curve germ. This germ is
contained in the conductor image of both
branches; the unique conductor germ at
each normalization point forces the two
image germs to be equal. Therefore the
image conductor is locally a single smooth
curve. Its components are disjoint, and
each conductor map is finite unramified.
Riemann--Hurwitz makes every image component
elliptic. Finally, normalization is an
isomorphism off these curves, so
\[
 e(S)=e(T)-e(D_{\mathrm{red}})
                 +e(\operatorname{Sing}S)=0.
\]
\end{proof}


\subsection{The reduced resolved-conductor case}

\begin{statement}{P4.25}{Proved}
Every semi-log-canonical reduced fibre support in
Part~4.2 has normal crossings. In particular, let
$\rho:T\to N\to S$ be a resolution of the normalization
of an irreducible support, and let $D$ be the integral
divisor defined by
$\omega_T(D)\simeq\rho^*\omega_S$. When this actual
divisor is reduced and has normal crossings, $S$
has normal crossings, its normalization is smooth,
and P4.6--P4.8 apply.
\end{statement}
\begin{proof}
We first explain why the stated resolved-conductor
case is semi-log-canonical. At a generic conductor
point, an exponent one on every normalization
branch forces the transverse plane curve to be an
ordinary node. Indeed the conductor exponent on
branch $i$ is
$2\delta_i+\sum_{j\ne i}I_{ij}$; positivity of
intersection multiplicities leaves exactly two
smooth transverse branches. The support $S$ is a
Cartier hypersurface, hence is Cohen--Macaulay and
has invertible dualizing sheaf. On its normal
normalization $N$, the pair consisting of $N$
and its reduced conductor is log canonical.
To check this last assertion, pull its canonical
divisor to $T$. The resulting boundary is exactly
$D$, by the displayed isomorphism and the
divisorial orders of a local generator. A smooth
surface with reduced normal-crossing boundary is
log canonical: blowing up a point gives new
boundary coefficient $\mu-1\le1$, and this
bound persists under successive point blowups.
Discrepancy comparison therefore gives the
assertion on $N$. This is precisely the
semi-log-canonical definition of
Liu--Rollenske, Definition~2.3 and Lemma~2.2,
which explicitly applies to complex analytic
spaces.

We now treat all local semi-log-canonical
hypersurface germs. We use the complete analytic
classification in Liu--Rollenske,
\emph{Two-dimensional semi-log-canonical
hypersurfaces}, main Theorem, Table~1,
Propositions~3.1 and~4.1. The normal singular
germs are isolated. Any such germ is excluded
in our family: a small generic linear
perturbation of its reduced equation has
Morse critical points, the number counted
with multiplicity being the positive
colength of its gradient ideal. The equality
of this number before and after perturbation
follows by conservation of the local
intersection multiplicity of the three
gradient equations. A constant perturbation
puts their critical values away from zero.
The original zero divisor is smooth near
a small transverse boundary sphere, so it
retains that boundary property under these
perturbations. The proper boundary
transport proved in P4.20 carries a
vanishing sphere of square $-2$ into a
nearby torus fibre, which is impossible.

For the nonnormal germs we give the exact
deformations and check their boundary strata.
The ordinary double point and $xyz=0$
already have normal crossings. A pinch point
$x^2+y^2z=0$ is excluded by P4.16.
The remaining double-point equations are
\[
 x^2+y^2z^2,\qquad
 x^2+y^2(z^2+y^{q-2}),\quad q\ge3.
\]
For the first equation the two normalization
branches are the smooth graphs $x=\pm i yz$.
Their tangent planes coincide at the origin,
whereas they meet transversely off that point
along the two axes. The local proof of P4.20
therefore supplies a vanishing sphere after
a deformation retaining its boundary
ordinary-double strata.

For finite $q$ put
\[
 G_s=x^2+y^2(z^2+y^{q-2}+s).
\]
For $s\ne0$ small, at each point
$(0,0,z_0)$ with $z_0^2=-s$ the function
$Z=z^2+y^{q-2}+s$ is a local coordinate:
its $z$ derivative is $2z_0\ne0$.
The actual equation there is exactly
$x^2+y^2Z$, a pinch point. At the original
boundary double axis $x=y=0$, $z\ne0$,
the transverse Hessian determinant is
$4(z^2+s)$, which stays nonzero on the
compact boundary when $s$ is small.
All boundary smooth points remain smooth.
Thus the zero divisor retains the same
boundary smooth and double strata.
The spherical class of square $-4$
from the pinch-point nearby fibre transports,
by the detailed P4.20 construction, to
the original nearby torus fibre.
This excludes every finite $q$.

The remaining triple-point equations are
\[
 xyz+x^p,\quad p\ge3,\qquad
 xyz+x^p+y^q,\quad q\ge p\ge3.
\]
For the first put
$G_s=xyz+x^p+s x^2
=x(yz+x^{p-1}+s x)$.
When $s\ne0$, its two local normalization
branches are smooth: the derivative in $x$
of the second factor at the origin is $s$.
Their tangent planes coincide, so the
origin is not a normal crossing. Along
the boundary $y$ axis the transverse
Hessian determinant is $-y^2$, and along
the boundary $z$ axis it is $-z^2$.
Both axes persist as ordinary-double
strata. The local argument of P4.20 on
the deformed germ supplies a compact
spherical class of nonzero square, and
the same boundary transport gives the
contradiction.

For the second equation put
\[
 G_s=xyz+x^p+y^q+s x^2+s y^3,
 \qquad 0<|s|\ll1.
\]
Its transverse Hessian along the boundary
$z$ axis has determinant $-z^2$,
so its sole boundary double stratum
persists. We check the smooth normalization
of its germ at the origin explicitly.
Since $\partial^2G_s/\partial x^2(0)=2s$,
the critical equation in $x$ has a unique
analytic solution $x=\xi(y,z)$ with
$\xi$ divisible by $yz$. The parametric
holomorphic Morse change in $x$ writes
\[
 G_s=X^2+y^2H(y,z).
\]
Divisibility follows directly on substituting
$\xi$: each of $\xi yz,\xi^p,s\xi^2$
is divisible by $y^2$, as are the two
pure $y$ terms. Moreover
\[
 H_y(0,0)=
 \begin{cases}1+s&q=3,\\s&q>3,\end{cases}
 \qquad H(0,z)=-z^2/(4s).
\]
The finite
birational map
\[
 \{v^2+H(y,z)=0\}\longrightarrow\{G_s=0\},
 \qquad (v,y,z)\longmapsto(X=yv,y,z),
\]
followed by the inverse Morse coordinate
change, is the normalization. Its source
is smooth because $H_y(0,0)\ne0$,
and the map is finite because $v$
satisfies the displayed monic equation.
It is birational on $y\ne0$ and is
therefore the normalization by normality
of its source. At the origin its
differential has rank one, so this germ
is not a normal crossing. P4.20 again
produces a spherical class, with the
already verified boundary transport.

These calculations exhaust the analytic
classification. Thus only the smooth,
ordinary-double, and triple-normal-crossing
germs remain. Their local equations give
the claimed normal-crossing support and
smooth normalization.
\end{proof}


\subsection{Numerical sign for arbitrary irreducible supports}

\begin{statement}{P4.26}{Proved}
In Part~4.2 every irreducible reduced fibre support
has nonnegative topological Euler number. For a
singular irreducible support, rational resolution
of the normalization gives strictly positive
Euler number; so does a resolution whose minimal
model is of class VII with positive second
Betti number. These assertions permit singular
normalization.

Every normal-crossing support has the
nonnegative even Euler number of P4.6.
Consequently, for a fibration all of whose
reduced supports are irreducible or have
normal crossings,
\[
 B=\mathbb P^1,\qquad c_3(Y)\in\{0,2\},
 \qquad b_3(Y)=2-c_3(Y).
\]
This last conclusion does not assume a
classification of the equality cases for
irreducible nonnormal supports.
\end{statement}
\begin{proof}
For an irreducible singular support use the
minimal resolution of its normalization
$\rho:T\to N\to S$, and put
$A=\rho^{-1}(\operatorname{Sing}S)_{\rm red}$,
$C=\operatorname{Sing}(S)_{\rm red}$.
P4R.6 gives an effective integral divisor
$D$ with support $A$ and $K_T+D$ torsion.
The resolution $T$ has Kodaira dimension
$-\infty$, by P4.12. We explain the two
remaining transfers from smooth
normalization to this resolution.

Suppose first that $T$ is rational.
Push $D$ to a minimal rational surface.
The complete anticanonical-support calculation
in P4.21 has the following exact alternatives.
The reduced support is connected; either
its fundamental group is trivial and all
its normalized components are rational,
or the divisor itself is reduced with
normal crossings. On $\mathbb P^2$ these
are the plane cubic configurations
enumerated there. On a ruled minimal
surface they are the degree-two horizontal
configurations there. Those arguments use
only the effective integral divisor and
$K+D\equiv0$, so apply to the present
resolved divisor.

We check their preservation through the
actual blowup morphisms. A blowup at
conductor multiplicity $\mu$ gives
exceptional coefficient $\mu-1$.
For $\mu=1$, the centre is a smooth
point of the reduced support and its
strict transform is isomorphic to that
support. For $\mu\ge2$, the whole
exceptional $\mathbb P^1$ belongs to
the new support. Its contraction is
the blowup map on supports. Collapsing
this connected simply connected
subcomplex does not change the
fundamental group, by van Kampen.
The reduced normal-crossing case stays
reduced normal crossing: a smooth
point gives coefficient zero and a
double point gives coefficient one.
Thus the same alternatives hold for
the actual divisor $D$ on $T$.

In the reduced normal-crossing alternative,
P4.25 makes $S$ normal crossing. Its
normalization is smooth, so this minimal
normalization resolution is the
normalization itself. The rational case
of P4.6--P4.8 gives $e(S)=2$.
In the other alternative $\pi_1(A)=1$
and $\pi_1(T)=1$. The proper-map pushout
$T\cup_A C\to S$ is a homeomorphism:
it is a continuous compact-to-Hausdorff
bijection, since $\rho$ is an isomorphism
off $A$ and $A$ is the full inverse
image of $C$. Compatible triangulation
makes $A\to T$ a cofibration.
Van Kampen therefore gives
$\pi_1(S)\simeq\pi_1(C)$.
Every normalized component of $C$ is
rational: a nonconstant map from a
normalized component of $A$ onto it
is a finite map from $\mathbb P^1$.
Contracted components of $A$ contribute
no exception to this statement.
The curve $C$ is connected, and its
fundamental group is a free group
$F_b$ on the cycles of its
point-identification graph. P4R.1
implies $b\le1$. If $c\ge1$ is the
number of its components, then
\[
 e(C)=c+1-b\ge1.
\]
P4.13 gives $e(S_{\rm reg})\ge0$,
so $e(S)>0$.

Now suppose that the minimal model
$T_0$ is of class VII with
$b_2(T_0)>0$. The divisor pushed to
$T_0$ satisfies the hypotheses of
Dloussky's numerically anticanonical
theorem exactly as in P4.23.
The verified sources and their
normalization-independent hypotheses
are specified there. The
Inoue--Hirzebruch and parabolic Inoue
alternatives give a reduced
normal-crossing divisor, preserved
under the blowups just checked.
P4.25 then makes $S$ normal crossing
and $T=N$ smooth. The immersive
normalization Chern formula would
give $e(T)=K_T^2$, whereas
$e(T)=b_2(T)>0$ and
$K_T^2=-b_2(T)$. These alternatives
cannot occur.

In the intermediate alternative
$A$ is a connected rational cycle
with attached rational trees, and
\[
 \pi_1(A)=\mathbb Z,\qquad
 \pi_1(T)=\mathbb Z.
\]
Both statements survive the actual
blowups; they concern the resolution
and do not require $N$ to be smooth.
Every normalized component of $C$
is again rational and $C$ is connected.
With $\pi_1(C)=F_b$, the pushout
presentation is exactly
\[
 \pi_1(S)=
 \langle t,c_1,\ldots,c_b
 \mid t^k=w(c_1,\ldots,c_b)\rangle
\]
for the actual two conductor-group
homomorphisms. A component contracted
to a point merely induces the zero
homomorphism on its group. The
amalgam and free-product normal-form
argument in P4.23, which allowed
arbitrary $k$ and $w$, gives $b\le1$
from virtual abelianness.
Thus $e(C)=c+1-b\ge1$, and
$e(S)>0$ follows again from P4.13.

This exhausts these two resolution
classes. Ruled resolutions of positive
base genus are covered by P4R.7
and P4R.9, and Hopf minimal models
by P4R.8. The class-VII minimal
model with $b_2=0$ is Hopf or Inoue,
by the verified classification in
P4.8. An Inoue model has no curves,
so it cannot carry the nonempty
effective divisor pushed down from
$D$. Therefore it does not occur
here. All irreducible singular
supports have now been covered.

A normal irreducible support is
smooth by P4.12. The cyclic
base-change comparison and torus
deformation theorem in P4.8 show
that it has a finite unramified
torus cover, hence Euler number
zero. This proves the full
irreducible assertion.

For the final global statement,
$B=\mathbb P^1$ follows from
$b_1(Y)=0$ and connected fibres,
as proved in P4.7. Euler
additivity over the finitely many
exceptional base values gives
\[
 e(Y)=\sum_{b\in B}
 e\bigl((f^{-1}(b))_{\rm red}\bigr),
\]
because a smooth torus has Euler
number zero and the complement
is a proper smooth fibre bundle.
Every summand is nonnegative
under the stated support
hypothesis. Poincar\'e duality
and $b_1=b_2=0$ give
$e(Y)=2-b_3(Y)$; the nondegenerate
skew-symmetric cup-product
pairing on $H^3(Y;\mathbb Q)$
makes $b_3(Y)$ even. Thus
$0\le e(Y)\le2$ and $e(Y)$ is
even, which proves the asserted
values and their Chern-number
interpretation.
\end{proof}


\subsection{Equality in the Hopf-resolution case}

\begin{statement}{P4.27}{Proved}
Let \(S\) be an irreducible singular reduced fibre support in Part 4.2.
Suppose that the minimal resolution \(T\to N\) of its normalization has a
Hopf minimal model \(T_0\). Then
\[
 e(S)=0\quad\Longleftrightarrow\quad
 N\text{ is a smooth minimal Hopf surface}.
\]
In this case the reduced conductor and its image are disjoint unions of
smooth elliptic curves, and all component maps are finite unramified covers.
\end{statement}
\begin{proof}
Write \(\rho:T\to N\to S\), \(A=\rho^{-1}(\operatorname{Sing}S)_{\rm red}\),
and \(C=\operatorname{Sing}(S)_{\rm red}\). Let \(k\) be the number of
blowups in \(\pi:T\to T_0\). P4R.6 and P4R.8 give an effective integral pole
divisor whose reduced support \(A\) has \(r\) connected components. Each is
a smooth elliptic root with a rational tree attached. Let \(v\) be the
number of its rational components. Every such component is one of the
\(k\) blowup exceptional curves. We have
\[
 e(S)=k-v+e(C),\qquad k-v\ge0,\qquad e(C)\ge0.
\]
Thus equality forces \(k=v\) and \(e(C)=0\).

Let \(c,s,g,b\) denote the numbers of irreducible components of \(C\),
connected components of \(C\), elliptic normalized components of \(C\), and
the sum of its point-identification graph cycle ranks. P4R.8 gives
\[
 e(C)=c+s-2g-b,\qquad
 b+r-s\le1,\qquad c\ge g,\quad r\ge g.
\]
Consequently \(e(C)=0\) implies \(c+r-2g\le1\).

First suppose \(c=g\). Every normalized component of \(C\) is elliptic.
Every rational component of \(A\) is contracted by \(T\to N\), since a
finite nonconstant map from \(\mathbb P^1\) to an elliptic normalization
does not exist. A root cannot be contracted. Indeed all rational curves
attached to it are already contracted, so its whole rooted connected
component would be a connected exceptional block disjoint from the strict
conductor. P4R.6 excludes precisely such a block. Hence all roots are
noncontracted.

The equality \(k=v\) now says that every blowup exceptional curve of
\(\pi:T\to T_0\) belongs to \(A\) and is contracted by \(T\to N\).
Conversely every exceptional curve of \(T\to N\) belongs to \(A\), by
P4R.6, and cannot be a root. The two proper birational maps therefore have
exactly the same exceptional curves. Their fibres are identical: each
nontrivial fibre is a connected component of the exceptional locus, and
their fibres are connected by Stein factorization and normality of the
targets.

Here is the holomorphic comparison, rather than only a topological
identification. Take the image of
\[
 (\pi,\rho_N):T\longrightarrow T_0\times N
\]
and then its normalization \(G\). The two projections \(G\to T_0,N\)
are proper birational maps with zero-dimensional fibres, because the
original two maps have identical fibres. They are finite. A finite
birational morphism to a normal complex space is an isomorphism, as follows
on local rings from integral closedness. Thus \(N\simeq T_0\)
biholomorphically. It is a smooth minimal Hopf surface.

The only other possibility is \(c=g+1,r=g\). There is one rational
normalized component of \(C\), and all \(g=r\) roots are noncontracted and
map to distinct elliptic normalized components of \(C\). Each such map is
an isogeny. Also
\[
 b+r-s=1.
\]
We exclude this possibility by a quotient of the exact fundamental-group
pushout; the quotient is constructed explicitly below.

Represent \(C\) as its disjoint normalized smooth curves with a finite
graph of point identifications. This is a graph of spaces with elliptic
vertex groups \(\mathbb Z^2\), trivial rational vertex groups, and trivial
edge groups. Its graph has \(s\) connected components and total cycle rank
\(b\). Add a vertex for \(T\), and add one root edge for each connected
component of \(A\). The root edge group is \(\mathbb Z^2\), since its
rational tree is simply connected. Its map to the \(C\) side factors
through the elliptic normalization receiving that root; its map to the
\(T\) side is the actual inclusion. This graph-of-spaces presentation is
the ordinary van Kampen presentation of the homotopy pushout
\(T\cup_A C=S\). It has total graph cycle rank \(b+r-s=1\).

Let \(\Gamma=\pi_1(T)=\pi_1(T_0)\). It is virtually infinite cyclic, and
\(b_1(T)=1\). The primitive map from its abelianization to its free part
therefore gives a surjection
\[
 h:\Gamma\longrightarrow\mathbb Z
\]
with finite kernel. Indeed a cyclic subgroup of finite index maps
nontrivially to \(\mathbb Z\), so the kernel meets that subgroup trivially
and is finite.

For root \(H_i\), write
\[
 a_i:\pi_1(H_i)=\mathbb Z^2\longrightarrow\Gamma,\qquad
 q_i:\pi_1(H_i)=\mathbb Z^2\longrightarrow
       \pi_1(\widetilde C_i)=\mathbb Z^2.
\]
The second map is injective of finite index because it is induced by an
isogeny. The composite \(h a_i\) is nonzero. To see this, a finite image
of \(a_i\) would become trivial after a finite cover of \(H_i\); that
compact elliptic cover would lift holomorphically to the universal Hopf
cover \(\mathbb C^2\setminus\{0\}\), where its coordinate functions are
constant. Thus \(a_i\) has infinite image; since \(\ker h\) is finite,
\(h a_i\ne0\).

Put \(K_i=\ker(h a_i)\), a primitive rank-one subgroup of
\(\mathbb Z^2\). Quotient the \(i\)-th elliptic vertex group by the
saturation of \(q_i(K_i)\). The resulting vertex group is \(\mathbb Z\).
Quotient the root edge group by \(K_i\); it too becomes \(\mathbb Z\).
Its two induced maps into the \(T\) vertex \(\mathbb Z\) and the elliptic
vertex \(\mathbb Z\) are nonzero multiplications and hence injective.
The latter assertion follows because an element outside \(K_i\) cannot
map under the injective rational map \(q_i\otimes\mathbb Q\) into the
rational span of \(q_i(K_i)\). This verifies compatibility of the chosen
quotients with both edge homomorphisms.

We have constructed a surjective homomorphism from \(\pi_1(S)\) onto the
fundamental group \(G\) of a graph of groups. Its nontrivial vertex groups
are infinite cyclic, its root edges are infinite cyclic with injective
maps, and all point-identification edges are trivial. The root edges
form a star: the roots map to distinct elliptic normalizations. They
therefore contain no graph cycle. The unique graph cycle of the complete
graph contains a point-identification edge, whose edge group is trivial.

Delete one such edge on the cycle. The remaining graph is a connected
tree, and the normal form for iterated amalgamated products along
injective edge maps embeds its \(T\) vertex group \(\mathbb Z\) into its
fundamental group \(G_0\). Restoring the removed edge is an HNN extension
over the trivial group, hence
\[
 G\simeq G_0*\mathbb Z.
\]
Choosing a nontrivial element of the embedded \(\mathbb Z\subset G_0\)
and the new stable letter gives an embedded free group \(F_2\), by reduced
free-product normal form. But every quotient of a virtually abelian group
is virtually abelian, and a virtually abelian group has no \(F_2\)
subgroup. This contradicts P4R.1. The case \(c=g+1,r=g\) is impossible.

We have proved that equality makes \(N\) smooth minimal Hopf. Conversely
P4.24 applies to such an actual normalization and proves the unramified
elliptic-conductor description and \(e(S)=0\). This completes both
directions.
\end{proof}


\subsection{Cuspidal transverse type and the ruled equality case}

\begin{statement}{P4.28}{Proved}
Let $p$ lie on a reduced fibre support of Part~4.2.
Suppose its reduced singular locus is a smooth curve
near $p$, and its generic transverse plane-curve germ
is the cusp $a^2+b^3=0$. Then the actual germ at $p$
is analytically isomorphic to
$\{a^2+b^3=0\}\times(\mathbb C,0)$.
In particular its normalization is smooth.
\end{statement}
\begin{proof}
We extend the boundary argument of P4.20 in the
precise form used here. Suppose a deformation retains
the smooth critical curve near a fixed transverse
boundary sphere and retains transverse type $A_2$.
Near a boundary point, the holomorphic Morse lemma
in the nondegenerate transverse variable followed
by extraction of a cube root of a unit gives
coordinates with defining function $a^2+b^3$.
The third coordinate parametrizes the critical
curve. Its transversality to the boundary lets
one real third-coordinate variable be solved
from the boundary equation; the complex
coordinates $a,b$ remain free. Thus the
restriction to the boundary is submersive
at every nonzero level. Compactness and
the smooth part of the boundary give the
same uniform nonzero-level submersion as
in P4.20. Its finite-critical-value and
proper Ehresmann transport proof applies
unchanged. A compact spherical class of
nonzero square in a deformed nearby
fibre would therefore lie in an actual
nearby torus fibre, which is impossible.
All deformations below retain this
boundary transverse type.

Choose holomorphic coordinates $(x,y,z)$
so the smooth critical curve is $x=y=0$.
The transverse Hessian along it is a
symmetric rank-one matrix over the
discrete valuation ring $\mathbb C\{z\}$.
It is not the zero matrix generically.
It can be diagonalized holomorphically
to $h(z)x^2$. Here is the elementary
ring calculation. The determinant is
zero. Choose a diagonal entry of
minimal order; one exists since the
square of the off-diagonal entry is
the product of the two diagonal entries.
The ratio of the off-diagonal entry
to the chosen entry is holomorphic.
Completing its square gives the desired
invertible coordinate change. The
coefficient $h$ is not identically zero.

First suppose that the Hessian has
rank one at $p$. The Morse lemma with
parameters $(y,z)$ writes the reduced
equation as
\[
 g=X^2+y^3h_1(y,z).
\]
Indeed after removing the $X$ direction,
the remaining function and its first
two $y$ derivatives vanish on the
critical curve; the third derivative
is not identically zero there. If
$h_1(0,0)\ne0$, extracting its cube
root in the $y$ coordinate proves the
asserted product form.

Otherwise $h_1(0,z)$ has a zero of
finite positive order. Use
$g_\eta=X^2+y^3(h_1+\eta)$.
For arbitrarily small generic $\eta$,
the function $h_1(0,z)+\eta$ has a
simple zero $z_\eta$ near zero.
This follows directly from one-variable
Weierstrass preparation, avoiding the
finite critical values of $h_1(0,z)$.
At $(0,0,z_\eta)$ the function
$Z=h_1(y,z)+\eta$ is a coordinate,
since its $z$ derivative is nonzero.
The actual germ is
$X^2+y^3Z$, which is excluded by
P4.16: its nearby fibre has a
spherical class of square $-6$.
On a fixed sufficiently small boundary,
$h_1(0,z)$ is nonzero at the critical
curve, and this persists for small
$\eta$. The deformed transverse type
is therefore still $A_2$ there.
The preceding boundary transport gives
the contradiction.

It remains to exclude Hessian rank
zero at $p$, so $h(0)=0$. In the
diagonalizing coordinates add
$\varepsilon y^3$ to $g$. The
critical curve and its transverse
Hessian do not change. The cubic
coefficient in the transverse kernel
remains nonzero near the original
compact boundary for small
$\varepsilon$, so this deformation
has the boundary property just proved.
Choose $\varepsilon$ so that the
coefficient of $y^3$ at the origin
is nonzero.

On a smaller neighbourhood, Weierstrass
preparation in $y$ gives a unit times
\[
 P=y^3+a(x,z)y^2+b(x,z)y+c(x,z).
\]
Restriction to $x=0$ is a unit times
$y^3$, so all three coefficients
vanish at $x=0$. The critical-curve
equation $g_x(0,0,z)=0$ first gives
$c_x(0,z)=0$. In the identity
$g=U P$ the mixed derivative on
the critical curve is then
$g_{xy}=U b_x$, since the terms
containing $c,c_x,b$ vanish there.
The zero $xy$ Hessian entry therefore
gives $b_x(0,z)=0$. Thus
$x^2\mid b$, $x^2\mid c$, and
$x\mid a$.
The Tschirnhaus change $y=Y-a/3$
therefore gives exactly
\[
 P=Y^3+x^2H(x,Y,z).
\]
The coefficient $H(0,0,z)$ is a unit
multiple of the original Hessian
coefficient $h(z)$. It is not
identically zero but vanishes at
zero. Perturb it to $H+\eta$.
At a simple root of
$H(0,0,z)+\eta$, the coordinates
$(x,Y,Z=H+\eta)$ give
$Y^3+x^2Z$. P4.18 excludes this
binomial germ by its compact
spherical class of nonzero square.

There are two nested transport
neighbourhoods in this argument.
First fix the original ball for
$g\to g+\varepsilon y^3$.
Then choose a smaller ball in the
Weierstrass coordinate neighbourhood,
where $g+\varepsilon y^3$ has only
the original critical curve and its
boundary has transverse type $A_2$.
The second perturbation transports
its spherical class into a nearby
fibre on this smaller ball. Inclusion
puts it in a fibre on the original
ball; the first boundary transport
puts it in a fibre of $g$.
This does not require extending the
second perturbation outside its
coordinate neighbourhood.
Units in the reduced equations are
retained during the perturbations.
At a binomial special point a unit
can be absorbed by the positive
weighted coordinate scaling of
P4.16--P4.18.

We have ruled out both a zero
Hessian and a vanishing cubic
coefficient. The product form
follows. Its normalization is
the finite map
\[
 (u,z)\longmapsto (a=u^3,b=-u^2,z),
\]
whose source is smooth, whose
inverse rational parameter off
$b=0$ is $u=-a/b$, and whose
restriction to the conductor is
the identity on the $z$ coordinate.
\end{proof}

\begin{statement}{P4.29}{Proved}
Let $S$ be an irreducible singular support
in Part~4.2 whose normalization has ruled
resolution, and suppose $e(S)=0$.
In the double-horizontal-divisor case,
when the generic transverse type is
cuspidal $A_2$, the normalization is
a smooth minimal elliptic ruled surface.
The image conductor is a smooth elliptic
curve and every point along it has the
product cusp germ of P4.28.
\end{statement}
\begin{proof}
Use the notation of P4R.7 and P4R.9.
The actual resolved pole support is
a section $H$ with rational trees,
and its inclusion in the ruled
resolution $T$ induces an isomorphism
on fundamental groups. Thus
$\pi_1(S)=\pi_1(C)$, where
$C=\operatorname{Sing}(S)_{\rm red}$.
The virtually abelian curve-group
calculation gives $e(C)\ge0$, and
P4.13 gives $e(S_{\rm reg})\ge0$.
Equality makes both numbers zero.
The same connected-curve calculation
shows that $C$ has exactly one
normalized component, of genus one,
and has no point-identification
cycles.

All rational components of the
resolved pole support are contracted
to points: a nonconstant proper map
from a rational curve to the
elliptic normalization of $C$ is
impossible. Hence the strict
conductor consists only of $H$.
The generic cusp has one
normalization branch, so
$H\to C$ has generic degree one
and is the normalization of $C$.
In particular $H$ and the ruling
base both have genus one.

Every exceptional curve of
$T\to N$ is now rational and
vertical for the ruling, since
$H$ is not contracted. The ruling
$r:T\to H$ descends to a holomorphic
map $\bar r:N\to H$. For an explicit
descent, the image of
$(\rho_N,r):T\to N\times H$
projects properly and bijectively
to $N$, because $r$ is constant
on each exceptional fibre. Its
normalized graph projects finitely
and birationally to the normal
space $N$, hence is isomorphic
to $N$.

The normalization singularities
are rational. Indeed a sequence
of point blowups of a
$\mathbb P^1$-bundle has
$R^1r_*\mathcal O_T=0$.
For a blowup this follows from
the two affine charts, or from
the exact sequence of the
exceptional $\mathbb P^1$ with
$H^1(\mathcal O_{\mathbb P^1})=0$;
the assertion for the whole
sequence follows by Leray.
Since $\rho_{N*}\mathcal O_T
=\mathcal O_N$, the low-degree
Leray sequence for
$r=\bar r\circ\rho_N$ gives
\[
 \bar r_*R^1\rho_{N*}\mathcal O_T=0.
\]
Here $R^2\bar r_*\mathcal O_N=0$
because $\bar r$ has one-dimensional
fibres. The sheaf
$R^1\rho_{N*}\mathcal O_T$ is
supported at finitely many
normalization singular points,
and direct image of a nonzero
such sheaf to $H$ is nonzero.
Thus it vanishes. In particular
$\chi(\mathcal O_N)
=\chi(\mathcal O_T)=0$.

We next check the scheme structure
of the conductor, to prove that
its image curve is smooth.
Write $I$ for the conductor ideal
on $S$ and on $N$, and let
$C_{\rm sch}\subset S$ and
$D_{\rm sch}\subset N$ be its
two zero schemes. Locally
$\mathcal O_S$ is Gorenstein
of dimension two and
$\nu_*\mathcal O_N$ is a
maximal Cohen--Macaulay module.
Finite duality identifies
its dual module with $I$,
which is again maximal
Cohen--Macaulay. The depth
lemma in
$0\to I\to\mathcal O_S
\to\mathcal O_{C_{\rm sch}}\to0$
therefore gives depth one
for the conductor curve.
At a generic cusp its local
length in the transverse
direction is one. A
one-dimensional Cohen--Macaulay
scheme which is generically
reduced is reduced: any
nilpotent ideal would have
zero-dimensional support,
contradicting the absence of
embedded associated points.
Consequently $C_{\rm sch}=C$.

Put $L=\nu^*\omega_S$, a
torsion line bundle by P4.12.
Finite duality on $N$ gives
\[
 I=\omega_N\otimes L^{-1}.
\]
Serre duality on the compact
Cohen--Macaulay surface $N$
gives $\chi(I)=\chi(L)$.
The rational-singularity
comparison just proved and
the projection formula give
$\chi(L)=\chi(\rho_N^*L)$.
Riemann--Roch on the smooth
surface $T$, with the rational
Chern class of this torsion
line bundle equal to zero,
gives $\chi(\rho_N^*L)
=\chi(\mathcal O_T)=0$.
Therefore
$\chi(\mathcal O_{D_{\rm sch}})
=\chi(\mathcal O_N)-\chi(I)=0$.

The exact conductor square is
\[
 0\longrightarrow\mathcal O_S
 \longrightarrow
 \nu_*\mathcal O_N\oplus\mathcal O_C
 \longrightarrow\nu_*\mathcal O_{D_{\rm sch}}
 \longrightarrow0.
\]
Taking Euler characteristics and
using $\chi(\mathcal O_S)=0$
from P4.12 yields
$\chi(\mathcal O_C)=0$.
Its normalization is the
smooth elliptic curve $H$,
so the normalization sequence
gives $\chi(\mathcal O_C)
 =-\sum_p\delta(C,p)$.
All delta invariants vanish;
thus $C$ is smooth.

Every point of $C$ now satisfies
P4.28, so $S$ has the product
cusp germ there and $N$ is
smooth along the conductor.
It was already smooth away
from the conductor. P4.22
therefore applies and makes
this elliptic ruled
normalization minimal.
\end{proof}


\subsection{The conormal calculation and ordinary-triple equality}

\begin{statement}{P4.30}{Proved}
Let $C$ be a compact irreducible reduced curve in the smooth
threefold $Y$ of Part~4.2. Suppose its normalization is an elliptic
curve $E$, and every point of $C$ is unibranch. Write $P$ for its
ideal in $Y$ and $P^{(2)}$ for the ideal of functions vanishing to
order at least two at the generic point of $C$. Then
\[
 \chi(\mathcal O_Y/P^{(2)})\le0.
\]
Equality forces $C$ to be smooth. The assertion uses
$c_1(TY)\cdot C=0$, which follows from $b_2(Y)=0$.
\end{statement}
\begin{proof}
Let $j:E\to Y$ be the normalization map followed by inclusion.
The derivative map
\[
 j^*\Omega_Y\longrightarrow\Omega_E
\]
has image $\Omega_E(-R)$ for an effective divisor $R$.
Its kernel $K$ is locally free of rank two: dividing the derivative
row by its common zero gives a surjection onto a line bundle.
Since $\deg j^*\Omega_Y=0$ and $\deg\Omega_E=0$, we have
$\deg K=\deg R$.

Differentiation induces an injection
\[
 Q:=P/P^{(2)}\longrightarrow j_*K.
\]
Indeed its kernel consists precisely of functions with zero first
normal derivative at a generic smooth point of $C$, which is the
defining order-two condition for $P^{(2)}$. Both sides agree at the
generic point. Their quotient has finite length, denoted by $\tau$.
Riemann--Roch on the smooth elliptic curve gives
\[
 \chi(Q)=\deg R-\tau.
\]
We prove $\tau\ge\deg R$ by local equations, including at a
singular point of $C$.

Parametrize its sole branch by $t$. Choose ambient coordinates
vanishing at the point, and let
$m=\min_i\operatorname{ord}_t x_i(t)$ be the multiplicity.
The common order of the three derivative entries is $m-1$,
so the coefficient of $R$ is $m-1$. The saturated kernel $K$
is a direct summand of the free rank-three module over
$\mathbb C\{t\}$. Consequently a generator of $\bigwedge^2K$
has at least one unit Pl\"ucker coordinate. If $L$ is the
$\mathbb C\{t\}$-span of the differentiated defining equations,
then the colength of $L$ in $K$ is the minimum order of their
two-by-two minors. This is the diagonal elementary-divisor
calculation for a rank-two submodule over a discrete valuation
ring. The local length of $j_*K/Q$ is at least this colength,
since $Q\subset L\subset K$.

At a smooth point $m=1$ and there is nothing to prove.
At a singular point of embedding dimension three, every defining
equation belongs to the square of the ambient maximal ideal.
Every entry of its gradient therefore has order at least $m$
on the parametrization, and every two-by-two minor has order
at least $2m$. Hence the local quotient length is at least
$2m\ge m-1$.

At a singular point of embedding dimension two, an invertible
ambient coordinate change writes the ideal as $(z,h(x,y))$,
where $h$ is a reduced irreducible plane equation without a
linear term. The conormal image has a free $dz$ summand.
Its other generator is $dh$; both its plane coefficients have
order at least $m$. The primitive rank-one kernel of the plane
derivative row has a unit coefficient, so $dh$ is its generator
multiplied by a function of order at least $m$. The quotient
length is therefore at least $m\ge m-1$. Embedding dimension
one is the smooth case. These cases prove the local bound and
therefore $\chi(Q)\le0$.

The normalization exact sequence gives
$\chi(\mathcal O_C)=-\sum_p\delta(C,p)$.
Finally the exact sequence
\[
 0\longrightarrow P/P^{(2)}\longrightarrow
 \mathcal O_Y/P^{(2)}\longrightarrow\mathcal O_C
 \longrightarrow0
\]
gives
\[
 \chi(\mathcal O_Y/P^{(2)})
 =-\sum_p\delta(C,p)+\deg R-\tau\le0.
\]
Equality forces all delta invariants to vanish, so the
normalization is an isomorphism and $C$ is smooth.
\end{proof}

\begin{statement}{P4.31}{Proved}
Let $S$ be an irreducible singular support in Part~4.2 whose
normalization has ruled resolution, and suppose $e(S)=0$.
In the double-horizontal-divisor case, when the generic
transverse type is the ordinary triple point $xy(x-y)=0$,
the normalization is a smooth minimal elliptic ruled surface.
Its conductor image is a smooth elliptic curve, and its
reduced conductor is an unramified cover of that curve of
degree three. Every germ along the image is analytically
the product of an ordinary plane triple point with a disc.
\end{statement}
\begin{proof}
Use the minimal resolution $T\to N\to S$ and the section $H$
of P4R.7 and P4R.9. Put $g=g(H)$ and
$C=\operatorname{Sing}(S)_{\rm red}$.
The same fundamental-group and Euler argument as in P4.29
shows that $C$ is connected, has just one normalized
component, that this normalization $E$ has genus one, and
that its point-identification graph has no cycles. Thus
every point of $C$ is unibranch. All rational components
of the resolved pole support are contracted, and the
only strict conductor component is $H$. Its finite map
to $C$ has degree three, since an ordinary transverse
triple point has three normalization branches.
In particular $g\ge1$ by Riemann--Hurwitz for $H\to E$.

All exceptional curves of $T\to N$ are rational and
vertical for the ruling. Exactly the normalized-graph
descent and Leray calculation of P4.29 give a morphism
$N\to H$ and
\[
 R^1\rho_{N*}\mathcal O_T=0,\qquad
 \chi(\mathcal O_N)=\chi(\mathcal O_T)=1-g.
\]
Thus the normalization singularities are rational; no
smoothness assertion is used here.

Let $I$ be the conductor ideal, and let $C_{\rm sch}$
and $D_{\rm sch}$ be its zero schemes in $S$ and $N$.
Finite duality and the depth argument of P4.29 show that
$\mathcal O_{C_{\rm sch}}$ is Cohen--Macaulay of pure
dimension one. If $P\subset\mathcal O_Y$ is the reduced
ideal of $C$, the inverse image $J$ of $I$ in
$\mathcal O_Y$ therefore has no embedded associated
prime. Its only associated prime locally is $P$.
The generic ordinary triple point has conductor equal
to the square of its maximal ideal in the transverse
plane. To check this directly, normalize its three
lines. Constant functions on them have a common value;
their linear terms are the evaluations of one linear
form on three distinct directions, a two-dimensional
space. Quadratic forms evaluate surjectively on those
three directions. All terms of degree at least two
can therefore be chosen independently on the three
branches, and a nonzero constant or linear term cannot
belong to the conductor. This proves the stated square
ideal, including its scheme structure.

Consequently $J_P=P_P^2$, and in fact
\[
 J=P^{(2)}.
\]
For completeness, if $a\in P^{(2)}$, some $s\notin P$
satisfies $sa\in J$. Multiplication by $s$ is injective
on $\mathcal O_Y/J$, because its only associated prime
is $P$, so $a\in J$. The reverse inclusion follows by
localizing at $P$. Thus the conductor image is the
actual symbolic double curve to which P4.30 applies.

Put $L=\nu^*\omega_S$, which is torsion. Finite duality
gives $I=\omega_N\otimes L^{-1}$. Serre duality, the
rational-singularity comparison and Riemann--Roch on
$T$ give
\[
 \chi(I)=\chi(L)=\chi(\rho_N^*L)
   =\chi(\mathcal O_T)=1-g.
\]
Hence $\chi(\mathcal O_{D_{\rm sch}})=0$.
Taking Euler characteristics in the exact conductor
square, and using $\chi(\mathcal O_S)=0$, yields
\[
 \chi(\mathcal O_{C_{\rm sch}})=g-1.
\]
P4.30 now gives $g-1\le0$. Since $g\ge1$, equality
holds, $g=1$, and $C$ is smooth. Riemann--Hurwitz then
makes the degree-three map $H\to C$ unramified.

We finish by proving normalization smoothness from
the actual local morphisms. Over a point of $C$ the
three distinct points of $H$ remain distinct in $N$:
every exceptional block is vertical for the ruling
$T\to H$ and therefore cannot join different points
of $H$. The finite normalization consequently gives
three local irreducible branches of $S$. Each branch
is a hypersurface in $Y$ and is smooth at the generic
point of $C$. Its singularities are isolated, so the
hypersurface Cohen--Macaulay property and regularity
in codimension one make it normal. The corresponding
component of $N$ is that branch itself.

Write its three local equations as $f_1,f_2,f_3$.
The intersection of any two branches is supported on
$C$ and has generic multiplicity one, because the
generic transverse triple point is ordinary. On a
normal branch the other equation is a nonzero divisor.
Its quotient is a one-dimensional Cohen--Macaulay
ring, generically reduced with support $C$, so it is
reduced and is exactly $\mathcal O_C$. Thus $C$ is a
smooth Cartier curve in each branch. A local ring of
dimension two with a nonzero divisor whose quotient
is a regular one-dimensional local ring has embedding
dimension at most two: lift its one maximal-ideal
generator and add the nonzero divisor. It is regular.
All three branches are therefore smooth.

Any two branches meet transversely along $C$.
Choose their equations as coordinates $x,y$, with
$z$ a coordinate on $C$. The third equation is
$a(x,y,z)x+b(x,y,z)y$, where both $a$ and $b$ are
units by the pairwise transversality just proved.
The invertible coordinate change $X=ax$, $Y=-by$
writes the reduced union exactly as
$XY(X-Y)=0$. This is the asserted product triple
point, and its normalization is smooth. Away from
the conductor, normalization singularities were
already excluded by P4R.6. Finally P4.22 applies to
this smooth elliptic ruled normalization and proves
its minimality.
\end{proof}


\subsection{Complete zero-Euler classification in the requested classes}

\begin{statement}{P4.32}{Proved}
An irreducible support in Part~4.2 has Euler number zero exactly
in the following cases: it is smooth with a finite unramified
torus cover, or it is singular with smooth minimal elliptic ruled
normalization, or it is singular with smooth minimal Hopf
normalization. In the ruled double-horizontal case the conductor
image is a smooth elliptic curve; the local germs are product
cusps, product tacnodes, or product ordinary plane triple points,
and the reduced conductor maps have degrees one, two, or three
respectively and are unramified.

For a fibration whose supports are all irreducible or normal
crossing, $c_3(Y)=0$ exactly when every irreducible support has
one of these forms and every reducible support has the
zero-Euler normal-crossing form of P4.8.
\end{statement}
\begin{proof}
Normal irreducible supports are smooth with the stated finite
unramified cover by P4.8 and P4.12. For a singular irreducible
support, use its minimal normalization resolution $T$.
P4.26 excludes rational resolutions and class-VII minimal
models with positive second Betti number from equality.
The Inoue minimal model has no conductor curves and is
excluded there as well. A Hopf minimal model is covered
completely by P4.27: equality makes the actual normalization
a smooth minimal Hopf surface, and the converse is P4.24.

It remains to use the positive-genus ruled resolution.
The reduced-horizontal alternatives of P4R.9 already give
smooth minimal elliptic ruled normalization. In the remaining
case the pole support is a section $H$ with rational trees,
and the coefficient of $H$ is two. Its inclusion in $T$
induces an isomorphism on fundamental groups. The conductor
pushout and the curve-group calculation of P4R.7--P4R.9 give
$\pi_1(S)=\pi_1(C)$ and $e(C)\ge0$ for the reduced image
conductor $C$. Equality in
$e(S)=e(S_{\rm reg})+e(C)$ makes $C$ a connected curve
with elliptic normalization and no point-identification
cycles. Every rational pole component is contracted, since
there is no finite nonconstant map from a rational curve
to an elliptic curve. Thus $H$ is the sole strict conductor
component, and every branch of a generic transverse
plane-curve singularity has conductor exponent two.

We enumerate those transverse germs explicitly. For a branch
$B_i$ of a reduced plane curve, its conductor exponent is
\[
 c_i=2\delta(B_i)+\sum_{j\ne i}I(B_i,B_j).
\]
If a branch is singular, $c_i=2$ forces
$\delta(B_i)=1$ and no other branch. The one-gap
parametrization of a delta-one irreducible plane curve
has orders two and three and gives the ordinary cusp.
If all branches are smooth, two branches have intersection
multiplicity two and give a tacnode, or three branches
have all pairwise intersection multiplicities one and
give an ordinary triple point. Four branches would give
each conductor exponent at least three. One smooth branch
is nonsingular and has exponent zero. These exhaust the
possibilities.

P4.29, P4R.11, and P4.31 respectively treat the cusp,
tacnode, and ordinary-triple cases, without assuming that
the initial normalization or image conductor was smooth.
They give the asserted smooth minimal elliptic ruled
normalization and the precise local conductor maps.
The converse for a smooth minimal elliptic ruled
normalization is P4.22. This proves the entire irreducible
zero-Euler classification.

Finally P4.26 expresses $c_3(Y)$ as a finite sum of
nonnegative fibre Euler numbers. Its vanishing is therefore
equivalent to vanishing of every summand. The irreducible
summands have just been classified, and P4.8 gives the
complete reducible normal-crossing zero case. This proves
the final assertion.
\end{proof}


\subsection{Normalization cohomology for the remaining positive cases}

\begin{statement}{P4.33}{Proved}
Let $S$ be an irreducible singular support in Part~4.2, and let
$\rho:T\to N\to S$ be its minimal normalization resolution.
Then $H^2(N,\mathcal O_N)=0$. If
$p=\dim H^0(N,R^1\rho_{N*}\mathcal O_T)$ is the sum of the
geometric genera of the normalization singularities, then
\[
 p=h^1(T,\mathcal O_T)-h^1(N,\mathcal O_N).
\]
In particular a rational resolution forces every normalization
singularity to be rational and gives $\chi(\mathcal O_N)=1$.
For the scheme-theoretic image conductor $C_{\rm sch}$ we have
\[
 \chi(\mathcal O_{C_{\rm sch}})
   =-\chi(\mathcal O_N)=h^1(N,\mathcal O_N)-1.
\]

In the ruled double-horizontal case, with section $H$ of genus
$g\ge1$, a contracted $H$ gives $h^1(N,\mathcal O_N)=0$ and
$p=g$. A noncontracted $H$ gives $p=0$ and
$h^1(N,\mathcal O_N)=g$.
\end{statement}
\begin{proof}
Write $I_S\subset\mathcal O_S$ for the conductor ideal.
Finite duality gives
\[
 \nu_*\omega_N=I_S\otimes\omega_S.
\]
The dualizing line bundle $\omega_S$ is torsion by P4.12.
A nonzero section of a torsion line bundle on a compact
irreducible reduced complex space has no zero: raising it
to a trivializing positive power gives a nonzero holomorphic
function, hence a nonzero constant. Such a section cannot
belong to the proper conductor ideal along the nonempty
conductor. Therefore $H^0(N,\omega_N)=0$. Serre duality
on the compact Cohen--Macaulay normal surface gives
$H^2(N,\mathcal O_N)=0$.

The resolution is birationally ruled or of class VII and
has $H^2(T,\mathcal O_T)=0$, by P4.12 and the classification
already used in P4.26. The low-degree Leray sequence,
using $\rho_{N*}\mathcal O_T=\mathcal O_N$ and the
finite support of $R^1\rho_{N*}\mathcal O_T$, is therefore
\[
 0\longrightarrow H^1(N,\mathcal O_N)
 \longrightarrow H^1(T,\mathcal O_T)
 \longrightarrow H^0(N,R^1\rho_{N*}\mathcal O_T)
 \longrightarrow0.
\]
This proves the formula for $p$. If $T$ is rational,
$h^1(\mathcal O_T)=0$, so $p=0$. Each nonnegative local
geometric genus is zero, giving the asserted rational
singularities and Euler characteristic.

For completeness the conductor Euler identity does not
require rational singularities. Put $L=\nu^*\omega_S$,
and write $I_N=\omega_N\otimes L^{-1}$ for the conductor
on $N$. Serre duality gives $\chi(I_N)=\chi(L)$.
Riemann--Roch on $T$, applied to the torsion line bundle
$\rho_N^*L$, gives
$\chi(\rho_N^*L)=\chi(\mathcal O_T)$.
The projection formula makes
$R^1\rho_{N*}\rho_N^*L=L\otimes R^1\rho_{N*}\mathcal O_T$;
its finite length is still $p$. Leray therefore gives
$\chi(L)=\chi(\mathcal O_N)$.
Consequently the conductor scheme $D_{\rm sch}\subset N$
has $\chi(\mathcal O_{D_{\rm sch}})=0$.
The exact conductor square, together with
$\chi(\mathcal O_S)=0$, proves
$\chi(\mathcal O_{C_{\rm sch}})=-\chi(\mathcal O_N)$.
Since $N$ is compact connected and $H^2(\mathcal O_N)=0$,
the remaining equality follows.

Finally suppose the ruled pole support is $H$ with
rational vertical trees. If $H$ is contracted, the
composition $H\to T\to N$ is constant. Pullback from
$H^1(N,\mathcal O_N)$ to $H^1(H,\mathcal O_H)$ is zero.
But restriction
$H^1(T,\mathcal O_T)\to H^1(H,\mathcal O_H)$ is an
isomorphism: the ruling has section $H$, a
$\mathbb P^1$-bundle has zero first higher direct image
of its structure sheaf, and the actual point blowups
preserve these groups. The Leray injection above
therefore forces $H^1(N,\mathcal O_N)=0$, so $p=g$.
If $H$ is not contracted, every exceptional curve is
vertical. The ruling descends by the normalized-graph
comparison of P4.29. Its Leray calculation gives $p=0$,
and the same exact sequence gives $h^1(\mathcal O_N)=g$.
\end{proof}


\subsection{The value one in the rational ordinary-triple sector}

\begin{statement}{P4.34}{Proved}
An irreducible singular support with rationally resolved
normalization and generic transverse ordinary triple point
cannot have Euler number one. For a fibration all of whose
supports are irreducible or normal crossing, its Euler number
is therefore two.
\end{statement}
\begin{proof}
Suppose its Euler number were one. The rational
normal-crossing case has Euler number two by P4.8, so we
are in the simply connected resolved-conductor alternative
of P4.26. That proof gives $\pi_1(S)=\pi_1(C)=F_b$,
$b\le1$, and
$e(C)=c+1-b$ for the number $c$ of normalized rational
components. Since $e(S_{\rm reg})\ge0$, the value one
forces $c=b=1$ and $e(S_{\rm reg})=0$.
Thus $C$ has normalization $E=\mathbb P^1$, exactly one
point with two analytic branches, and only unibranch
singular points elsewhere.

P4.33 proves that $N$ has rational singularities and
that its image conductor scheme has
$\chi(\mathcal O_{C_{\rm sch}})=-1$.
The generic ordinary-triple conductor and pure
Cohen--Macaulay argument of P4.31 identify this scheme
with the symbolic double curve $\mathcal O_Y/P^{(2)}$.
At a point with two analytic branches, the symbolic ideal
means order at least two along each branch. The same
argument applies: a nonzero finite-support defect in the
conductor quotient would violate depth one.

We compute its Euler characteristic and the possible
local defects. On $E$ use the bundles $K$ and the
ramification divisor $R$ of P4.30. This time
$\deg K=2+\deg R$, since $E=\mathbb P^1$.
Write $\tau_p$ for the local length of the conormal
quotient $j_*K/(P/P^{(2)})$ and $r_p$ for the sum of
the coefficients of $R$ over $p$. Normalization and
Riemann--Roch give
\[
 \chi(\mathcal O_Y/P^{(2)})
 =5-\sum_p d_p,
 \qquad d_p=\delta(C,p)+\tau_p-r_p.
\]
We prove the precise gap needed here: at an ordinary
node $d_p=5$; at every planar singularity
$d_p=5\delta(C,p)$; and at every nonplanar singularity
having at most two branches $d_p\ge7$.

First consider a planar curve, with ideal $(z,h(x,y))$.
Its conormal is free over $\mathcal O_C$, generated by
$dz,dh$. On each normalized branch, write
$dh=a_i u_i$ for a primitive plane kernel generator,
and put $v_i=\operatorname{ord}a_i$.
The inclusion of each of these two free summands into
its normalized module has length $\delta(C,p)$.
The second summand has the additional vanishing orders
$v_i$. Hence
\[
 \tau_p=2\delta(C,p)+\sum_i v_i.
\]
We verify $\sum_i(v_i-r_i)=2\delta(C,p)$ directly by
curve duality. Let $A$ be the reduced plane-curve ring,
$B$ its normalization, and $I=\operatorname{Hom}_A(B,A)$
its conductor. Applying $\operatorname{Hom}_A(-,A)$ to
$0\to A\to B\to B/A\to0$ gives
$A/I\simeq\operatorname{Ext}^1_A(B/A,A)$.
Here $\operatorname{Ext}^1_A(B,A)=0$ because $B$ is
maximal Cohen--Macaulay over the one-dimensional
Gorenstein ring $A$. The dual of a finite-length module
has the same length, so $\operatorname{length}(A/I)=\delta$,
and $\operatorname{length}(B/I)=2\delta$.
If the conductor exponent on branch $i$ is $c_i$,
this says $\sum_i c_i=2\delta$.
Choose a generic plane coordinate $x$ with minimal
valuation on every normalized branch. The dualizing
generator $dx/h_y$ then pulls back with
order $r_i-v_i$. The identity
$\omega_B=I\omega_A$ therefore gives
$c_i=v_i-r_i$. Combining these equalities proves
$\tau_p-r_p=4\delta$ and the stated planar formula.

Now consider a nonplanar unibranch curve of multiplicity
$m$. Its embedding dimension is three. As in P4.30,
all gradient entries have valuation at least $m$, so
$\tau_p\ge2m$, while $r_p=m-1$ and
$\delta\ge m-1$. A nonplanar unibranch curve has
$m\ge3$. To see the last assertion, if $x=t^2$, write
each other coordinate as $a(x)+t b(x)$. After subtracting
the even parts, choose the nonzero $b$ of least order.
The other odd coordinate is an analytic multiple of it
over $\mathbb C\{x\}$. This is a smooth ambient equation
containing the curve, so its embedding dimension is at
most two. The case where both odd parts vanish contradicts
the birational normalization parametrization.

The estimates give $d_p\ge2m\ge6$. Equality at six
would require $m=3$, $\delta=2$, and $\tau_p=6$.
The valuation semigroup of a unibranch curve has exactly
$\delta$ gaps, as follows by filtering its local ring
and its normalization by order. Thus the two gaps here
are one and two, and we can choose generators with
initial orders three, four, and five. They form an
ambient coordinate system: their classes generate the
three-dimensional cotangent space, since the square
of the maximal ideal has order at least six.
In a defining relation, the quadratic monomials
$x^2$ and $xy$ have respectively the unique least
orders six and seven, so their coefficients vanish.
Every gradient with valuation three consequently has
its initial vector in the single $z$ direction
(from a possible $xz$ term). Two such initial
gradient vectors are dependent. Every two-by-two
minor therefore has valuation at least seven,
contradicting $\tau_p=6$. Hence $d_p\ge7$ in this case.

Finally consider a nonplanar point with two branches
of multiplicities $m_1,m_2$. Two smooth curve branches
in a smooth threefold are contained in a smooth
surface. Indeed choose a common immersive coordinate
$x$ and write them as graphs
$(y,z)=(a_i(x),b_i(x))$. Divide their difference pair
by its common power of $x$; the equation
\[
 \frac{b_1-b_2}{x^k}(y-a_1)
 -\frac{a_1-a_2}{x^k}(z-b_1)=0
\]
has a nonzero linear derivative and contains both.
Thus in the nonplanar case $m_1+m_2\ge3$.
The absence of a linear defining equation gives,
on the two branches separately,
$\tau_p\ge2(m_1+m_2)$.
The reduced-curve normalization sequence gives
$\delta(C,p)=\delta_1+\delta_2+
\operatorname{length}\mathcal O_Y/(P_1+P_2)
\ge m_1+m_2-1$.
Also $r_p=m_1+m_2-2$.
It follows that
$d_p\ge2(m_1+m_2)+1\ge7$.
This proves the gap in all cases.

At the unique two-branch point, a planar singularity
has defect five if its delta invariant is one, and
at least ten otherwise. A nonplanar point has defect
at least seven. Every other singular point is
unibranch and contributes at least five if planar,
and at least seven if nonplanar. Smooth points
contribute zero. Thus the sum of the defects can
be five or at least seven, but never six.
The required value $\chi(\mathcal O_{C_{\rm sch}})=-1$
would give exactly six. This contradiction excludes
$e(S)=1$.

The final statement uses P4.26: the sum of all the
nonnegative fibre Euler numbers is at most two.
The rational-resolution case is strictly positive,
and we have just excluded one, leaving exactly two.
\end{proof}


\subsection{The conductor graph excludes Euler number one in the rational case}

\begin{statement}{P4.35}{Proved}
No irreducible support with rationally resolved normalization
has Euler number one. Consequently, in a fibration whose
supports are all irreducible or normal crossing, every
irreducible singular support with rationally resolved
normalization has Euler number exactly two.

More generally, an irreducible support cannot have reduced
image conductor equal to a rational curve with exactly one
point-identification cycle when its reduced normalization
conductor has no graph cycles. This excludes that conductor
case for every ruled-resolution double-horizontal pole tree.
\end{statement}
\begin{proof}
We first prove the graph assertion independently of
the transverse singularity type and of ramification.
Let $D\subset N$ be the reduced normalization conductor,
and suppose its point-identification graph is a tree.
Let $c$ be its number of normalized compact components.
Suppose its image $C$ is irreducible with rational
normalization and exactly one point-identification
cycle. There is then a unique point $p\in C$ with
two analytic branches; all other points are unibranch.
Denote the two preimages of $p$ on $C^\nu$ by
$p_+,p_-$. Every normalized component of $D$ maps
finitely and surjectively to $C^\nu$.

The finite set $\nu^{-1}(p)$ has $n$ points and
parametrizes the local irreducible hypersurface
branches of $S$ at $p$. Every pair of such branches
intersects along a curve: two nonunit local
hypersurface equations in a regular local ring of
dimension three generate an ideal of height at most
two. Their common curve lies in $C$. Consequently
the nonempty subsets of $\{+,-\}$ recording which
branches of $C$ are contained in each hypersurface
branch pairwise intersect. Such a family of subsets
of a two-element set has a common element. If a
singleton occurs its element belongs to every set;
if no singleton occurs every set is the whole set.
Interchanging the labels, every point of
$\nu^{-1}(p)$ therefore has at least one conductor
branch mapping to the plus branch of $C$.

Let $r_+$ and $r_-$ count the points on $D^\nu$
over $p_+$ and $p_-$ respectively. The preceding
paragraph gives $r_+\ge n$. Surjectivity of the
map from each of the $c$ normalized components
gives $r_-\ge c$. The contribution of the points
over $p$ to the point-identification graph of $D$
is exactly
\[
 \sum_{q\in\nu^{-1}(p)}(r_q-1)
       =r_++r_--n\ge c.
\]
For a connected reduced compact curve with $c$
normalized components, the cycle rank of this graph
is $1-c+\sum_q(r_q-1)$. All other contributions
are nonnegative. Its rank is therefore at least
one, contradicting the tree hypothesis.
This proves the graph assertion.

Now let the normalization have rational resolution
and suppose $e(S)=1$. The normal-crossing alternative
has Euler number two by P4.8. In the other alternative
P4.26 gives a connected simply connected resolved
pole support $A\subset T$, and
\[
 \pi_1(S)=\pi_1(C)=F_b,\qquad b\le1,
 \qquad e(C)=c_C+1-b.
\]
Its normalized components are rational.
Since $e(S_{\rm reg})\ge0$, Euler number one forces
$c_C=b=1$ and $e(S_{\rm reg})=0$.
It remains to verify the tree property for the
actual conductor $D\subset N$, allowing singular $N$.

P4.33 makes every singularity of $N$ rational.
Every connected exceptional fibre of $T\to N$
is simply connected. Here is the precise curve
calculation. Over a small Stein neighbourhood,
$R^1\rho_{N*}\mathcal O_T=0$ gives
$H^1(\mathcal O_T)=0$ on its inverse image.
This proper map has every fibre of dimension at
most one over that Stein neighbourhood. Its coherent
cohomological dimension is consequently at most
one, so the second higher direct image of every
coherent sheaf vanishes. Applied to the ideal of the reduced
exceptional curve $E$, Leray and the Stein property
give $H^2(\mathcal O_T(-E))=0$ there. Its ideal
sequence consequently gives $H^1(E,\mathcal O_E)=0$.
The curve normalization exact sequence now implies
that all normalized components are rational and
that their point-identification graph is a tree.
Thus $E$ is simply connected.

By P4R.6, $A$ is the full inverse image of $D$.
The proper map $A\to D$ is the quotient which
collapses these connected exceptional curves to
points and is an isomorphism elsewhere.
Compatible triangulation makes these curve
inclusions cofibrations. Collapsing a connected
simply connected subcomplex preserves fundamental
groups by van Kampen. Collapsing the finitely many
exceptional fibres therefore gives $\pi_1(D)=0$.
Each normalized component of $D$ is rational,
being birational to a strict pole component in
$T$. Hence its point-identification graph is a
tree. The already proved graph assertion contradicts
the description of $C$. This excludes $e(S)=1$.

For the final more general application, the
resolved pole support in the ruled double-section
case is $H$ with rational trees attached.
Every exceptional fibre is connected and is a
union of its components. If it does not contain
$H$, it is a connected rational subtree. If it
contains $H$, all the remaining components form
rational trees attached to that component once.
Contracting such a connected subcurve cannot
create a cycle in the graph: contraction of a
connected subgraph of a tree again has a tree
as quotient. The only positive-genus vertex,
if contracted, is replaced by a point, and this
also introduces no graph edge. Thus the image
conductor $D$ has no graph cycles, and the same
contradiction applies.

Finally, for a fibration in the two requested
support classes, P4.26 makes the total sum of
nonnegative fibre Euler numbers at most two.
The rational-resolution case is strictly positive;
excluding one leaves exactly two.
\end{proof}


\subsection{Eliminating all positive-Euler Hopf resolutions}

\begin{statement}{P4.36}{Proved}
Every irreducible singular support in Part~4.2 whose
normalization has a Hopf minimal model has smooth minimal
Hopf normalization and Euler number zero. The assertion
does not assume initially that the normalization is smooth
or that its Euler number vanishes.
\end{statement}
\begin{proof}
Let $T\to N\to S\subset Y$ be the minimal normalization
resolution, and let $\pi:T\to T_0$ contract its $k$ blowups
to a minimal Hopf surface. We first take a finite unramified
primary Hopf cover $T'_0\to T_0$ and pull back the actual
sequence of point blowups. This gives a finite unramified
map $T'\to T$ and a blowup morphism $T'\to T'_0$.
Compose with $T\to Y$. All Chern calculations below take
place on this explicit smooth cover; no cover of the
singular normalization $N$ is asserted. To avoid a second
set of symbols we write $T,T_0,k,\rho$ for these covered
objects until the last paragraph.

The primary Hopf normal form and proof of P4.24 give a
nowhere-zero holomorphic foliation on $T_0$ tangent to
every compact curve. Its saturated pullback to $T$ is a
line subbundle away from finitely many simple saddle
zeros. It is tangent to every strict conductor component
and every created exceptional curve. The actual conductor
centres lie on these invariant curves. Every blowup centre
is either a regular point of this foliation or one of its
saddles. Let $\varepsilon=0$ in the first case and
$\varepsilon=1$ in the second. The same explicit two-chart
calculation as in P4.24 gives
\[
 c_1(L)=-K_T+V,\qquad
 V'=\mu^*V+\varepsilon E,
 \qquad K_T^2=-k,\qquad K_TV-V^2=0.
\]
Topologically trivial line-bundle factors from $T_0$ do
not change these integral intersection calculations.

Put $B=\rho^*TY$. Its first and second rational Chern
classes vanish: $b_2(Y)=0$ gives $c_1(TY)=0$ rationally,
and Poincar\'e duality gives $H^4(Y;\mathbb Q)=0$ and
hence $c_2(TY)=0$. The derivative $L\to B$ can now have
a divisorial zero, since $\rho$ contracts exceptional
curves. Denote that effective integral divisor by $Z$.
Dividing by it gives the saturated line injection in
codimension one
\[
 L'=L(Z)\longrightarrow B,
 \qquad c_1(L')=-K_T+V+Z.
\]
The quotient $Q=B/L'$ is torsionfree of rank two.
Write $z_0\ge0$ for the length of $Q^{**}/Q$.
The double dual is locally free on the smooth surface,
and the Chern sequence gives
\[
 c_2(Q^{**})=(-K_T+V+Z)^2-z_0.
\]

Let $E_B\subset B$ be the saturated rank-two image
of $d\rho:TT\to B$. Its determinant is
$\det E_B=\omega_T^{-1}(R)$, where the effective
integral divisor $R$ is the divisorial zero of the
determinant map. Since $E_B$ is saturated in $B$,
the divided foliation line $L'$ is contained in it.
The image of $E_B/L'$ in $Q^{**}$ is saturated in
codimension one. Its rank-one double dual
$J=(E_B/L')^{**}$ is a line bundle. The map into
$Q^{**}$ extends across the finitely many omitted
points by Hartogs, since both sheaves are locally
free there. Its saturation is therefore a line bundle
of first Chern class $R-V-Z$; the rank-one quotient
is an ideal sheaf of a zero-dimensional scheme
times a line bundle of first Chern class $K_T-R$.
Let its length be $z_1\ge0$. Thus
\[
 (-K_T+V+Z)^2-z_0
       =(R-V-Z)(K_T-R)+z_1.
\]
Using $K_T^2=-k$ and $K_TV-V^2=0$, this is precisely
\[
 \mathcal B(R,Z)
 :=R^2-K_TR-VR-ZR-K_TZ+2VZ+Z^2
 =k+z_0+z_1.                                      \tag{P4.36.1}
\]

We now evaluate the left side through the actual
point blowups, retaining all integral coefficients.
Push $R$ and $Z$ down along the chosen sequence.
At one reversed blowup write
\[
 R'=\mu^*R+cE,\quad Z'=\mu^*Z+dE,
 \quad K'=\mu^*K+E,\quad V'=\mu^*V+\varepsilon E.
\]
The integers $c,d$ are not assumed positive.
Since $E^2=-1$ and a pulled-back divisor has zero
intersection with $E$, direct expansion gives
\[
 \mathcal B(R',Z')-\mathcal B(R,Z)
 =-c^2+(1+\varepsilon)c+cd+(1-2\varepsilon)d-d^2.
\]
Set $a=c-1$ and $b=d-(1-\varepsilon)$. For each
$\varepsilon\in\{0,1\}$ this is exactly
\[
 1-(a^2+b^2-ab)\le1.                            \tag{P4.36.2}
\]
The quadratic form is
$\frac12(a^2+b^2+(a-b)^2)$; it is nonnegative
and is zero exactly when $a=b=0$.
All divisor intersections on the initial Hopf
surface are zero because $b_2(T_0)=0$.
Summing (P4.36.2) therefore gives
$\mathcal B(R,Z)\le k$.
Equation (P4.36.1) forces equality, $z_0=z_1=0$,
and at every single blowup
\[
 c=1,\qquad d=1-\varepsilon.                    \tag{P4.36.3}
\]

Every created exceptional curve has positive
coefficient in $Z$. At a regular centre, (P4.36.3)
adds one to its nonnegative pulled-back coefficient.
At a saddle centre it adds zero, but at least one
previously created exceptional curve passes through
that saddle. The original Hopf foliation has no
zero; the two-chart blowup calculation places each
new saddle on the new exceptional curve. Induction
therefore gives a positive pulled-back coefficient
at every later saddle blowup as well.

The foliation is tangent to each exceptional curve.
If its image under $\rho$ were a curve, the derivative
along its tangent line would be nonzero at a generic
point: a nonconstant holomorphic map of smooth
normalized curves in characteristic zero is
generically immersive. Positive divisorial order
of $d\rho|_L$ contradicts that fact. Thus every
created exceptional curve is contracted by $\rho$,
and hence by the normalization resolution, since
the normalization map is finite.

Return to the original, possibly secondary Hopf
surface. Every exceptional curve downstairs has
an exceptional lift in the finite unramified cover;
the preceding argument makes its image a point.
Thus all exceptional curves of the original
$\pi:T\to T_0$ would be contracted by the minimal
resolution $T\to N$. If $k>0$, the exceptional curve
of its last blowup is a smooth rational curve of
self-intersection $-1$, which a minimal resolution
cannot contract. Hence $k=0$.
There is no further exceptional contraction from
a minimal Hopf surface: every compact curve has
self-intersection zero by $b_2=0$, whereas a
nonempty exceptional configuration has negative
definite intersection matrix, as proved in P4R.6.
Therefore $T=N=T_0$ is a smooth minimal Hopf
surface. P4.24 gives $e(S)=0$ and its complete
elliptic-conductor description.
\end{proof}


\subsection{Excluding the remaining class-VII minimal models}

\begin{statement}{P4.37}{Proved}
No irreducible singular support in Part~4.2 has normalization
whose minimal resolution has a class-VII minimal model with
positive second Betti number. Consequently every class-VII
normalization occurring here is a smooth minimal Hopf surface
and has Euler number zero.
\end{statement}
\begin{proof}
The numerically anticanonical classification used in P4.23
and P4.26 reduces this question to an intermediate Kato
minimal model $T_0$: the Inoue--Hirzebruch and parabolic
alternatives were excluded there by their reduced
normal-crossing conductor and the Chern identity.
Write $n=b_2(T_0)>0$, let $A_0$ be its maximal reduced
curve divisor, and let $D_0$ be the effective integral
pole divisor pushed down from the normalization resolution.
Thus $K_{T_0}+D_0$ is numerically zero. The support of
$D_0$ is $A_0$.

We state the precise additional primary input.
Favre, \emph{Dynamique des applications rationnelles},
Theorem~1.2.24(1)--(2), printed p.~35, says that the
$n$ rational curve classes span $H^2(T_0;\mathbb Q)$
and that their intersection form is negative definite.
Theorem~1.2.30, printed p.~42, and its proof on
pp.~44 and~46 give a holomorphic foliation with conormal
line
\[
 N_{\mathcal F}^*=\mathcal O_{T_0}(-A_0)
                       \otimes\text{a flat line}.
\]
Theorem~1.2.32 and its proof on p.~46 show that every
component of $A_0$ is invariant and that its singularities
are precisely the intersection points, with local
first integrals $z^\alpha w^\beta$, $\alpha,\beta>0$.
The invariant-curve and order-one assertions are also
Dloussky--Oeljeklaus, Theorem~3.1, pp.~1528--1529.
The actual source URLs and exact theorem locators are
recorded in the bibliography and literature log.

It follows from the determinant of tangent and normal
lines that the foliation tangent line has numerical
first Chern class $-K_{T_0}-A_0$. The local generator
at an intersection is
$\beta z\partial_z-\alpha w\partial_w$.
Blowing up a regular point or one of these saddles
again gives only regular points and saddles with
opposite-sign nonzero eigenvalues. Thus every actual
blowup centre has the same two possibilities
$\varepsilon=0,1$ as in P4.36, with the same formulas.

Let $T\to T_0$ contain $k$ blowups. Write
$c_1(L)=-K_T+V$ for the saturated foliation line;
initially $V_0=-A_0$, and
$V'=\mu^*V+\varepsilon E$ at a blowup.
The reduced divisor $A_0$ is a rational cycle with
rational trees attached. Its normalization has
$n$ rational components and $n$ nodes, including
self-nodes, so $\chi(\mathcal O_{A_0})=0$.
Adjunction, or the divisor exact sequence on $T_0$,
therefore gives
\[
 (K_{T_0}+A_0)A_0=0.
\]
Hence $K_{T_0}V_0-V_0^2=0$.
Under each blowup this quantity changes by
$\varepsilon^2-\varepsilon=0$.
Since $K_{T_0}^2=-n$, we obtain
\[
 K_TV-V^2=0,\qquad K_T^2=-n-k.
\]

Define $B=\rho^*TY$, the divisorial determinant zero
$R$, the divisorial foliation zero $Z$, and the
nonnegative lengths $z_0,z_1$ exactly as in P4.36.
Its saturated sheaf calculation, including the
Hartogs extension of the rank-one double dual,
now gives
\[
 \mathcal B(R,Z)
 =R^2-K_TR-VR-ZR-K_TZ+2VZ+Z^2
 =n+k+z_0+z_1.                                  \tag{P4.37.1}
\]
Every reversed blowup changes $\mathcal B$ by at
most one, by the unrestricted integer identity
(P4.36.2). It remains to evaluate the initial term.

Push $R,Z$ down to effective divisors $R_0,Z_0$
on $T_0$. Their supports are contained in $A_0$:
the derivative is generically injective away from
the resolved conductor, and in any event all
compact curves of the intermediate surface are
components of its maximal divisor. Put $D=D_0$,
$A=A_0$. Using $K=-D$ numerically and $V=-A$ gives
\[
 \mathcal B_0
 =R_0^2+DR_0+AR_0-Z_0R_0+DZ_0-2AZ_0+Z_0^2.
\]
For divisor classes define
$q(X,Y)=X^2-XY+Y^2$.
Adjunction above gives $DA=A^2$, while $D^2=-n$.
Expansion consequently gives the exact completion
\[
 \mathcal B_0
       =n+q(R_0+D,Z_0+D-A).                     \tag{P4.37.2}
\]
The form on divisor classes is negative definite.
Moreover
\[
 q(X,Y)=\tfrac12\bigl(X^2+Y^2+(X-Y)^2\bigr),
\]
so it is nonpositive, and it vanishes only for
$X=Y=0$. The first argument $R_0+D$ is nonzero:
it is a nonzero effective combination of the
$n$ independent curve classes. Thus
$\mathcal B_0<n$.
After the $k$ blowups we still have
$\mathcal B(R,Z)<n+k$, contradicting (P4.37.1).
This excludes the intermediate minimal model.

All class-VII cases with positive second Betti
number have now been excluded. The zero-Betti
Inoue case has no conductor curves and was
excluded in P4.26. P4.36 treats the remaining
Hopf case and proves the final assertion.
\end{proof}


\subsection{Two exact local exclusions for ordinary-double curves}

\begin{statement}{P4.38}{Proved}
An actual irreducible support in Part~4.2 cannot have a singular
planar reduced critical curve at a point at which every generic
transverse germ is an ordinary double point. Nor can its reduced
critical curve there be the transverse union of a plane cusp
and a smooth branch, again with ordinary-double transverse germs
on both components. The latter assertion also holds with the
cusp replaced by a plane tacnode.
These exclusions allow singular normalization.
\end{statement}
\begin{proof}
We retain the boundary transport of P4.20 throughout. At a fixed
compact boundary, ordinary-double transverse germs form an open
condition. A sufficiently small deformation retaining that
condition gives the same nearby nonzero-level boundary
submersion. The finite-critical-value argument in that proof
then transports any compact vanishing sphere from the deformed
germ into the preceding nearby fibre. When a coordinate change
or a unit requires a smaller ball, we perform this construction
on that ball first, and then transport through the earlier
larger ball. Thus every sphere constructed below reaches the
original nearby four-torus.

First suppose that the reduced critical curve has ideal
$P=(z,h(x,y))$, where $h$ is reduced and singular at the
origin. The equation $f$ lies in $P^{(2)}$: its generic
transverse ordinary double point vanishes to order two on
each component. The complete-intersection square $P^2$ is
unmixed, so $P^{(2)}=P^2$. We therefore have the exact expression
\[
 f=A z^2+Bzh+Ch^2.
\]
Adding a small constant multiple of $z^2$ preserves the critical
curve and the ordinary-double boundary condition. We choose it
so that the $z^2$ coefficient at the origin is nonzero. The
Morse lemma in $z$, with parameters $x,y$, then gives
\[
 f=Z^2+h(x,y)^2u(x,y).
\]
Indeed the critical $z$ coordinate vanishes on $h=0$ and
is divisible by $h$, so its substitution has the displayed
square factor. Add a small constant to $u$ and shrink the
ball so that $u$ is a unit. Absorbing its square root in $h$
gives $Z^2+H(x,y)^2$, with the same singular reduced zero
curve. A reduced singular plane curve has an isolated critical
point. A generic small morsification $H_s$ has a Morse critical
point, and a small additive constant makes all these critical
values nonzero. At any such point $Z^2+H_s^2$ has a
nondegenerate isolated critical point: its plane Hessian is
$2H_s\operatorname{Hess}H_s$, and its $Z$ Hessian is $2$.
Its vanishing sphere has self-intersection $-2$.
The moving curve $H_s=0$ remains transverse to the chosen
boundary, with nondegenerate normal quadratic form there.
The preceding boundary transport therefore contradicts
$\pi_2(T^4)=0$.

For the second assertion write
\[
 h=y^2-x^k,\qquad P=(xz,yz,h),\qquad k=3\text{ or }4.
\]
This is the union of the plane curve $z=h=0$ and the
transverse line $x=y=0$. Its symbolic square is
\[
 P^{(2)}=(z,h)^2\cap(x,y)^2
  =(x^2z^2,xyz^2,y^2z^2,zh,h^2).
\]
For $k=4$, the first equality uses the two smooth branches of
$h=0$ separately; their complete-intersection square is unmixed
and hence is still $(z,h)^2$. To verify the intersection
formula directly, write an element of $(z,h)^2$ as
$az^2+bzh+ch^2$. The last two terms belong to $(x,y)^2$,
and $z$ is a nonzero divisor modulo this ideal. The first term
belongs to it exactly when $a\in(x,y)^2$.
Consequently
\[
 f=Azh+z^2(Bx^2+Cxy+Dy^2)+Gh^2.
\]
Small constant changes in $A,B$ retain the boundary condition
and make $A(0)B(0)\ne0$. Replace $h$ in this expression by
$h_s=y^2-x^k-sx^2$. The new critical curve contains
$z=h_s=0$ and $x=y=0$, and its boundary ordinary-double
germs persist. On the latter axis the transverse quadratic
matrix has coefficients
\[
 a=-sAz+Bz^2,\qquad b=Cz^2,\qquad c=Az+Dz^2,
\]
with off-diagonal entry $b/2$. Here the four coefficient
functions are restricted to that axis. Its determinant
divided by $z^2$ is
\[
 F(s,z)=(-sA+Bz)(A+Dz)-C^2z^2/4.
\]
At $(s,z)=(0,0)$ we have $F_z=AB\ne0$ and
$F_s=-A^2$. The implicit-function theorem gives a simple root
$z(s)=sA(0)/B(0)+O(s^2)$, nonzero for small nonzero $s$.
There the $y^2$ coefficient is nonzero, the Hessian has rank
one, and its determinant has a simple zero in the axis
direction. The parametrized Morse lemma gives
$Y^2+x^2H(x,z)$, with $H_z\ne0$. The exact coordinate
change $Z=H(x,z)$ yields $Y^2+x^2Z$.
P4.16 gives its compact spherical class of square $-4$.
Nested boundary transport once more places this class in
the original nearby four-torus, a contradiction.
\end{proof}

The exact quadratic matrix, determinant, and the two nonzero
derivatives in this proof are verified by
\texttt{code/part4/critical\_curve\_node.py}; its saved output
reports PASS.


\subsection{A coefficient-two exceptional curve over an ordinary-triple germ}

\begin{statement}{P4.39}{Proved}
Let $S=(f=0)$ be a reduced hypersurface germ in a smooth
threefold, with one-dimensional reduced critical locus $C$.
Suppose that every generic transverse germ along $C$ is an
ordinary triple point. On the minimal resolution $T$ of its
normalization, let $D$ be the integral total pole divisor
defined by the residue of an ambient volume form divided by
$f$. No exceptional curve of positive genus has coefficient
two in $D$.

In particular, the contracted horizontal section of a
positive-genus ruled resolution in Part~4.2 cannot lie over
a point where all local critical-curve branches have generic
type $D_4$. This assertion does not replace that local
hypothesis by a mixture of transverse types.
\end{statement}
\begin{proof}
We blow up the ambient origin, write $S_1$ for the strict
transform and $N_1$ for its normalization, and take the
minimal resolution $T_1$ of $N_1$. The proper map $N_1\to S$
lifts to the normalization $N$ by its universal property.
Thus $T_1\to N$ is a resolution and dominates $T$ by a
sequence of point blowups. The strict transform of a proposed
positive-genus exceptional curve $H$ remains present and has
the same pole coefficient two.

Put $d=\operatorname{mult}_0f$, and denote the ambient
exceptional divisor by $E$. Actual adjunction gives
\[
 \omega_{S_1}=\pi^*\omega_S\otimes
        \mathcal O((2-d)E)|_{S_1}.
\]
Let $D_1$ be the total conductor divisor on $T_1$ relative
to $\omega_{S_1}$, and let $e$ be the Cartier pullback of
$E$. With the specified residue comparison, the total pole
of the original germ on $T_1$ is
\[
                    D_1+(d-2)e.                       \tag{4.39a}
\]
The divisor $D_1$ is effective. Its strict conductor
coefficients are positive integers. For exceptional curves
of the minimal resolution of $N_1$, the negative-matrix
argument of P4R.6 applies locally: $K_{T_1}F\ge0$, since
there is no smooth rational exceptional $(-1)$-curve;
adding the strict conductor retains nonnegative intersection.
The inverse exceptional intersection matrix has nonpositive
entries, so all exceptional coefficients are nonnegative.

If $H$ is exceptional over $N_1$, its coefficient in $D_1$
is strictly positive. A zero coefficient would give
\[
 0=(K_{T_1}+D_1)H
   =2g(H)-2-H^2+D_1H>0.
\]
If that coefficient is one, adjunction gives
\[
 0=2g(H)-2+(D_1-H)H.
\]
Hence $H$ is elliptic and is disjoint from every other
positive component of $D_1$. A zero-coefficient exceptional
curve cannot meet a positive one, by its own intersection
equation and $K_{T_1}F\ge0$. The connected exceptional
fibre is consequently $H$ alone and meets no strict conductor.
The entire germ of $S_1$ at its image is normal with an
isolated simple elliptic singularity. Indeed, any additional
hypersurface branch would intersect this branch in a curve
and would give strict conductor meeting $H$.

Set $n=\operatorname{ord}_H(e)>0$ and
$a=\operatorname{coeff}_H(D_1)\ge0$. Formula (4.39a)
gives the exact identity
\[
                       2=(d-2)n+a.                    \tag{4.39b}
\]
Thus $d\ge5$ is impossible. For $d=4$ we have
$n=1,a=0$, and $H$ is not exceptional over $N_1$.
At its generic point the normalization is an isomorphism,
so $H$ is birational to a reduced component of the projective
quartic tangent cone. Every second derivative of $f$ vanishes
on $C$. Its leading term along any analytic arc of $C$
therefore makes that quartic have a point of multiplicity
at least three. An irreducible quartic has arithmetic genus
three and loses at least three there. A smooth elliptic cubic
component together with the remaining line could have
multiplicity at most two at the point. All components of
degree at most two are rational. Thus every reduced tangent
component is rational, excluding $H$ for $d=4$.

We are left with $d=3$ and
\[
                   (n,a)=(2,0)\quad\hbox{or}\quad(1,1).
\]
Choose a coordinate $u$ with nonzero pure cubic initial
coefficient, apply Weierstrass preparation, and depress the
monic cubic. Multiplication by the preparation unit does not
affect Hessian vanishing on $C$, because $f=df=0$ there.
We obtain
\[
 f=\text{unit}\cdot\bigl(u^3+b(x,z)u+c(x,z)\bigr).
\]
The second $u$ derivative is $6u$, so $C$ lies in the smooth
surface $u=0$. Write its reduced ideal as $(u,h(x,z))$.
The complete-intersection cube is unmixed and equals its
symbolic cube. Generic order three along $C$ and coefficient
comparison therefore give the exact expression
\[
 f=\text{unit}\cdot
       \bigl(u^3+h^2B(x,z)u+h^3C(x,z)\bigr).           \tag{4.39c}
\]
If $C$ is singular, $\operatorname{ord}h\ge2$ and the
tangent cone is the triple plane $u^3$. The degree-four
term restricts to zero on this plane. After the ambient
blowup the entire reduced exceptional line is critical,
so none of its points is an isolated normal singularity.
If $C$ is smooth, take $x=h$. Every homogeneous term $f_j$
then has degree at least three in $u,x$. The cubic tangent
cone is a binary cubic in these two variables, hence three
lines, counted with multiplicities, through the tangent
direction of $C$.

First suppose that $H$ dominates a tangent-cone line.
For $(n,a)=(2,0)$ the normalization is generically an
isomorphism, and $H$ is birational to that line. For
$(n,a)=(1,1)$, a transverse normalization branch has conductor
exponent one. It is smooth and meets exactly one other smooth
branch transversely: the intrinsic conductor exponent of a
singular plane branch is at least two, and each other branch
contributes a positive intersection multiplicity. Thus the
generic transverse germ of $S_1$ is a node. A multiplicity-one
tangent line is generically smooth and cannot have $a=1$.
For a multiplicity-three line, the two geometric branches of
that node have exceptional orders whose weighted sum is three.
A connected double cover would give equal orders and an even
sum. Thus the order-one branch has degree one over the line
and is rational.

The remaining possibility is a double line together with a
distinct simple line. Here $C$ is smooth by (4.39c). Make a
linear change so that $f_3=xu^2$. On the $x$ chart of the
blowup put $u=xv,z=xw$. The strict equation is
\[
 v^2+x f_4(1,v,w)+x^2f_5(1,v,w)+\cdots.
\]
An ordinary double curve on $x=v=0$ requires
$f_4(1,0,w)=0$ identically. Its transverse quadratic form is
\[
 v^2+xvB(w)+x^2C(w),
 \qquad \deg B\le1,\quad\deg C\le2.                  \tag{4.39d}
\]
The degree bounds follow from the degree-at-least-three
condition in $u,x$ on every $f_j$: here
$B=\partial_v f_4(1,0,w)$ and $C=f_5(1,0,w)$.
The normalization over the exceptional line adjoins a root
of $q^2+B(w)q+C(w)$. Its discriminant has degree at most
two. The resulting quadratic extension of $\mathbb C(w)$
is branched at at most two points, including infinity.
Riemann--Hurwitz makes its compact normalization rational;
the components of a split cover are rational as well.

Finally suppose that $H$ is exceptional over a point of $N_1$.
We must have $n=a=1$, and the first part of the proof would
make $S_1$ normal with an isolated simple elliptic singularity
there. This cannot happen. For singular $C$ the whole
exceptional line is critical, as already proved. For smooth
$C$ its strict transform passes through the common point of
the tangent lines, so that point is not an isolated normal
singularity. Away from it, at a line of multiplicity
$e=2$ or $3$, divide by the nonzero factor from any other
line. The strict equation has the form
\[
                         v^e+xG(v,w,x),
\]
where $G(0,w,0)$ is affine linear in $w$, by the same
normal-degree bound. If this function is nonzero at a point,
the surface is smooth. If it is identically zero, the entire
line is critical. At a simple zero, $G_w$ is a unit and the
exact coordinate change $W=G(v,w,x)$ gives $v^e+xW$, a
rational double point of type $A_{e-1}$. A multiplicity-one
line is smooth away from intersections. There is therefore
no isolated simple elliptic singularity. All possibilities
for $H$ have been excluded.
\end{proof}

The positive-integer solutions of (4.39b), the monomial degree
bounds in (4.39d), and the forced fourth-order factor of the
homogeneous quadratic discriminant are checked in
\texttt{code/part4/d4\_pole.py}. Its saved output reports PASS.


\subsection{The excess-one curve calculation and the ordinary-double sector}

\begin{statement}{P4.40}{Proved}
In the scope of Part~4.2 where every fibre support is
irreducible or normal crossing, no irreducible support with
positive-genus ruled resolution has generic transverse type
$A_1$ on every component of its critical curve.
\end{statement}
\begin{proof}
We first give the exact local curve calculation needed in
the proof. Let $A$ be a reduced complex curve local ring with
$r$ branches, let $B=\bigoplus_{i=1}^r\mathbb C\{t_i\}$
be its normalization, and let $R\subset B$ be the ring whose
constant terms agree. Then
\[
 \operatorname{length}(R/A)=\delta(A)-(r-1).
\]
Suppose this length is one. The simple $A$-module $R/A$
is annihilated by $\mathfrak m_A$, so
$\mathfrak m_A R\subset A$. Choose
$v\in\mathfrak m_R$ spanning $R/A$ and write
$v^2=a+cv$, where $a\in\mathfrak m_A$ and $c\in\mathbb C$.
If $c\ne0$, the convergent square-root series solves
$b^2-cb-a=0$ with $b\in\mathfrak m_A$. Then
$(v-b)(v+b-c)=0$ and the second factor is a unit, forcing
$v=b$, a contradiction. Thus $v^2\in A$ and
\[
                         \mathfrak m_R^2\subset A.
\]
Consequently $A$ is exactly the inverse image of a hyperplane
\[
                         \sum_i a_i f_i'(0)=0
\]
in $\mathfrak m_R/\mathfrak m_R^2=\mathbb C^r$.
Rescaling the parameters with nonzero $a_i$ gives actual
analytic ring isomorphisms with the following models.

With one nonzero coefficient, that branch has ring
$\mathbb C\{t^2,t^3\}$ and the others are independent smooth
branches joined at their constant terms. The embedding
dimension is $r+1$: for $r=1$ this is a plane cusp, for
$r=2$ a cusp and transverse line, and for $r=3$ it does
not embed in a smooth threefold. With two nonzero
coefficients, the pair has equal first derivatives. Its
ring is generated by $(t,t)$ and $(t^2,-t^2)$, giving the
plane tacnode $y^2=x^4$. The other $r-2$ branches are
independent transverse lines. For $r=3$ and three nonzero
coefficients, the first jets are the restrictions of linear
forms to three distinct lines in a plane. Their products
span all three second-jet coordinates, so the ring is
exactly the ordinary plane triple point. Equivalently,
the latter ring has precisely the same constant and
first-derivative conditions and colength three in $B$.
These cases exhaust $r\le3$. The analytic ring
isomorphisms lift to ambient coordinate changes: choose
minimal algebra generators, lift their independent linear
parts, and apply the inverse-function theorem. In embedding
dimension two we first straighten the smooth containing
surface. Every realizable model is therefore one of the
local classes excluded in P4.38 for ordinary-double
transverse type.

Now suppose the stated ruled support exists. Its horizontal
pole section has coefficient two. It cannot map
nonconstantly to the conductor, since an ordinary double
point has conductor exponent one on each normalization
branch. The section is contracted. P4.33 gives
$\chi(\mathcal O_{C_{\rm sch}})=-1$.
The conductor image is Cohen--Macaulay and generically
reduced in type $A_1$. It is therefore reduced, so
$C_{\rm sch}=C$ and $\chi(\mathcal O_C)=-1$.
All normalized components of $C$ are rational.

The already proved group and Euler calculations leave
exactly two normalized rational components and one
point-identification cycle. Indeed the rational-tree case
has Euler number at least three by P4R.4; the
one-component cycle is excluded by P4.35; and the global
nonnegative Euler sum is at most two. The normalization
sequence and the graph formula consequently give
\[
 \sum_p\delta_p=3,\qquad
 \sum_p(r_p-1)=2,
 \qquad
 \sum_p\bigl(\delta_p-(r_p-1)\bigr)=1.
\]
Every summand in the last sum is nonnegative. There is
exactly one excess-one point and every point has at most
three local branches. The local calculation classifies that
point by an actual analytic isomorphism, and P4.38 excludes
each possible model. This is the required contradiction.
\end{proof}

The first-jet cases, the second-jet product determinant, and
the global excess computation are verified in
\texttt{code/part4/curve\_excess.py}; its saved output reports
PASS. The proof retains the ordinary-double boundary
condition and makes no substitution of transverse type
$A_2$ for $A_1$.


\subsection{The ordinary-node defect in transverse type \texorpdfstring{$A_3$}{A3}}

\begin{statement}{P4.41}{Proved}
At a point of an actual support in Part~4.2, suppose the
entire reduced critical curve is the ordinary node
$C=\{z=xy=0\}$ and both branches have generic transverse
type $A_3$. Let $d_i,b_i,q$ be the Hessian zero orders,
projected conormal colengths, and normalization descent
defect of P4R.10. Then
\[
 \delta(C)+\sum_{i=1}^2\left(\frac{d_i}{2}+b_i\right)+q
                         \ge\frac92.                 \tag{4.41a}
\]
Here $R=0$ and $\delta(C)=1$. Equality makes the normalization
of the hypersurface smooth over the point.
\end{statement}
\begin{proof}
Put $P=(z,xy)$, and let $I$ be the ambient ideal of the
scheme-theoretic conductor image. We have $P^2=P^{(2)}\subset I$,
and $J=P/I$ is torsionfree of rank one on both branches.
After dividing its normalized image by the two branchwise
vanishing orders, it is an $A=\mathcal O_C$ submodule of
$B=\mathbb C\{x\}\oplus\mathbb C\{y\}$ generating both
summands. The conductor of $A$ in $B$ is $\mathfrak m_A$,
so $\mathfrak m_AB\subset J$. Thus $J$ is determined by
a subspace of $\mathbb C^2$ projecting onto both coordinates.
If that subspace is a line, rescaling gives $J\simeq A$ and
$q=1$. If it is the whole space, $J\simeq B$ and $q=0$.
These are actual module isomorphisms. The divided orders are
the $b_i$ already in the formula.

We use the ring condition of the actual conductor:
\[
 \operatorname{Hom}_{\mathcal O_S}(I,I)=
 \operatorname{Hom}_{\mathcal O_S}(I,\mathcal O_S).
                                                               \tag{4.41b}
\]
The right side is the normalization algebra by finite duality,
and $I$ is an ideal of that algebra. For the determinant
presentation of a Cohen--Macaulay codimension-two conductor
inside a Gorenstein hypersurface, this condition is equivalent
to membership of the adjugate entries in $I$. This is precisely
de Jong--van Straten, Theorem~1.12, printed p.~532
\cite{DeJongVanStraten1990Normalization}. Its hypotheses hold
here: the ambient space is regular, the reduced support is a
hypersurface, and its conductor quotient is Cohen--Macaulay.
We give all the matrices used below.

Suppose first that $q=1$. The surjection
$P/P^2\to J$ is a surjection $A^2\to A$, with primitive
kernel generator $g$. Taking a complementary generator $k$
of $P$ gives $I=(g,k^2)$. A determinant presentation
$f=ag+bk^2$ has two-by-two adjugate containing $a,b,g,k^2$.
Condition (4.41b) says $a,b\in I$, so the exact equation is
\[
                   f=Ag^2+Bgk^2+Ck^4.                \tag{4.41c}
\]
We claim that $A(0)$ cannot be a unit. If $g_z(0)\ne0$,
straighten $g$, take $k$ to be a nodal equation in $g=0$,
and apply the parameter Morse lemma in $g$. The residual
is $k^4$ times a unit after a small change in its fourth-order
coefficient. A morsification of the node and a shift of its
critical value away from zero give an isolated Morse point
of nonzero critical value.

If $g_z(0)=0$, its $xy$ coefficient is a unit and we take
$k=z$. Completing the square in (4.41c), including the
square root of $A$ in its first factor, gives exactly
\[
                         f=G^2+z^4U,
\]
with $G(x,y,0)=xy$ times a unit and $G_z(0)=0$.
Add a small constant to $U$, then take the actual coordinate
$Z=zU^{1/4}$. The parameter Morse lemma for $G$ in $x,y$
gives $G=XY+a(Z)$, where $a\in(Z^2)$. A small change
$a\mapsto a+\eta Z^2$ makes $a''(0)\ne0$. Replacing
$G$ by $G-\epsilon$ produces at the origin a critical point
of value $\epsilon^2$. The $XY$ Hessian entry is
$-2\epsilon$, and the $ZZ$ entry is
$-2\epsilon a''(0)$, so the point is Morse.
These deformations retain type $A_3$ along the moving
critical curve at each fixed boundary. The nested boundary
transport of P4.28 and P4R.14 places its square-$-2$ sphere
in the actual nearby torus, a contradiction. Hence $A(0)=0$.

If $g_z(0)$ is a unit, both saturated Hessian lines approach
$dz$, so $b_1=b_2=1$ and $d_i\ge1$. If $g_z(0)=0$,
the differential of $g$ on each branch has exactly the
common factor $t_i$ and a remaining nonzero coefficient
in the other normal direction. Thus $b_i=0,d_i\ge3$.
In either case the left side of (4.41a) is at least five,
including $q=1$.

Now let $q=0$. Choose actual generators $a,b$ of $P$ mapping
to $(1,0),(0,1)$ in $J\simeq B$. Their coefficient matrix
is invertible modulo the maximal ideal, hence invertible
locally. We obtain
\[
 I/P^2=(ya,xb),\qquad
 f=xU(x)b^2+yV(y)a^2\pmod {P^3}.                      \tag{4.41d}
\]
The second identity records the Hessian kernel $b$ on the
$x$ branch and $a$ on the $y$ branch. Write the exact
relation $xy=\alpha a+\beta b$. Near a general point of
the $x$ axis, solve this for $y$ using coordinates $x,a,b$.
The null cubic coefficient contains
$V(0)\alpha(x,0)/x$; the remaining contribution from
$P^3$ is holomorphic at the node. Its identical vanishing
in type $A_3$ forces $V(0)\alpha(0)=0$. The other branch
gives $U(0)\beta(0)=0$.

If both $a_z(0),b_z(0)$ are nonzero, the relation for $xy$
has $\alpha(0)\beta(0)\ne0$. Hence $U(0)=V(0)=0$,
both $b_i=1$, and both $d_i\ge2$. The defect is at least
five. Otherwise interchange the branches so that $a_z(0)$
is a unit and $b_z(0)=0$. Then $\alpha(0)=0$ and
$\beta(0)$ is a unit, forcing $U(0)=0$. On the $x$ axis,
$db$ has exactly the common factor $x$ and remaining
$dy$ coefficient a unit. On the $y$ axis, $da$ is primitive
with a unit $dz$ coefficient. Thus
\[
 b_x=0,\quad d_x=3+\operatorname{ord}_xU\ge4,
 \qquad
 b_y=1,\quad d_y=1+\operatorname{ord}_yV\ge1.
\]
This proves (4.41a). Equality requires
$\operatorname{ord}_xU=1$ and $V(0)\ne0$.

We identify the normalization in this exact equality case.
Use $a$ as the ambient coordinate $z$. The generator relation
with $\beta$ a unit gives
\[
                    I=(a^2,ya,x^2y-x\alpha a).
\]
Modulo its first two generators, replace $\alpha$ by
$\alpha(x,0,0)=x\ell(x)$. The actual coordinate
$Y=y-\ell(x)a$ puts the ideal into
$I=(z^2,Yz,x^2Y)$. Renaming $Y$ as $y$, its Hilbert--Burch
matrix and determinant extension are
\[
 \begin{pmatrix}y&0\\-z&x^2\\0&-z\end{pmatrix},
 \qquad
 M=\begin{pmatrix}y&0&a\\-z&x^2&b\\0&-z&c\end{pmatrix},
 \qquad
 \det M=az^2+byz+cx^2y.
\]
The new two-by-two minors are
$yb+za,yc,-zc,-x^2a,-za,x^2c+zb$, up to signs.
Their membership in $I$ gives exactly
\[
             a=yA,\qquad c=zC,\qquad b=-x^2C+zB
                          \pmod I.
\]
Therefore $f\in(yz^2)+I^2$, so
\[
           f=Uyz^2+Vz^4+Wx^2y^2z+Tx^4y^2.            \tag{4.41e}
\]
The equality orders $d_y=1,d_x=4$ say precisely
$U(0)T(0)\ne0$.

Substitute $y=z^2v$ and divide (4.41e) by $z^4$.
At the finite root its $v$ derivative is $U(0)$.
The implicit-function theorem supplies the actual smooth
factor $g_1=y-z^2v(x,z)$. Division by it gives
\[
                  g_2=z^2A+x^2yzB+x^4yC,
                  \qquad A(0)C(0)\ne0.
\]
For clarity, this division retains the factors $x^2,x^4$:
subtracting the expression at $y=z^2v$ supplies one factor
$y-z^2v$; every extra substituted $y$ supplies $z^2$ and
is absorbed into the first term. The parameter Morse
centre in $z$ is divisible by $x^2y$. Substitution leaves
$x^4y$ times a unit. An actual coordinate change therefore
gives $g_2=Z^2+x^4Y$. Its finite birational normalization is
\[
                 (x,u)\longmapsto(x,-u^2,x^2u),
\]
with target coordinates $(x,Y,Z)$. It is smooth; finiteness
follows from $u^2=-Y$ and birationality from $u=Z/x^2$
away from $x=0$. The other branch $g_1$ is smooth already.
Thus the full normalization is smooth over the node.
\end{proof}

All determinant minors, the four minimum-defect cases,
the factor substitution, and the normalization equation
are verified by \texttt{code/part4/a3\_node.py}.
The output reports PASS. Its initial failed assertion
compared an expanded polynomial with a factored expression;
that output is retained, and the corrected assertion tests
their exact zero difference.

\begin{statement}{P4.42}{Proved}
In the requested global support classes, a positive-genus
ruled-resolution support whose critical curve has generic
type $A_3$ everywhere cannot have two two-branch points as
its point-identification cycle. The two points may be
nonordinary. This includes additional
unibranch singular points of the critical curve.
\end{statement}
\begin{proof}
After P4R.13, the remaining positive ruled case has two
rational normalized conductor components and one graph
cycle. In the stated instance the cycle is provided by
two two-branch points. We first extend the local lower
bound of P4.41 to these points.

At a planar curve point, write its ideal as $(z,h)$.
Its symbolic square equals its ordinary square, so the
rank-one Hessian belongs to $\operatorname{Sym}^2 E_0$,
where $E_0$ is the normalized conormal image. On a branch
put $E_0=\langle t^a e_1,t^b e_2\rangle$ and write the
primitive Hessian direction as $\ell=u e_1+v e_2$.
One of $u,v$ is a unit. Membership of
$t^d\ell^2$ in $\operatorname{Sym}^2E_0$ and projection
onto the quotient by $\ell$ give
\[
 \frac d2\ge\max\{a-\operatorname{ord}u,
                         b-\operatorname{ord}v\},\qquad
 b_{\rm proj}=\min\{a+\operatorname{ord}v,
                         b+\operatorname{ord}u\}.
\]
The sum of the maximum and minimum on the right is
$a+b$. Thus $d/2+b_{\rm proj}\ge a+b$.
The plane different formula proved in P4.34 is
$\sum(a+b)=2\delta+\deg R$. The whole local defect
is consequently at least $3\delta+q$.
A nonordinary planar two-branch point has $\delta\ge2$
and contributes at least six.

For a nonplanar point, the ambient Hessian vanishes: a
nonzero rank-one Hessian would give $f=X^2+g(Y,Z)$ by
the parameter Morse lemma and would put the entire
critical curve in the smooth plane $X=0$.
On a branch of multiplicity $m$, all conormal generators
have order at least $m$, and every Hessian entry has
order at least $m$. Thus $b_{\rm proj}\ge m$, $d\ge m$,
and, since $r=m-1$,
\[
                  d/2+b_{\rm proj}-r\ge1+m/2\ge1.
\]
When $\delta\ge3$, a nonplanar two-branch point
therefore contributes at least five.
The only remaining possibility has $\delta=2$.
The exact excess-one classification in P4.40 identifies
it with a cusp and transverse line,
\[
                  P=(xz,yz,h),\qquad h=y^2-x^3.
\]
Here $P^{(2)}/P^2$ is generated by $zh$ and is
annihilated by the maximal ideal. Write $f=Azh+P^2$.
Along the cusp, the normal quadratic determinant is
$4KG-A^2$, with $K(0)=0$. Its rank is one generically,
so $A(0)=0$ and in fact $f\in P^2$.
The conormal lattices on the two normalized branches
are respectively $\operatorname{diag}(t^2,t^3)$ and
$\operatorname{diag}(t,t)$. Their total colength is seven,
while $\deg R=1$ and $\delta=2$.
Applying the same symmetric-square inequality gives
defect at least $2+7-1=8$.
Every two-branch point thus contributes at least $9/2$,
and equality is possible only at the ordinary node of
P4.41, whose normalization is smooth.

Apply the rank-one Hessian calculation
of P4R.10 to the two normalized projective lines. On each
line the rank-two conormal kernel has degree $2+\deg R$.
The rank-one quotient has Euler characteristic
$3+\deg R-\deg L$. Combining both lines with the
normalization sequence of $C$ gives the exact formula
\[
 \chi(\mathcal O_{C_{\rm sch}})
   =8-\sum_p\left[
       \delta_p+\sum_{i\text{ over }p}
          \left(\frac{d_i}{2}+b_i-r_i\right)+q_p
                 \right].                            \tag{4.42a}
\]
Every extra unibranch singular point has a strictly positive
bracket by the local conormal inequalities of P4R.10.
At a smooth point the bracket is nonnegative. The two
two-branch points contribute at least nine together by the
local calculation just given.

If the horizontal section is not contracted, P4.33 gives
$\chi(\mathcal O_{C_{\rm sch}})=g-1\ge0$, contradicting
(4.42a), whose right side is at most $-1$.
If the section is contracted, that Euler characteristic is
$-1$. Equality must then hold at both two-branch points and every
other bracket must vanish. There are no additional
unibranch singular points. Both cycle points are ordinary
nodes and the normalization is smooth over them by P4.41.
At every smooth point of the
critical curve, P4R.14 gives either the product tacnode
or the displayed $A_1$ normalization degeneration. The
latter has Hessian zero order one and contributes $1/2$,
so it is excluded by the equality in (4.42a). Thus the
normalization is smooth along all of the conductor, and
it is smooth away from the conductor by P4R.6. A contracted
positive-genus horizontal section in its minimal resolution
is impossible. This proves the assertion.
\end{proof}


\subsection{The uniform cuspidal sector}

\begin{statement}{P4.43}{Proved}
In the requested full-fibration classes, a support with
positive-genus ruled resolution cannot have positive
Euler number and generic transverse type $A_2$ along
every component of its critical curve.
\end{statement}
\begin{proof}
We use the exact normalization map
$\nu:N\longrightarrow S$ and its minimal resolution
$\rho:T\longrightarrow N$. A germ of $S$ cannot have
two distinct analytic hypersurface branches. Indeed,
their intersection has dimension one in the smooth
ambient threefold. A general point of an irreducible
component of that intersection has at least two local
hypersurface branches, so its transverse curve germ is
reducible. The transverse cusp $A_2$ is irreducible.
This contradicts the stated uniform transverse type.
The points of the finite normalization over a germ are
its analytic branches. Consequently $\nu$ is bijective.
It is proper, so its continuous map from a compact space
to the Hausdorff space $S$ is a homeomorphism.

Let $B$ be the base curve of the minimal ruled surface,
of genus $g\ge1$, and let $H$ be the horizontal section
of coefficient two in the total pole divisor. Point
blowups preserve the fundamental group, and the ruling
gives $\pi_1(T)=\pi_1(B)$. By the fibre-tree calculation
in P4R.7, every vertical exceptional configuration is a
tree of rational curves. Collapsing any such connected
tree preserves the fundamental group: apply van Kampen
to its tubular neighbourhood and the complement, or
equivalently to the cofibration quotient by that
simply connected tree. These quotients are the actual
topological contractions of the resolution map.

If $H$ is not contracted by $\rho$, all resolution
exceptional configurations are of this vertical kind.
It follows that $\pi_1(N)=\pi_1(B)$. If $H$ is
contracted, its connected exceptional configuration
contains $H$. The inclusion $H\hookrightarrow T$
induces an isomorphism on fundamental groups, since
the ruling restricts to the identity section and all
point blowups preserve that isomorphism. The quotient
which collapses this configuration therefore kills
the whole group. The other contracted configurations
are rational trees, giving $\pi_1(N)=1$.

On the other hand, the proved positive ruled reduction,
including P4R.13, gives a conductor curve $C$ with two
rational normalized components and one graph cycle.
The conductor cofibration pushout in P4R.7 gives
$\pi_1(S)=\pi_1(C)=\mathbb Z$. This contradicts the
homeomorphism $N\simeq S$: neither the trivial group
nor the genus-$g$ surface group is $\mathbb Z$, as their
abelianizations have ranks zero and $2g$, respectively.
\end{proof}


\subsection{Further uniform transverse sectors}

\begin{statement}{P4.44}{Proved}
In the remaining positive ruled case with generic transverse
type $A_3$ on every conductor component, the horizontal
section of the pole divisor is contracted by the resolution
map. The point-identification cycle occurs at a single
three-branch point of the critical curve.
\end{statement}
\begin{proof}
P4.42 excludes a cycle at two two-branch points. The
remaining graph has two rational normalized components
and one cycle, so its cycle occurs at one three-branch
point $p$. Suppose that the horizontal section $H$ is
not contracted. It maps finitely to one normalized
conductor component $C_H^\nu=\mathbb P^1$.
The generic tacnode has two conductor preimages.
The degree of $H\to\mathbb P^1$ is greater than one,
since $H$ has genus $g\ge1$, and hence is exactly two.
No other, vertical conductor component can map onto
$C_H$: the horizontal map already supplies both generic
preimages.

At a point where the entire critical curve is smooth,
P4R.14 gives either the product tacnode or the germ
$X(Y^2+XZ)$. Each has two distinct normalization points
over the critical-curve point. Each point lies on the
strict conductor, by P4R.6. Thus the two preimages must
both lie on $H$, and the degree-two map is unramified
there. All ramification is over singular points of $C$.

By (4.42a) and P4.33 the sum of the local defects is
$9-g\le8$. A nonplanar three-branch point has
$\delta\ge2$, and the branchwise estimate in P4.42
adds at least one per branch; its defect is at least
five. A planar three-branch point has $\delta\ge3$
and contributes at least $3\delta\ge9$, which is
already impossible. Every additional unibranch singular
point contributes at least two: in the planar case
the bound is $3\delta$, and in the nonplanar case it is
at least $\delta+1$. There is consequently at most one
additional unibranch singular point.

The normalization $C_H^\nu$ has at most two points
above $p$, since the other global conductor component
supplies at least one of its three branches. Including
the possible additional unibranch point leaves at most
three possible branch values for $H\to\mathbb P^1$.
Every ramification point of a degree-two map has index
two, and its branch value has only that ramified
preimage. Riemann--Hurwitz requires $2g+2\ge4$ distinct
branch values. This contradiction proves that $H$ is
contracted.
\end{proof}

\begin{statement}{P4.45}{Proved}
In an actual torus degeneration, a point of a planar
singular critical curve cannot have generic transverse
type $D_4$ along all its local branches.
\end{statement}
\begin{proof}
Write the full reduced critical-curve ideal as
$P=(u,h(x,y))$, where $h$ is reduced and singular.
Generic multiplicity three gives $f\in P^{(3)}$.
Powers of this complete-intersection prime-by-components
ideal are unmixed, so $P^{(3)}=P^3$. Thus there are
holomorphic coefficients such that
\[
 f=A u^3+B u^2h+Cuh^2+Dh^3.
\]
Introduce an independent variable $H$ and the actual
holomorphic function
\[
 F(u,x,y,H)=A(u,x,y)u^3+B(u,x,y)u^2H
               +C(u,x,y)uH^2+D(u,x,y)H^3.
\]
Then $f=F(u,x,y,h)$. Arbitrarily small constant changes
of the four coefficients preserve the $D_4$ transverse
germs on a fixed boundary and make the binary cubic
at the origin have distinct roots and a simple critical
point as a polynomial in $u/H$. Its critical value is
nonzero: a zero critical value would be a repeated root.
We also choose its leading coefficient nonzero.

After substituting $u=Hv$, the quotient $F_u/H^2$
is holomorphic. The implicit-function theorem gives
a holomorphic critical root $v=v(x,y,H)$ and
\[
          F(Hv(x,y,H),x,y,H)=H^3\beta(x,y,H),
                       \qquad \beta(0)\ne0.
\]
Choose the local cube root of the unit $\beta$.
The map $H\mapsto H\beta(x,y,H)^{1/3}$ has nonzero
$H$ derivative at the origin and has a holomorphic
inverse with parameters $x,y$. Put
\[
             \widehat h(x,y)=h(x,y)
                           \beta(x,y,h(x,y))^{1/3}.
\]
It defines the same singular reduced plane curve as $h$.
Choose a small morsification $\widehat h_s$ with a
Morse critical point of nonzero critical value. This
exists by the positive local degree of the gradient
of the singular reduced plane function; a small constant
shift avoids zero critical values. Use the displayed
inverse to solve exactly
\[
         H_s\beta(x,y,H_s)^{1/3}=\widehat h_s,
             \qquad f_s=F(u,x,y,H_s(x,y)).
\]
On its $u$-critical branch the critical-value function
is exactly $\widehat h_s^3$. At the selected plane
Morse point, $H_s\ne0$ and $F_{uu}\ne0$. The Schur
complement of the full Hessian is
$3\widehat h_s^2\operatorname{Hess}(\widehat h_s)$.
Consequently $f_s$ has a nondegenerate three-variable
critical point with nonzero critical value.

All deformations above can be chosen successively small
on the fixed boundary. There the moving curve $H_s=0$
has the same ordinary-triple transverse germ: its binary
cubic has three distinct linear factors. The boundary
transport used in P4.28 therefore identifies the local
nearby fibres through nested representatives. The Morse
vanishing sphere embeds in a nearby global torus fibre
and has self-intersection $-2$, contradicting its
asphericity. This proves the assertion.
\end{proof}

\begin{statement}{P4.46}{Proved}
In the requested full-fibration classes, a support with
positive-genus ruled resolution cannot have positive
Euler number and generic transverse type $D_4$ along
every component of its critical curve.
\end{statement}
\begin{proof}
If the horizontal section is contracted, P4.39 already
gives a contradiction. Suppose it is not contracted.
The remaining critical curve has two rational normalized
components and one graph cycle. For the actual
conductor scheme, ordinary-triple transverse type gives
the ideal $P^{(2)}$, as in P4.31. Let $Q=P/P^{(2)}$,
and let $\tau_p$ be the full colength of its differentiated
image in the direct image of the normalized rank-two
kernel $K$. This includes both the branchwise lattice
colength and the further descent colength. Riemann--Roch
on the two normalized projective lines gives
\[
 \chi(\mathcal O_{C_{\rm sch}})
       =10-\sum_p(\delta_p+\tau_p-r_p).
\]
P4.33 identifies the left side with $g-1$, so the total
defect is $11-g\le10$.

We prove that a three-branch point contributes at least
eleven. A planar point contributes $5\delta\ge15$ by
the exact different calculation in P4.34. Consider a
nonplanar point. If its three branches are smooth, the
curve has multiplicity three and embedding dimension
three. A generic linear parameter has an Artinian
section of length three and embedding dimension two,
with ideal $(v,w)^2$. The defining curve ideal therefore
has exactly three minimal generators. To justify the
lifting, the parameter is a nonzerodivisor on the
Cohen--Macaulay curve ring, so the Tor term vanishes in
the ideal sequence modulo that parameter; Nakayama
then lifts precisely the three generators.
Each branchwise conormal lattice has rank two and
colength at least two, since every defining gradient
has positive order. Their sum has colength at least
six. Modulo the three branch maximal ideals its
normalization has six generator values. The image of
$Q$, generated by three elements over the local curve
ring, supplies a vector space of dimension at most
three. The additional descent colength is therefore
at least three. Since $\delta\ge2$ and $R=0$, the
full defect is at least $2+6+3=11$.

If one of the three branches is singular, let $m\ge4$
be total multiplicity. The branchwise gradient bound
gives $\tau\ge2m$, while $\deg R=m-3$.
The general inequality $\delta\ge m-1$ leaves only
$m=4,\delta=3$ as a possible value below the required
bound. That value has excess $\delta-(3-1)=1$.
The exact classification in P4.40 shows that its only
singular-branch model is a cusp and two independent
lines, of embedding dimension four. It does not embed
in the ambient smooth threefold. Thus $\delta\ge4$
in the present case, and the full defect is at least
$\delta+2m-(m-3)=\delta+m+3\ge11$.
A cycle at one three-branch point is impossible.

The other possible cycle has two two-branch points.
Each contributes at least five by P4.34; equality is
possible only for a planar ordinary node. A nonplanar
two-branch point contributes at least seven, and a
nonordinary planar point at least ten. The total bound
of ten therefore forces two ordinary planar nodes
and $g=1$. Each such node is forbidden by P4.45.
This proves the assertion.
\end{proof}

The local Hessian identities, rank-one lattice identity,
and exact integer defect bounds are checked in
\texttt{code/part4/uniform\_sectors.py}; the saved output
reports PASS.


\subsection{The full quadratic conductor at three axes}

\begin{statement}{P4.47}{Proved}
Consider an actual torus degeneration at a point where the
entire reduced critical curve is the three coordinate axes,
with generic transverse type $A_3$ on each axis. Suppose
the local Hessian--conductor defect is nine, the three
Hessian orders are two, and the quartic equation is the
square of $a xy+b xz+c yz$ with $abc\ne0$.
Then the normalization has only rational singularities
over that point. In particular it cannot contain the
contraction of a positive-genus horizontal section.
\end{statement}
\begin{proof}
Put $P=(xy,xz,yz)$ and $h=xyz$. The symbolic square is
$P^{(2)}=P^2+(h)$, and the actual conductor image ideal
$I$ contains $P^{(2)}$. Also $f\in P^{(2)}$.
Write $f=Ah+P^2$. On the $x$-axis, a nonzero constant
term of $A$ would give a leading rank-two transverse
Hessian, since the quadratic term is then a nonzero
multiple of $xyz$. Generic type $A_3$ has rank one.
Thus $A(0)=0$. Since $\mathfrak m h\subset P^2$,
we have $f\in P^2$.

The normalized conormal lattice is $tK$ on each axis,
so its projected colength is one. The curve has
$\delta=2$ and $R=0$. The stated defect therefore
forces the remaining normalization descent colength
to be one. The normalized quotient values form a
two-dimensional subspace of $\mathbb C^3$ with one
relation involving all three coordinates. Multiplying
a source generator that is a unit on one branch by
that branch's axis coordinate generates its positive
orders alone. Hence all positive-order branch functions
occur in the image, and the leading kernel vector
lifts to an actual element $g\in I$ with quadratic
part $a xy+b xz+c yz$.

We prove the exact ideal equality $I=(g,h)$.
On each coordinate plane, $g$ is the product of its
two coordinate functions times a unit. Thus
$I_0=(g,h)$ has radical $P$. Furthermore $P^2\subset I_0$:
modulo $h$, multiplication of $g$ by $xy$ gives a unit
times $(xy)^2$, and the other two squares follow in
the same way; all mixed products are multiples of $h$.
At the generic point of every axis, $I_0$ is the
tacnode conductor ideal, generated by a normal coordinate
and the square of the null coordinate. It agrees
generically with $I$. The ideal $I_0$ is a
codimension-two complete intersection, so
$\mathcal O/I_0$ is Cohen--Macaulay and $P/I_0$ is
torsionfree on the reduced curve. The map
$P/I_0\to P/I$ has finite-length kernel, which must
vanish. This proves $I=I_0$.

The actual conductor ring condition, with the hypotheses
verified in P4.41, gives $f\in I^2$ by the two-by-two
adjugate calculation of de Jong--van Straten,
Theorem~1.12. Write $f=A g^2+Bgh+Ch^2$.
Its quartic part makes $A(0)\ne0$. Completing the square
and absorbing $\sqrt A$ into the first factor gives
the exact equation
\[
                         f=G^2+h^2U.                 \tag{4.47a}
\]
The quadratic part of $G$ is nondegenerate, and
$(G,h)=(g,h)$. If $U(0)\ne0$, the two factors
$G\pm h\sqrt{-U}$ have nondegenerate quadratic
Hessians. Each is an isolated $A_1$ surface singularity,
so the full normalization has rational singularities.

We exclude $U(0)=0$ by a deformation that keeps the
three axes fixed. Add an arbitrarily small generic
linear form to $U$, and then replace $U$ by
$U-\epsilon$. On a fixed boundary the transverse
$A_3$ coefficient is a unit, so these changes retain
the boundary transverse type. A diagonal linear change
puts the quadratic part of $G$ in the form
$G_2=xy+xz+yz$, changing $h$ only by a nonzero constant.
Let $L$ be the resulting linear part of $U$.
Under $(x,y,z)=\epsilon(X,Y,Z)$, equation (4.47a) becomes
\[
 \frac{f_\epsilon}{\epsilon^4}
   =G_\epsilon(X,Y,Z)^2+\epsilon^3(XYZ)^2
                          \bigl(L(X,Y,Z)-1+O(\epsilon)\bigr),
 \qquad G_\epsilon=G_2+O(\epsilon).
\]
Near a nonzero point of the smooth cone $G_2=0$,
use $G_\epsilon$ as the normal coordinate. The
parameter Morse centre is $O(\epsilon^3)$.
Substitution leaves the tangential function
$\epsilon^3(\Phi+O(\epsilon))$, where
\[
                 \Phi=(XYZ)^2(L-1)|_{G_2=0}.
\]

We verify a nonempty open set of choices of $L$ for
which $\Phi$ has a nondegenerate critical point of
nonzero value. Parametrize the projective conic by
\[
              [X:Y:Z]=[t(t+1):t+1:-t].
\]
For $L=X+Y+Z$, set
$H(t)=-t^2(t+1)^2$ and $\ell(t)=t^2+t+1$.
In radial coordinate $r$,
$\Phi=r^6H(t)^2(r\ell(t)-1)$.
Radial criticality gives $r=6/(7\ell(t))$, and
\[
 \frac{d}{dt}\left(\frac{H}{\ell^3}\right)
  =\frac{t(t-1)(t+1)(t+2)(2t+1)}{(t^2+t+1)^4}.
\]
The three critical points $t=1,-2,-1/2$ are simple
and avoid zeros and poles. Their radial coordinates
are $2/7,2/7,8/7$. At each point,
\[
       \Phi=-\frac{1024}{823543},\qquad
       \det\operatorname{Hess}_{r,t}\Phi
                     =\frac{4718592}{1977326743}.
\]
These numbers are nonzero. Simplicity and avoidance
of zeros and poles are nonvanishing algebraic conditions
on the coefficients of $L$. The explicit example proves
that they hold on a nonempty Zariski open set, so an
arbitrarily small generic linear perturbation of any
initial linear part is sufficient.

After eliminating the normal Morse coordinate, divide
the two tangential critical equations by $\epsilon^3$.
They are analytic in $\epsilon$ and have invertible
Jacobian at the displayed points. The implicit-function
theorem lifts them to nondegenerate three-variable
critical points of $f_\epsilon$. Their critical values
have leading term $\epsilon^7$ times the displayed
nonzero value. They tend to the original point.
Since the deformation retains boundary $A_3$, the
nested boundary transport embeds a Morse vanishing
sphere of square $-2$ into a nearby global torus fibre.
This contradicts asphericity. Therefore $U(0)=0$
cannot occur, proving the assertion.
\end{proof}

The conic derivative, all three exact critical values
and Hessian determinants, and the related conductor
matrices are executed in
\texttt{code/part4/a3\_three\_axes.py}.
The output reports PASS. The working record explicitly
rejects the earlier proposed change $G\mapsto G-\epsilon$:
with $h$ fixed that change replaces the boundary tacnode
by two ordinary double points and does not supply the
required boundary transport.


\subsection{The three-branch tacnode defect budget}

\begin{statement}{P4.48}{Proved}
At a planar three-branch critical-curve point with uniform
generic transverse type $A_3$, the local conductor defect
is at least ten. At an ordinary nonplanar three-branch
point, the defect is at least nine. Equality in the
second assertion forces Hessian orders $(2,2,2)$,
projected colengths $(1,1,1)$, descent colength one,
and a quartic square $(a xy+b xz+c yz)^2$ with at
least two nonzero coefficients. The case $abc\ne0$
has rational normalization by P4.47.
\end{statement}
\begin{proof}
For a planar critical curve the ideal is $P=(u,h)$.
The rank-one nilradical $J=P/I$ of the conductor image
is generated by at most two elements over $\mathcal O_C$.
Its normalized image has at least one generator value
on each of the three branches. Modulo the three branch
maximal ideals, the source supplies a vector space of
dimension at most two, while the normalized image
supplies dimension three. Hence its descent colength
$q$ is at least one. The planar symmetric-square and
different calculation in P4.42 gives local defect
$\ge3\delta+q$. Since a planar three-branch point has
$\delta\ge3$, its defect is at least ten.

For an ordinary nonplanar point use actual coordinates
with $P=(xy,xz,yz)$. The argument in P4.47 gives
$f\in P^2$. On each axis, its normalized conormal
lattice is exactly $tK$, so the projected colength
is exactly one. Every Hessian order is at least two.
Write the quartic form in the basis $(xy,xz,yz)$ as
\[
 A(xy)^2+C(xz)^2+E(yz)^2
      +2B(xy)(xz)+2D(xy)(yz)+2F(xz)(yz).
\]
The rank-one Hessians give
$AC=B^2$, $AE=D^2$, $CE=F^2$.
If a diagonal coefficient vanishes, its adjacent
off-diagonal coefficients vanish and the remaining
two-by-two block has rank at most one. If all three
diagonal coefficients are nonzero, the vanishing null
cubic on the $x$-axis gives $F=BD/A$. Higher-degree
terms cannot alter the coefficient linear in $x$
of that normal cubic. Thus the whole quartic is a
square $(a xy+b xz+c yz)^2$.

If the quartic is zero, all three Hessian orders are
at least three, and the local defect is at least
$2+3+9/2=19/2$. If the square has just one nonzero
coefficient, say $a xy$, the $z$-axis Hessian order
is at least three. In the source-generator basis
$(xy,xz,yz)$, the three leading quotient-value rows
have the forms
\[
                      (0,1,0),\qquad(0,0,1),
                      \qquad(0,*,*).
\]
Their rank is at most two, so $q\ge1$ and again
the defect is at least $19/2$.

If at least two coefficients are nonzero, all three
Hessian orders are two. The leading quotient-value
matrix has rank two. For example, up to rescaling
its rows it is
\[
 \begin{pmatrix}-b&a&0\\-c&0&a\\0&-c&b\end{pmatrix}.
\]
It has zero determinant and rank two whenever at
least two of $a,b,c$ are nonzero. Therefore $q\ge1$,
and the local defect is at least $2+3+3+1=9$.
Equality forces $q=1$ and the displayed Hessian and
projection orders. P4.47 applies when all three
quadratic coefficients are nonzero.
\end{proof}


\subsection{The remaining ordinary-axis conductor}

\begin{statement}{P4.49}{Proved}
At an ordinary three-axis critical-curve point with uniform
generic transverse type $A_3$ and local defect nine, the
normalization has rational singularities. This includes the
quartic square with exactly one missing coefficient left
by P4.47--P4.48.
\end{statement}
\begin{proof}
Only the one-missing-coefficient case remains. Rescale
its leading conductor form to $x(y+z)$.
As in P4.47, the descent colength one lifts the leading
kernel vector to $g\in I$ with this quadratic part.
Write $g=x(\alpha y+\beta z)+yz\gamma$, where
$\alpha,\beta$ are units and $\gamma(0)=0$.
Replacing $y,z$ by $\alpha y,\beta z$ is an analytic
coordinate change preserving the axes. Put $k=y+z$,
$h=yz$. Then $g=xk-hc$, with $c(0)=0$.
The ideal $I_0=(g,xh,h^2)$ lies in the actual conductor
$I$ and has Hilbert--Burch matrix
\[
                 \begin{pmatrix}h&0\\-k&h\\c&-x\end{pmatrix}.
\]
It has radical $(xy,xz,yz)$, is Cohen--Macaulay, and
agrees with $I$ generically on each axis. The
torsionfree finite-kernel comparison used in P4.47
therefore gives the exact equality $I=I_0$.

There are exactly two required ideal forms. Since
$\mathcal O/(x,k,h)=\mathbb C\{y\}/(y^2)$, analytic
division writes $c=\lambda y+xA+kB+hC$ for a scalar
$\lambda$. Replace $g$ by $g+xhA+h^2C$, and then
set $x'=x-hB$. These generator and coordinate changes
give
\[
                    I=(x'k-\lambda yh,x'h,h^2).
\]
If $\lambda=0$ this is $(xk,xh,h^2)$. If
$\lambda\ne0$, rescaling $x'=\lambda X$ gives
$(Xk-yh,Xh,h^2)$. Every change is invertible near
the origin. We treat the two forms separately.

\paragraph{The ideal $(xk,xh,h^2)$.}
Appending $(a,b,c)$ to its Hilbert--Burch matrix gives
determinant $axk+bxh+ch^2$. The adjugate ring condition
of de Jong--van Straten, Theorem~1.12, applies to the
actual conductor as in P4.41. The colon ideals
\[
 I:h=(x,h),\qquad I:x=(k,h),\qquad I:k=(x,h^2)
\]
give exactly
$f=A x^2k^2+B x^2kh+Cxh^2+Dh^4$, with $A(0)\ne0$.
Complete the square in $k$, absorb $\sqrt A$ into $x$,
and replace $y$ by $y(1+\gamma z)$ to make
$k+\gamma h$ the new $y+z$. These are analytic
coordinate maps; $h$ changes only by a unit.
Absorbing that unit in the coefficients gives
\[
                       f=x^2k^2+xh^2C+h^4D.          \tag{4.49a}
\]
If $C(0)\ne0$, substitute $x=h^2v$ and divide by
$h^4$. The finite $v$-root is simple, giving a smooth
factor $x-h^2v(y,z)$. Its quotient is $xk^2+h^2U$
with $U$ a unit. The coordinate $X=x/U$ turns it
into $Xk^2+h^2$. Its finite birational normalization is
\[
                 u=h/k,\qquad X=-u^2,
                 \qquad u(y+z)=yz.
\]
The last equation is an $A_1$ quadric, so the
normalization has rational singularities.

Suppose $C(0)=0$. Perturb $D$ by a small constant
and the linear part of $C$ generically, retaining
$C(0)=0$, and then replace $C$ by $C-\epsilon$.
On a fixed boundary the transverse tacnode coefficients
are units, so these changes retain boundary $A_3$.
Under $x=\epsilon^3X$, $y=\epsilon Y$,
$z=\epsilon Z$, the leading scaled equation is
\[
 \frac f{\epsilon^8}=F+O(\epsilon),\qquad
 F=X^2K^2+XH^2(L-1)+D_0H^4,
 \quad K=Y+Z,\ H=YZ.
\]
For $L=K$, $D_0=1$, put $W=Y-Z$.
At $W=0$ the $X$-critical coordinate is
$X=-K^2(K-1)/32$, and the remaining derivative is
$K^5(12K^2+7K-3)/512$.
At either root of $12K^2+7K-3$, the full critical
value and Hessian determinant are respectively
\[
 \frac{K^6(K+1)(3K-1)}{1024},\qquad
 -\frac{3K^{10}(K+1)(3K-1)(28K^2+14K-5)}{32768}.
\]
They are nonzero: their numerator resultants with the
critical polynomial are $8748$ and $4923269424$.
Nondegeneracy, nonzero value and avoidance of the axes
hold on a nonempty algebraic open set of coefficients.
An arbitrarily small generic perturbation is therefore
available. The implicit-function theorem lifts these
Morse points through the higher weighted terms, with
nonzero critical values of order $\epsilon^8$.
The boundary comparison embeds their square-$-2$
vanishing spheres in a nearby torus fibre, a
contradiction. Thus $C(0)=0$ is excluded.

\paragraph{The ideal $(g,xh,h^2)$, where $g=xk-yh$.}
Append $(a,b,d)$ to
\[
                         \begin{pmatrix}h&0\\-k&h\\y&-x\end{pmatrix}.
\]
Its determinant is $ag+bxh+dh^2$. The new adjugate
entries give the conditions
\[
       ax,\ ah,\ bx+dh,\ kd+by,\ dh-ay,\ ak+bh\in I.
\]
The exact colon identities are
\[
 I:h=(x,h),\qquad I:x=(xk,h,z^2),\qquad I:x^2=(k,h).
\]
The first two conditions give $a=hA$ modulo $I$;
then $dh-ay$ gives $d=Ay+xD+hE$.
Initially $bx+dh$ permits
$b=-Ak+hB+z^2B_2$ modulo $I$.
The condition $ak+bh$ forces $z^2B_2\in(x,h)$,
so $B_2\in(x,y)$. Its $x$-part lies in $I$, since
$xz^2=h^2-xh$ modulo $I$, and its $y$-part is
absorbed in $hB$. Thus $b=-Ak+hB$ modulo $I$.
The remaining condition is $y(D+B)+kE\in(x,h)$.

Write $T=D+B$. Modulo $(x,h,k)$, the equality
$yT=0$ in $\mathbb C\{y\}/(y^2)$ forces $T(0)=0$.
Write $T=xW+yU+kV$. Since $k$ is a nonzerodivisor
modulo $(x,h)$, the general solution is
$E=-yU-yV+xE_1+hE_2$. Substitution in the determinant
gives $f\in(p)+I^2$, where
\[
                            p=yh^2(x-h).
\]
The exact identities
$xh^3=ygh^2-yp$ and $h^4=ygh^2-kp$ then put the
equation in the form
\[
 f=A g^2+B gxh+Cgh^2+D(xh)^2+\kappa p,
                       \qquad A(0)\ne0.              \tag{4.49b}
\]

We first treat $\kappa(0)\ne0$.
Use the appended coefficients
$a=Ag+Bxh+Ch^2$, $b=Dxh+\kappa yh$,
$d=-\kappa yh$. The actual conductor identifies
$\operatorname{Hom}(I,\mathcal O_S)$ with the
normalization, generated by $1,n_1,n_2$, where the
adjugate fractions are
\[
          n_1=-a/h,\qquad n_2=-(ak+bh)/h^2.
\]
Since $A$ is a unit, the same actual normalization
algebra is generated by $1,v,w$, where
$v=(n_1+Bx+Ch)/A=-g/h$ and $w=n_2$.
Direct multiplication gives the exact relations
\begin{align*}
 Av^2-Bxv-Chv+Dx^2+\kappa xy-\kappa yh&=0,\\
 xw&=Ayv-Bxy-(\kappa+C)yh,\\
 vw&=\kappa h-Dxy,\\
 w^2+[(\kappa+B)y+Ck]w-A\kappa k+ADy^2
             -D(\kappa+C)yh-B\kappa h&=0.
\end{align*}
At the origin they give $v=w=0$, so there is one
normalization point. The last equation eliminates $k$
analytically, with $k$ in the maximal-ideal square.
Thus $h=y(k-y)$ has quadratic initial term $-y^2$.
Write $a=A_0$, $b=B_0$, $d=D_0$, $u=\kappa_0$.
The three quadratic equations become
\[
 av^2-bxv+dx^2+uxy=0,\qquad
 xw-ayv+bxy=0,\qquad vw+uy^2+dxy=0.
\]
Set $V=av-(b/2)x$, $\delta=ad-b^2/4$,
$Y=\delta x+auy$, $W_0=w+(b/2)y$, and
$W=auW_0+\delta V$. This invertible linear change
turns them into
\[
                       V^2+xY=0,\qquad
                       VW+Y^2=0,\qquad xW-YV=0.
\]
They define the irreducible rational normal cubic cone.
The full tangent-cone ideal contains this prime ideal
and has the same dimension two, so it equals it.
Blowing up the normalization point is smooth along its
exceptional rational normal cubic: the two independent
initial equations have independent derivatives along
that smooth projective curve in every blowup chart.
The exceptional curve is $\mathbb P^1$ with normal
bundle $\mathcal O(-3)$. The exact sequences of its
thickenings have successive terms $\mathcal O(3n)$,
all with $H^1=0$. Formal functions gives
$R^1\rho_*\mathcal O=0$. The singularity is rational.

Finally suppose $\kappa(0)=0$. Perturb $D$ and the
linear part of $\kappa$ generically, and replace
$\kappa$ by $\kappa-\epsilon$. The fixed axes retain
boundary $A_3$: equation (4.49b) stays in the same
ring-condition module, and its boundary transverse
quartic coefficients remain units.
Scale $x=\epsilon^2X$, $y=\epsilon Y$, $z=\epsilon Z$.
Put $K=Y+Z$, $H=YZ$, $G=XK-YH$.
After division by $\epsilon^6$, the first terms are
\[
 A_0G^2+\epsilon(B_0GXH+C_0GH^2)
   +\epsilon^2\bigl[D_0(XH)^2+p(L-1)\bigr].
\]
Near a nonzero smooth point of $G=0$, eliminate the
normal Morse coordinate. On that cone the identities
$XH^3=-Yp$, $H^4=-Kp$ absorb the square-completion
corrections into the coefficients, leaving the leading
tangential function
\[
                    \Phi=D'(XH)^2+p(L'-1)|_{G=0}.
\]
Choose $D'=1$, $L'=-Z$ to exhibit a nonempty open
set of suitable coefficients. The parametrization
\[
 K=r,\quad Y=tr,\quad Z=(1-t)r,\quad
 X=t^2(1-t)r^2
\]
gives $\Phi=t^4(1-t)^4r^7[1+r(t^2-t+1)]$.
At $t=1/2$, $r=-7/6$, its value and Hessian are
\[
 -\frac{823543}{573308928},\qquad
 \operatorname{diag}\left(\frac{117649}{1990656},
                          \frac{15647317}{214990848}\right).
\]
The Hessian determinant is
$1840891197733/427972821516288$, nonzero.
The angular critical numerator factors as
$(2t-1)(3t^2-3t-4)$, so this point is simple and
avoids the axes and $K=0$. The generic coefficient
condition is algebraic and nonempty. After normal
elimination, divide the tangential critical equations
by $\epsilon^2$ and apply the implicit-function theorem.
The resulting full Morse points have nonzero critical
values of order $\epsilon^8$ and tend to the origin.
Boundary transport gives the forbidden square-$-2$
vanishing sphere in a nearby torus fibre. Thus
$\kappa(0)=0$ is excluded, completing both ideal forms.
\end{proof}

All adjugates, colon ideals, normalization relations,
linear coordinate identities, and both polar Hessian
calculations are executed in
\texttt{code/part4/a3\_missing\_coefficient.py}.
The saved output reports PASS.

\begin{statement}{P4.50}{Proved}
In the requested full-fibration classes, no support
with positive-genus ruled resolution has positive
Euler number and generic transverse type $A_3$
on every critical-curve component.
\end{statement}
\begin{proof}
P4.42 and P4.44 leave only a contracted horizontal
positive-genus section and a cycle at one three-branch
point of a conductor with two rational normalized
components. The total defect is nine by P4.33 and
(4.42a). P4.48 excludes planar triples; P4R.16
excludes the nonplanar excess-one tacnode-plus-line
case; the exact excess-two classification and conormal
calculation in P4R.17 exclude
every delta-four nonplanar triple. The branchwise bound
$\delta+3+m/2$ excludes all larger delta invariants,
where $m\ge3$ is total multiplicity.
The only remaining triple is the ordinary three-axis
point. P4.48 gives local defect at least nine, so it
uses the entire budget, and P4.47 and P4.49 give
rational normalization there.

Every additional singular critical-curve point has
strictly positive defect and therefore is absent.
At a smooth point of the entire critical curve,
P4R.14 gives the product tacnode or its $A_1$
normalization degeneration; both have rational
normalization. Away from the conductor, P4R.6
excludes exceptional configurations. The normalization
therefore has rational singularities everywhere.
Its minimal resolution cannot contract the horizontal
positive-genus section. This contradiction proves
the assertion.
\end{proof}


\subsection{Local intersection forms and the quartic tangent obstruction}

\begin{statement}{P4.51}{Proved}
Let $S$ be an irreducible reduced support in the torus
fibration of Part~4.2 with $b_1(S)>0$. Every local
Milnor fibre of a reduced local equation of $S$ has
identically zero rational intersection form. This applies
in particular to the remaining positive ruled support,
whose fundamental group is $\mathbb Z$.
\end{statement}
\begin{proof}
Let $U$ be a sufficiently small full tube and $F_t$ a
nearby compact smooth torus fibre. P4R.4 identifies the
homomorphism
\[
 \pi_1(F_t)=\mathbb Z^4\longrightarrow
 \pi_1(U)=\pi_1(S)
\]
and gives a finite cyclic cokernel, of order equal to
the fibre multiplicity. Its map on rational first
homology is therefore onto. Write
\[
 W=\ker\bigl(H_1(F_t;\mathbb Q)\longrightarrow
                  H_1(S;\mathbb Q)\bigr).
\]
Then $\dim W=4-b_1(S)\le3$.

Choose a coordinate ball about the point, with its
closed ball contained in $U$. The local full fibre
equation is $u g^m$, with $g$ reduced and $u$ a
holomorphic unit. We take a holomorphic $m$-th root of
$u$ on the ball and absorb it into $g$. The actual
nearby fibre in the ball is consequently the disjoint
union of the local Milnor fibres
$g=\zeta t^{1/m}$. Each such fibre $M$ is an oriented
codimension-zero submanifold with boundary of $F_t$.
Its inclusion into $U$ factors through the contractible
ball. Thus every loop of $M$ maps to zero in
$H_1(U;\mathbb Q)$, and the image subgroup
\[
 H=\operatorname{im}\bigl(\pi_1(M)\longrightarrow
                                 \mathbb Z^4\bigr)
\]
satisfies $H\otimes\mathbb Q\subset W$.

The covering criterion lifts $M\to F_t$ to the
connected cover $F_H\to F_t$ corresponding to $H$.
This cover is a $K(H,1)$, with $H$ free abelian. The
map on second rational homology is therefore exactly
\[
 \bigwedge^2(H\otimes\mathbb Q)\longrightarrow
              \bigwedge^2 H_1(F_t;\mathbb Q).
\]
Every class from $H_2(M;\mathbb Q)$ consequently maps
into $\bigwedge^2W$. The torus intersection pairing is
exterior multiplication followed by the orientation
isomorphism $\bigwedge^4H_1(F_t;\mathbb Q)=\mathbb Q$.
It vanishes on $\bigwedge^2W$ because $\dim W\le3$.

This ambient pairing is the local intersection
pairing on $M$. We represent finitely many cycles
inside its interior and make them transverse there;
the orientation-preserving inclusion into $F_t$
retains all intersection points and signs. Thus the
local pairing is zero. The symbolic exterior-square
matrix and its zero Gram matrix are recorded in
\texttt{code/part4/local\_intersection.py} and its
executed output.
\end{proof}

\begin{statement}{P4.52}{Proved}
In the remaining positive ruled support of Part~4.2,
the image of a contracted horizontal pole section has
ambient multiplicity exactly three. This restriction
imposes no generic-transverse-type hypothesis on
any component of the critical curve.
\end{statement}
\begin{proof}
We first give the comparison with a leading
homogeneous polynomial. Write the reduced equation
as $g=g_d+g_{d+1}+\cdots$, and suppose that $g_d$ is
reduced. Every finite collection of cycles in
$M_d=\{g_d=1\}$, together with its intersections,
transports into a local Milnor fibre of $g$.
Indeed, choose a compact neighbourhood of the cycle
representatives. Euler's identity shows that
$dg_d\ne0$ on $M_d$. The functions
\[
 G_\epsilon(x)=\epsilon^{-d}g(\epsilon x)
\]
converge with their derivatives to $g_d$ on that
neighbourhood. We solve
\[
 dG_\epsilon(V_\epsilon)
                     =-\partial_\epsilon G_\epsilon
\]
using the Hermitian gradient divided by its squared
norm. A cutoff supported in a slightly larger
compact neighbourhood, equal to one along the
trajectories under consideration, gives a flow for
small positive real $\epsilon$. It transports those
cycles to $G_\epsilon=1$. Holomorphic scaling by
$\epsilon$ then embeds them in $g=\epsilon^d$ inside
the fixed original ball, preserving orientations
and intersections. This compares smooth nonzero
level sets on compact neighbourhoods; no assertion
about equisingularity of the zero set is needed.
By P4.51 the pairing on $M_d$ is consequently zero.

For a reduced homogeneous polynomial its global
fibre and a sufficiently large compact Milnor fibre
have the same homology and scalar monodromy. One
can see this directly from the squared radius on
the affine fibre. This is a proper semialgebraic
function with finitely many critical values; above
the largest, its gradient supplies a product end.
A large invariant ball therefore gives a compact
core. Rescaling identifies that core with a small
local Milnor fibre. The Milnor-fibre CW dimension
bound gives zero homology above degree two, and
reducedness gives connectedness.

We need one exact variation identity. For such a
homogeneous polynomial, let $\mu_i$ be the
transverse plane-curve Milnor numbers at the
singular points of its reduced projective curve,
and put $M=\sum_i\mu_i$. Write $F$ for its compact
Milnor fibre and $\partial_2F$ for the boundary
part over the singular-link circles. Let
\[
 \begin{split}
 j&:H_2(F;\mathbb Q)\longrightarrow
                         H_2(F,\partial F;\mathbb Q),\\
 j_2&:H_2(F;\mathbb Q)\longrightarrow
                         H_2(F,\partial_2F;\mathbb Q).
 \end{split}
\]
Siersma's chain-level variation identity
\cite[\S3.6, pp.~451--452; \S7.1, p.~462]
{Siersma1991Variation} is
\[
                 j_2(T-I)=\operatorname{VAR}^{II}j.
                                                            \tag{4.52a}
\]
Here $T$ is the Milnor monodromy on $H_2(F;\mathbb Q)$.
The identity does not require the variation map to
be an isomorphism in dimension two. The dimension-two
Wang sequence is explicitly treated in that source,
\S6.3. Poincar\'e--Lefschetz duality identifies $j$
with the intersection form. Its vanishing makes
$\operatorname{im}(T-I)$ lie in $\ker j_2$.
The exact sequence of $(F,\partial_2F)$ identifies
that kernel with the image of $H_2(\partial_2F)$.
Each boundary piece is the mapping torus of a
transverse Milnor fibre. Its Wang sequence gives
\[
 H_2(\partial_2F;\mathbb Q)
       =\bigoplus_i\ker(A_i-I)\subset
                        \bigoplus_iH_1(F_i;\mathbb Q),
\]
where $A_i$ is the vertical monodromy. Thus zero
intersection form implies
\[
                 \operatorname{rank}(T-I)\le M.             \tag{4.52b}
\]
An empty $\partial_2F$ gives the same assertion with
$M=0$.

Suppose now that the reduced projective quartic
$D=\{g_4=0\}$ has a normalization component of
positive genus. By degree, either $D$ is an
irreducible quartic with normalization genus
$g\ge1$, or it is a smooth elliptic cubic and a
line. In the first case
$\sum\delta_i=3-g\le2$, so
$M=\sum(2\delta_i-r_i+1)\le4$.
In the second case B\'ezout gives total intersection
length three and at least one intersection point.
Hence $\sum\delta_i=3$ and
$\sum(r_i-1)\ge1$, giving $M\le5$.

The map $\{g_4=1\}\to\mathbb P^2\setminus D$ is
an unramified fourfold cover. The normalization
sequence and plane-curve genus formula give
\[
              \chi(F)=4\chi(\mathbb P^2\setminus D)
                      =4(7-M).
\]
The three nonidentity scalar monodromies act freely
on the compact invariant core, so their Lefschetz
numbers are zero. It follows by the character
formula for the cyclic group of order four that
\[
 H_0(F;\mathbb C)-H_1(F;\mathbb C)+H_2(F;\mathbb C)
       =(7-M)\,\mathbb C[\mathbb Z/4]
\]
as virtual representations. Every nontrivial
character therefore has multiplicity at least
$7-M$ in $H_2$. The finite-order monodromy is
semisimple, and hence
\[
             \operatorname{rank}(T-I)
                   \ge3(7-M)\ge6>M.                       \tag{4.52c}
\]
This contradicts (4.52b).

We apply this obstruction through the actual pole
comparison. Blow up the ambient origin, normalize
the strict transform, and take its minimal
resolution $T_1$. As in P4.39, this resolution maps
to the original normalization and dominates its
minimal resolution by point blowups. The strict
transform of the positive-genus horizontal section
$H$ retains pole coefficient two. If $E$ is the
ambient exceptional divisor, put
$n=\operatorname{ord}_H(E)>0$, and let $a\ge0$
be the coefficient of $H$ in the total conductor
divisor relative to the strict-transform dualizing
sheaf. Adjunction gives exactly
\[
                          2=(d-2)n+a.                     \tag{4.52d}
\]
Effectiveness follows from the exceptional
intersection matrix argument of P4R.6. An
exceptional positive-genus curve over the new
normalization has strictly positive coefficient,
by its adjunction equation. Thus $d\ge5$ is
impossible. If $d=4$, then $n=1,a=0$; the curve
$H$ is neither exceptional over the new
normalization nor over its conductor. The
normalization is an isomorphism at its generic
point. Its image is a multiplicity-one component
of the quartic exceptional curve, birational to
$H$. It has positive genus. Such a component has
degree three or four, leaving at most one linear
factor; the entire quartic polynomial is reduced.
The preceding contradiction excludes this case.

For clarity, an isolated critical point cannot
hide a contracted $H$. A generic small linear
perturbation of an isolated hypersurface critical
point has a positive finite number of Morse
critical points, the local degree of its gradient.
A small constant shift makes a chosen critical
value nonzero. On a fixed sufficiently small ball
the original zero-level boundary is smooth and
transverse; compactness retains the boundary
submersion over a uniform punctured value disc.
The finite-critical-value and path-transport
argument of P4.20 then carries the Morse
vanishing sphere, of square $-2$, into the
original local Milnor fibre. P4.51 excludes it.
This argument does not assume smoothness of the
normalization.

At a nonisolated point with $d=2$, the Hessian is
nonzero. P4R.15 gives the exact product equation
$X^2+Y^m$. For even $m$ its two analytic branches
are smooth. For odd $m$ its normalization is the
explicit smooth parametrization
$(Y,X,Z)=(u^2,\mathrm i u^m,Z)$.
Neither case has an exceptional curve in its
minimal normalization resolution. A point with
$d=1$ is smooth as well. Therefore the image of
$H$ has multiplicity exactly three.

The cyclic character calculation and all integer
bounds are independently recomputed in
\texttt{code/part4/quartic\_tangent.py}; no numerical
approximation enters the argument.
\end{proof}


\subsection{The cubic double line and its weighted leading fibre}

\begin{statement}{P4.53}{Proved}
In the remaining positive ruled support, blow up the
multiplicity-three image of the contracted horizontal
section $H$, normalize the strict transform $S_1$, and
take its minimal resolution $T_1$. The strict transform
of $H$ is exceptional over a point of $S_1$. The entire
germ of $S_1$ there is normal and has an isolated simple
elliptic singularity. In particular $H$ cannot dominate
a curve in the ambient exceptional plane. No generic
transverse type is imposed on the original critical curve.
\end{statement}
\begin{proof}
Use the comparison maps of P4.52, and let $D_1$ be the
effective total conductor relative to $\omega_{S_1}$.
Write $n=\operatorname{ord}_H(E)>0$ and
$a=\operatorname{coeff}_H(D_1)$. The pole identity is
$2=n+a$, so $(n,a)=(2,0)$ or $(1,1)$.

Suppose $H$ dominates a reduced component $L$ of the
projective cubic tangent curve. At its generic point
a transverse Weierstrass presentation has degree
$m$, the multiplicity of that component in the cubic.
Factoring over the completed valuation field of the
exceptional coordinate $r$ gives the exact degree
identity
\[
 \sum_{H'}\operatorname{ord}_{H'}(r)
                 [\mathbb C(H'):\mathbb C(L)]=m.             \tag{4.53a}
\]
The factors group the roots by valuation branches;
the two factors in each summand are the ramification
index and residue-field degree. Residue extensions
are separable in characteristic zero.

A positive-genus component of a cubic is a smooth
elliptic cubic of multiplicity one. The strict
transform is smooth along it, with $n=1,a=0$,
contrary to $2=n+a$. Every other reduced component
is rational. The finite map from $H$ then has degree
$h\ge2$. Equation (4.53a) excludes $n=2$, since
$m\le3$. Thus $n=a=1$.

At a generic point of $L$, the $h$ points on $H$
give distinct branches of the transverse plane germ,
each with conductor exponent one. The branch formula
\[
             c_i=2\delta_i+\sum_{j\ne i}I_{ij}
\]
forces $h=2$ and exactly two smooth branches meeting
once, with no additional branch. Each has $r$-order
one, so their total Weierstrass degree is two. Hence
$m=2$, and the tangent polynomial is $y^2\ell$,
where $y=0$ is the double line and $\ell$ a different
linear form.

Choose complementary coordinates $x,z$. The
restriction $\ell(x,z)$ is nonzero; any $y$ term in
$\ell$ has higher weight for the weights used below.
In a generic blowup chart along the double line,
write $v$ for the transverse coordinate. The
strict-transform equation begins
\[
                 \ell v^2+r(f_4|_L)+\cdots.
\]
A nonzero $f_4|_L$ would give ramification order two
and residue degree one. Therefore $f_4$ is divisible
by $y$. The two branches meet once, so after
$v=ru$ their distinct leading slopes solve
\[
                         \ell u^2+G_3u+G_5=0,
\]
where $G_3,G_5$ are forms of degrees three and five
in $x,z$. The discriminant
$\Delta=G_3^2-4\ell G_5$ is nonzero and nonsquare
in $\mathbb C(\mathbb P^1)$. The actual residue-field
extension of $H\to L$ is
$\mathbb C(\mathbb P^1)(\sqrt\Delta)$.

Consequently the original equation has weighted
leading polynomial
\[
          P_5=\ell(x,z)y^2+G_3(x,z)y+G_5(x,z),               \tag{4.53b}
\]
of degree five for weights $(1,2,1)$ on $(x,y,z)$,
with no smaller-weight terms. This identifies the
normalized projective component birationally with
$H$ through its function field.

Let $D=\{P_5=0\}\subset\mathbb P(1,2,1)$.
The quadratic is irreducible over $\mathbb C(x/z)$,
since the degree-two residue cover $H$ is connected.
The only possible further factor is a common
vertical factor of $\ell,G_3,G_5$. Since $\ell$
is linear, either $P_5$ is irreducible or
$P_5=\ell Q_4$, with $Q_4$ monic quadratic in $y$.
Both alternatives are reduced.

In the irreducible case, let $k_i$ be the
multiplicities of the distinct zeros of the
homogeneous degree-six discriminant. The normalized
double cover has genus
$(\#\{i:k_i\text{ odd}\}-2)/2$. Positive genus
therefore requires at least four odd multiplicities,
and gives
\[
              \sum_i(k_i-1)=6-\#\{i\}\le2.                  \tag{4.53c}
\]
Away from $\ell=0$ and the weighted vertex,
completing the quadratic identifies the projective
curve germ with $w^2-s^{k_i}$, up to a unit and
holomorphic coordinates. Its transverse Milnor
number is $k_i-1$. At $\ell=0$, if $G_3\ne0$ the
finite root is simple and the other point is at
the weighted vertex. If $G_3=0$, absence of a
common factor gives $G_5\ne0$, and $\Delta$ has
a simple zero. There is no additional singularity.
In the chart $y=1$ the vertex is the quotient of
$\mathbb C^2$ by $(x,z)\mapsto(-x,-z)$; the
lifted equation $\ell+G_3+G_5$ has nonzero
differential. The original polynomial also has
nonzero derivative in the $\ell$ direction on
the $y$-axis, so that axis is not critical.
The total transverse Milnor-number sum is thus
$M\le2$.

In the reducible case write
$G_3=\ell A_2$, $G_5=\ell A_4$. Then
$Q_4=y^2+A_2y+A_4$ has discriminant of degree four.
A connected positive-genus double cover requires
four distinct simple roots. Its normalization is
elliptic and $Q_4$ is smooth on the weighted plane;
it avoids the weighted vertex. The rational
component $\ell=0$ meets it in total length two,
read directly from the restricted monic quadratic.
Either there are two transverse intersections,
with total Milnor number two, or one tangency of
order two, with Milnor number three. There are no
other critical rays. Hence in all cases
\[
                              M\le3.                       \tag{4.53d}
\]

The weighted plane has Euler number three: its
strata $x\ne0$, $x=0,z\ne0$, and $x=z=0$ are
$\mathbb C^2$, $\mathbb C$, and a point.
In the irreducible case $e(D)\le e(H)\le0$.
In the reducible case
$e(D)=0+2-\#(Q_4\cap\{\ell=0\})\le1$.
Therefore
\[
               \chi\bigl(\mathbb P(1,2,1)\setminus D\bigr)
                         \ge2.                            \tag{4.53e}
\]

The weighted scalar action of a primitive fifth
root, with weights $(1,2,1)$, acts freely on
$F=\{P_5=1\}$. Its quotient is
$\mathbb P(1,2,1)\setminus D$: the only possible
nontrivial stabilizer of the weighted
$\mathbb C^*$ action is at the weighted vertex,
which belongs to $D$. This is an unramified
fivefold cover. The compact-core and character
calculation of P4.52 gives
\[
 \operatorname{rank}(T-I)
 \ge4\chi\bigl(\mathbb P(1,2,1)\setminus D\bigr)
 \ge8.
\]
The variation identity (4.52a) applies to every
one-dimensional hypersurface critical locus. Its
boundary mapping-torus estimate would give
$\operatorname{rank}(T-I)\le M\le3$ if the
intersection form were zero.

It is zero here by the actual anisotropic scaling
comparison. Set
\[
 D_\epsilon(x,y,z)=(\epsilon x,\epsilon^2y,\epsilon z).
\]
Then $\epsilon^{-5}f(D_\epsilon(x,y,z))$ converges
with derivatives to $P_5$. The compact-cycle flow
of P4.52 transports its pairings into the actual
local Milnor fibre, whose pairing vanishes by
P4.51. This contradiction excludes every dominating
$H$.

Thus $H$ is exceptional over the new normalization.
Its coefficient $a$ is positive, so $n=a=1$.
Adjunction gives
\[
                    0=2g(H)-2+(D_1-H)\cdot H.
\]
It follows that $H$ is elliptic and meets no other
positive pole component. A zero-coefficient
exceptional curve cannot meet a positive one:
its own intersection equation has nonnegative
canonical term on a minimal resolution and
nonnegative remaining terms, with a strictly
positive term at such a meeting. The connected
exceptional fibre is therefore $H$ alone and is
disjoint from the strict conductor. The
normalization is an isomorphism at its image;
another analytic hypersurface branch would
intersect this branch in a curve and create
conductor meeting $H$. Hence the entire
strict-transform germ is normal, with the single
smooth elliptic exceptional curve $H$. This is
precisely an isolated simple elliptic singularity.

The discriminant multiplicity partitions and the
exact rank bounds are recomputed in
\texttt{code/part4/cubic\_weighted.py} and its
executed output.
\end{proof}


\subsection{Exclusion of the contracted horizontal section}

\begin{statement}{P4.54}{Proved}
In the remaining positive ruled support of Part~4.2,
the horizontal pole section is not contracted by the
normalization resolution. It consequently maps with
degree two or three onto a rational normalized conductor
component whose generic transverse germ is $D_4$.
No transverse-type restriction is imposed on the other
conductor component.
\end{statement}
\begin{proof}
By P4.53, a contracted section would be the sole smooth
elliptic exceptional curve $H$ of an isolated normal
hypersurface singularity on the first ambient strict
transform. Its self-intersection is $-d<0$, and the
actual exceptional coordinate $r$ has order one along
$H$. Write $\pi:\widetilde V\to(V,0)$ for this local
resolution. We establish the following three coordinate
weights and leading degrees, retaining the actual $r$:
\[
 (1,1,1;3),\qquad(1,1,2;4),\qquad(1,2,3;6).
                                                        \tag{4.54w.1}
\]

\paragraph{1. Analytic formal functions and the exact associated graded ring}

Put \(L=\mathcal O_H(-H)\), of degree \(d\), and
\[
 I_m=(\pi_*\mathcal O_{\widetilde V}(-mH))_0
     =\{a\in\mathcal O_{V,0}:\operatorname{ord}_H(a)\ge m\}.
\]
For \(m\ge1\), the restrictions of
\(\mathcal O(-mH)\) to successive infinitesimal neighbourhoods
of \(H\) have graded factors \(L^{m+j}\), \(j\ge0\).
Each has zero first cohomology, by Serre duality on the elliptic
curve and its strictly positive degree. Their extensions also
have zero first cohomology.

The analytic formal-functions theorem is Grauert1960,
Section6, Hauptsatz IIa, printed p.59; Hauptsatz I on the same
page supplies coherence \cite[\S6, Hauptsätze I and IIa]{Grauert60}. The maximal-ideal pullback and the
ideal of \(H\) have the same radical. On a finite cover of
the compact exceptional curve, Noetherianity gives one uniform
power of the latter contained in the former. Their powers are
therefore cofinal. Applying the theorem with
\(\mathcal O(-mH)\) gives a zero completion for
\(R^1\pi_*\mathcal O(-mH)\). This coherent module is
supported at the origin, so it is zero. Consequently
\[
 I_m/I_{m+1}\simeq H^0(H,L^m)\qquad(m\ge1),             \tag{4.54w.2}
\]
through the actual restriction maps in
\[
 0\longrightarrow\mathcal O(-(m+1)H)
 \longrightarrow\mathcal O(-mH)
 \longrightarrow L^m\longrightarrow0.
\]
These maps respect multiplication. Since normality gives
\(\pi_*\mathcal O_{\widetilde V}=\mathcal O_V\), and a
function vanishes at the contracted point exactly when its
restriction to \(H\) vanishes, \(I_1=\mathfrak m\).
The degree-zero piece is \(\mathbb C\). Thus
\[
 \operatorname{gr}_{I}\mathcal O_{V,0}
   =\bigoplus_{m\ge0}H^0(H,L^m)                         \tag{4.54w.3}
\]
as graded algebras, not only as vector spaces.

The same formal-functions theorem in degree zero identifies
\[
 \widehat{\mathcal O}_{V,0}
      \simeq\varprojlim_m H^0(mH,\mathcal O_{mH})
      \simeq\varprojlim_m\mathcal O_{V,0}/I_m.          \tag{4.54w.4}
\]
The second equality uses the just-proved first-direct-image
vanishing. This supplies the exact completeness and separation
needed in the leading-equation argument below; no unproved equivalence between a
divisorial and an ordinary degree filtration is being used.

Because \(\mathfrak m^2\subset I_2\), (4.54w.2) gives a
surjection
\(\mathfrak m/\mathfrak m^2\to H^0(H,L)\).
The source has dimension three and the target dimension \(d\).
Hence \(d\le3\).

\paragraph{2. The three section rings and their actual elliptic curves}

For \(d=3\), Riemann--Roch shows that \(L\) separates
points and tangent vectors: subtracting any length-two
subscheme leaves a degree-one line bundle with one section.
The complete linear series therefore embeds \(H\) as a smooth
plane cubic. The restriction sequence on \(\mathbb P^2\),
using \(H^1(\mathbb P^2,\mathcal O(k))=0\), gives its whole
section ring as
\[
 \mathbb C[x,y,z]/(F_3),\qquad \deg(x,y,z)=(1,1,1),
\]
where \(F_3=0\) is that smooth cubic.

For \(d=2\), the complete pencil is base-point free and gives
a degree-two map \(p:H\to\mathbb P^1\), with
\(L=p^*\mathcal O(1)\). Trace splits
\(p_*\mathcal O_H=\mathcal O\oplus B^{-1}\).
Its Euler characteristic is zero, so \(\deg B=2\).
The algebra structure is therefore a double-cover equation
\(z^2=F_4(x,y)\), where the binary quartic has four distinct
zeros because the double cover is smooth. Projection formula
gives the entire section ring with generator weights
\((1,1,2)\) and that single degree-four relation.

For \(d=1\), write \(L=\mathcal O_H(p)\).
Riemann--Roch gives a degree-one section \(x\), one new
degree-two section \(y\), and one new degree-three section
\(z\). Dividing by powers of \(x\) gives functions of pole
orders two and three at \(p\). Their Weierstrass relation
has weighted degree six. Equivalently, after graded changes
only for describing this section ring, it is
\[
 z^2=y^3+A x^4y+B x^6,
 \qquad 4A^3+27B^2\ne0,
 \qquad\deg(x,y,z)=(1,2,3).
\]
The pole-order count shows that the monomials with at most one
factor \(z\), reduced by this relation, give every space
\(H^0(H,L^m)\). Thus this is the whole section ring. Its
weighted projective curve is the actual smooth \(H\).
The coordinate choice below does not require completing the
square or applying any of these optional graded simplifications.

\paragraph{3. Retaining the given coordinate using only permitted changes}

Start with analytic ambient coordinates \((r,u,v)\), retaining
the given \(r\). Their images form a basis of
\(\mathfrak m/\mathfrak m^2\). The map to
\(I_1/I_2=H^0(L)\) has the rank just computed.

When \(d=3\), it is an isomorphism, so any linear complement
to \(r\) supplies the other two valuation-one coordinates.

When \(d=2\), choose a linear coordinate \(s\) whose
leading section complements that of \(r\), and choose the
remaining independent linear coordinate \(v_0\) in the
kernel of the degree-one restriction. Thus
\(\operatorname{ord}_H(v_0)\ge2\). Its degree-two leading
section must be independent of \(r^2,rs,s^2\). Otherwise
all degree-two initial sections of analytic functions would
lie in that three-dimensional span: a Taylor expansion shows
that every term of ordinary order at least three has valuation
at least three, and the linear terms of valuation two are
accounted for by \(v_0\). This contradicts
\(h^0(L^2)=4\). Therefore \(v_0\) has valuation exactly
two and is the required indecomposable degree-two generator.
The coordinates \((r,s,v_0)\) were obtained by an invertible
linear change.

When \(d=1\), first choose two independent linear coordinates
\(u_0,v_0\) in the kernel of the degree-one restriction.
At least one linear combination, denoted \(s\), has a
degree-two initial section independent of \(r^2\), since
\(h^0(L^2)=2\). Choose the remaining independent linear
coordinate \(v_1\), and subtract a scalar multiple of \(s\)
so that its degree-two initial term is \(\alpha r^2\).
Set
\[
                         v_2=v_1-\alpha r^2.
\]
This is a triangular analytic coordinate change, and
\(\operatorname{ord}_H(v_2)\ge3\). Its degree-three
initial section must be independent of \(r^3,rs\).
Indeed, if it were absent or in that span, any Taylor series
of valuation three would have degree-three initial section
in that same span: the possible terms are precisely the
linear \(v_2\) term and the weighted-degree-three monomials
\(r^3,rs\); its valuation-one and valuation-two terms
must first vanish. This contradicts \(h^0(L^3)=3\).
Thus \(v_2\) has valuation exactly three and supplies the
indecomposable degree-three generator. The three ambient
coordinates remain analytically invertible, because only
an invertible linear change and the displayed quadratic
translation were performed.

This argument proves directly that the chosen analytic
coordinates have the prescribed orders and generate the
graded ring. It avoids a circular appeal to filtered
Nakayama or to a coordinate normalization that might move
the fixed \(r\).

\paragraph{4. The actual weighted leading equation}

Write the defining analytic equation in the coordinates just
constructed, with the weights in (4.54w.1). Its weighted initial
form is the single defining relation of the graded section
ring, up to a nonzero scalar. Here is an explicit proof of that
last assertion.

Let \(F_D\) be the graded relation from the section-ring calculation above, expressed
in the actual chosen initial generators. The analytic function
\(F_D(r,s,v)\) has valuation at least \(D+1\).
At each successive valuation, (4.54w.3) writes its initial
section as a weighted homogeneous polynomial in the three
generators. Subtract that polynomial. This produces a formal
series \(G=F_D+F_{D+1}+\cdots\) whose image in every
\(\mathcal O_{V,0}/I_m\) is zero. Since all generator
weights are at most three, the correction polynomials tend
to zero in ordinary power-series degree; their sum is a
well-defined formal power series. By (4.54w.4), its image in
\(\widehat{\mathcal O}_{V,0}\) is zero.

The completed hypersurface ideal is principal, generated by
the actual analytic equation \(f\). Thus \(G=h f\) for
a formal series \(h\). Taking weighted initial forms gives
\(F_D=\operatorname{in}(h)\operatorname{in}(f)\).
The polynomial \(F_D\) is irreducible because its weighted
projective curve is the smooth connected elliptic curve
\(H\). Also \(f\) is nonconstant and vanishes at the
origin. Hence \(\operatorname{in}(h)\) is a nonzero
constant, and \(\operatorname{in}(f)\) is a scalar multiple
of \(F_D\). This proves (4.54w.1) for the actual defining
equation, with all higher terms retained.

\paragraph{5. The exact descent to the original blowdown coordinates}

In the chart with original coordinates
\(x=r u\), \(y=r v\), every change above has the form
\[
 s=a u+b v+c r,\qquad
 v'=a'u+b'v+c'r-\alpha r^2,
 \qquad ab'-a'b\ne0.
\]
It descends to the actual holomorphic coordinate map
\[
 \begin{aligned}
 x'&=r s=a x+b y+c r^2,\\
 y'&=r v'=a'x+b'y+c'r^2-\alpha r^3,\\
 r'&=r.
 \end{aligned}                                        \tag{4.54w.5}
\]
Its Jacobian determinant at the origin is \(ab'-a'b\),
so the inverse-function theorem supplies an analytic inverse.
This is the required morphism preserving the original
parameter and blowdown structure.


\paragraph{Leading polynomials in the original coordinates.}
Apply this calculation to the simple elliptic point supplied
by P4.53. The original equation is $r^3f_1$, and its original
coordinates are $(x,y,z)=(rs,rv,r)$. The three pairs of
original weights and degrees are
\[
 (2,2,1;6),\qquad (2,3,1;7),\qquad (3,4,1;9).
\]
For a strict-transform monomial $r^as^bv^c$, multiplication
by $r^3$ gives $x^by^cz^{3+a-b-c}$. Its original weighted
degree is three plus its strict-transform weighted degree.
Analyticity of the original equation makes every
negative-power coefficient zero. Thus the actual weighted
leading term is a polynomial of the indicated degree.
On $z\ne0$, the ratios
\[
     \left(x/z^{1+w_s},\ y/z^{1+w_v}\right)
\]
identify its projective component with the strict-transform
elliptic relation on $r\ne0$. This is a birational map
of the actual residue-field curves; their smooth
normalizations are the same elliptic curve $H$.

Anisotropic scaling with these weights and the compact-cycle
flow in P4.52 transports every pairing of the leading
Milnor fibre into the actual local torus fibre. By P4.51,
all these pairings must vanish. We now contradict this
in each of the three cases.

\paragraph{Degree six.}
The leading polynomial is exactly
\[
                    P_6(x,y,z)=G_3(x,y,z^2),
\]
where $G_3(X,Y,Z)$ defines a smooth projective cubic.
The map $(x,y,z)\mapsto(x,y,z^2)$ gives a finite double
cover from $F=\{P_6=1\}$ to $\{G_3=1\}$, branched over
\[
                       B=\{G_3(x,y,0)=1\}.
\]
The homogeneous cubic fibre has Euler number nine: it
is a threefold unramified cover of the complement of
an elliptic cubic in $\mathbb P^2$, whose Euler number
is three. Let $r\in\{1,2,3\}$ be the number of distinct
zeros of the nonzero binary cubic $G_3(x,y,0)$ on
$\mathbb P^1$. The curve $B$ is a threefold cover of
that punctured projective line. Hence
\[
          \chi(B)=3(2-r),\qquad \chi(F)=18-\chi(B)=12+3r.
\]

Weighted monodromy has order six. Its only
nonidentity element with a fixed set is its cube,
which negates $z$ and fixes $x,y$; its fixed set is
$B$. The other four nonidentity elements have no
fixed point on $F$. We choose an invariant
semialgebraic compact core large enough for $F$
and $B$. The finite-order Lefschetz character of
$H_0-H_1+H_2$ is therefore $\chi(F)$ at the identity,
$\chi(B)$ at the cube, and zero elsewhere. Its
trivial-character multiplicity is
\[
                       \frac{\chi(F)+\chi(B)}6=3.
\]
Since $H_0$ has only the trivial character, the
nontrivial-character dimension in $H_2$ is at least
$\chi(F)-3=9+3r\ge12$.

Every nonzero critical orbit of $P_6$ lies on $z=0$
at a tangency of the cubic with $Z=0$. If the
intersection orders are $e_i$, then $\sum e_i=3$.
At a tangency, smoothness of $G_3$ makes its
$Z$-derivative a unit, and the exact transverse
germ is $u^{e_i}+z^2$, after holomorphic coordinates
and a unit. Consequently
\[
                       M=\sum(e_i-1)=3-r\le2.
\]
This also shows reducedness of $P_6$: a repeated
factor would give a two-dimensional critical set,
whereas the derivative calculation gives only
these critical orbits. The variation inequality
(4.52b) contradicts the lower bound twelve for
the nontrivial monodromy rank.

\paragraph{Weighted curve cohomology.}
We record the two curve facts needed for the last
cases, with their quotient comparisons. For
$\mathbb P(a,b,c)=\operatorname{Proj}\mathbb C[u,v,w]$
and a reduced equation of weighted degree $e$,
the exact graded sequence sheafifies to
\[
 0\longrightarrow\mathcal O_{\mathbb P}(-e)
 \longrightarrow\mathcal O_{\mathbb P}
 \longrightarrow\mathcal O_D\longrightarrow0.
\]
The three standard affine charts give its \v Cech
complex. The middle cohomology of every
$\mathcal O_{\mathbb P}(k)$ is zero. The top
cohomology of $\mathcal O_{\mathbb P}(-e)$ has
basis the Laurent monomials $u^{-i}v^{-j}w^{-k}$,
$i,j,k\ge1$, of weighted degree $-e$. Thus
\[
 h^1(\mathcal O_D)=
 \#\{\hbox{nonnegative weighted monomials of degree }
                         e-(a+b+c)\},
 \qquad h^0(\mathcal O_D)=1.                         \tag{4.54a}
\]
In particular the arithmetic genus is zero below
degree $a+b+c$, and is one in degrees $a+b+c$
and $a+b+c+1$ when there is exactly one variable
of weight one.

At a cyclic quotient vertex of prime order two
or three, with no invariant linear ambient
coordinate, a smooth coarse curve has a smooth
plane pullback. To prove this, its normalized
branches form one orbit. A single singular
branch of multiplicity $m\ge2$ has a faithful
normalization action of order $p$; every invariant
ambient function has order at least $2m>p$,
whereas a regular quotient requires a function
of normalization order $p$. For $p$ transitively
permuted branches the stabilizer is trivial,
and every invariant function has order at least
two, whereas a regular quotient requires order
one. Both alternatives are impossible. The
remaining single branch is smooth, and its
reduced equation has the character of a nonzero
linear ambient coordinate. These arguments
use the actual analytic invariant rings:
averaging makes the invariants functor exact,
so every quotient-curve function lifts to an
invariant ambient function.

\paragraph{Degree seven.}
Consider weights $(1,2,3)$ and degree seven.
Formula (4.54a) gives arithmetic genus zero for
curves of degree at most five and arithmetic
genus one in degrees six and seven. Thus a
positive-genus component has degree six or
seven, genus one, and multiplicity one. There
is at most one further degree-one factor;
the polynomial is reduced.

In degree seven the whole curve $D$ is
irreducible with arithmetic and normalization
genus one. Its total delta invariant is zero,
so it is a smooth coarse elliptic curve.
At the quotient vertices of orders two and
three, the preceding argument makes its plane
pullback smooth. Elsewhere the ambient plane
is smooth. Thus the affine cone has no
nonzero critical orbit and $M=0$.

In degree six write $D=E\cup\{z=0\}$, with
$z$ the weight-one coordinate. The curve $E$
is smooth coarse elliptic. It cannot contain
either quotient vertex: its degree-six
equation has trivial character there, while
both local-coordinate characters are
nontrivial, so a smooth lift would require
an impossible invariant linear equation.
The genus formula for $D$ gives total delta
one. Its two smooth components therefore
meet once transversely in the smooth part
of the ambient weighted plane. There is
one ordinary-double critical orbit and
$M=1$.

The weighted plane has Euler number three,
from its affine-cell stratification. In the
two cases $e(D)=0$ or one, so its complement
has Euler number three or two. Both quotient
vertices belong to $D$, since degree seven
is divisible by neither two nor three.
The weighted $\mathbb C^*$ action is free
off $D$, so $\{P_7=1\}$ is an unramified
sevenfold cover of that complement. The
free order-seven monodromy gives nontrivial
$H_2$ rank at least eighteen or twelve,
respectively. This contradicts (4.52b),
whose upper bound is zero or one.

\paragraph{Degree nine.}
For weights $(1,3,4)$, (4.54a) gives
arithmetic genus zero below degree eight
and one in degrees eight and nine.
A positive-genus component must have one
of these two degrees, and multiplicity one.

An elliptic degree-eight component is
impossible. It necessarily contains the
weight-three vertex, since eight is not
divisible by three. The local plane
characters are $1,1$ modulo three, whereas
its equation has character two. The
prime-order quotient argument would give
a smooth lift, whose equation must have
character one. This is a contradiction.

Thus the entire degree-nine curve $D$ is
irreducible and smooth coarse elliptic.
It avoids the weight-three vertex: its
equation there has character zero,
incompatible with a smooth plane lift.
At the weight-four vertex the plane
characters are $1,3$ and its equation
has character one. We check the lift
without assuming that four is prime.
The transitive orbit of normalized
branches has size one, two, or four.

For four branches the stabilizer is
trivial and every invariant maximal-ideal
function has order at least two on a
branch, contradicting regularity of the
quotient. For two branches the stabilizer
has order two. A singular branch gives
invariant orders at least four, exceeding
two. Two smooth branches would have a
quadratic initial equation, but the
quadratic characters of the ambient
action are only zero and two, never one.
For one branch the normalization action
is faithful of order four. Both ambient
coordinate orders are odd, since their
characters are one and three. A singular
branch therefore has multiplicity at
least three, and all invariant functions
have order at least six, exceeding four.
These contradictions leave a single
smooth branch.

The degree-nine cone is therefore smooth
away from its vertex. It is singular
at the vertex, because no variable has
weight nine. Its hypersurface critical
point is isolated, so the isolated-point
Morse transport of P4.52 gives a spherical
class of square $-2$ in its Milnor fibre.
This again contradicts the zero pairing
transported from the actual torus.

All three possibilities for the sole
simple elliptic exceptional curve have
now been excluded. Together with P4.53
this proves that the horizontal pole
section cannot be contracted. P4R.18
then gives the asserted generic $D_4$
image and degree two or three. The
other conductor component retains its
unrestricted transverse type.

The exact weighted cohomology counts,
order-six characters, quotient quadratic
characters and an explicit smooth cubic
are recomputed in
\texttt{code/part4/simple\_elliptic\_weighted.py}
and its executed outputs.

\end{proof}


\subsection{The conductor graph and the horizontal spectral cover}

\begin{statement}{P4.55}{Proved}
In the remaining positive ruled case, the image $A$ of the
horizontal pole section is a smooth rational curve. The map
$H\to A$ has degree exactly two. The third generic
normalization preimage along $A$ is supplied by one rational
vertical conductor component, mapping isomorphically to $A$.
The other image component retains its unrestricted transverse
type. A separated self-identification point on either image
component is impossible.
\end{statement}
\begin{proof}
We write $T\to N\to S$ for the minimal normalization
resolution. By P4.54 and P4R.18, the pole section $H$ is
noncontracted, has genus $g\ge1$, and maps with degree two
or three to $A^\nu=\mathbb P^1$. The full generic transverse
germ along $A$ is $D_4$. The reduced image conductor is
$C=A\cup B$, with two rational normalized components and
one graph cycle.

The reduced normalization conductor $D$ has a tree as its
point-identification graph. Its positive-genus component
$H$ is smoothly embedded: the ruling descends to $N$, and
restriction of that ruling is an inverse to $H\to N$.
In particular no two points of $H$ are identified by the
normalization resolution. Every other normalized conductor
component is rational. These are the actual conductor
curves; vertical components and their identifications are
retained throughout the argument.

At any point of $C$, points of $N$ over it correspond to
local irreducible hypersurface branches of $S$. Two such
branches intersect along a curve, by the height bound for
two equations in a regular local ring of dimension three.
Their nonempty supports among the local analytic branches
of $C$ therefore pairwise intersect.

\paragraph{Separated self-identifications.}
Suppose $p\in B\setminus A$ has two analytic branches.
Let $D_B$ be the union of the conductor components mapping
to $B$, with $c_B$ normalized components and $s_B\ge1$
connected components. Its graph is a forest. Every component
maps finitely and surjectively to $B^\nu$, so the numbers
$r_+,r_-$ of its normalized points over the two preimages
of $p$ satisfy $r_+,r_-\ge c_B$.

There are $n$ normalization points of $S$ over $p$.
Their nonempty supports in $\{+,-\}$ pairwise intersect.
They have a common element: a singleton forces its element
into every support, and without a singleton all supports
are the entire two-element set. After interchanging signs,
$r_+\ge n$. The graph contribution at $p$ is consequently
\[
                  r_++r_--n\ge c_B.
\]
The whole forest has contribution $c_B-s_B<c_B$, a
contradiction. The same proof applies to a self-identification
of $A$ away from $B$, using all components mapping to $A$.
Thus the image cycle consists either of two distinct
intersections of $A$ and $B$, or of one three-branch point
belonging to both components.

\paragraph{The actual primary conductor and smoothness of $A$.}
Let $C_{A,\mathrm{sch}}$ be the primary component of the
actual image conductor supported on $A$. The primary
restriction maps of P4R.13, used in P4R.18, give
\[
                    g-1\le\chi(\mathcal O_{C_{A,\mathrm{sch}}}).
                                                            \tag{4.55a}
\]
The generic ordinary-triple conductor is the transverse
maximal ideal squared. Purity of the actual conductor
quotient identifies this primary ideal with $P_A^{(2)}$,
as in P4.31. The exact normalized conormal calculation of
P4.34 therefore gives
\[
 \chi(\mathcal O_{C_{A,\mathrm{sch}}})
       =5-\sum_p\bigl(\delta(A,p)+\tau_p-r_p\bigr).
                                                            \tag{4.55b}
\]
A planar point contributes exactly $5\delta(A,p)$, and a
nonplanar singular point with at most two branches contributes
at least seven. There cannot be three branches of $A$ at one
point: the identity
$\sum_p(r_p(C)-1)=2$ would then leave no identification for
the necessary connection with $B$; placing that connection
at the same point gives at least four branches and also
exceeds the identity.

Since $g\ge1$, (4.55a)--(4.55b) bound the total defect by
five. If $A$ is singular, it is planar everywhere, has total
delta invariant one, and $g=1$. It is then a locally complete
intersection curve in $Y$. Its conormal
$Q=P_A/P_A^2$ is locally free of rank two. Let
$j:\mathbb P^1\to A$ be its normalization and put $E_0=j^*Q$.
Adjunction gives
\[
             \omega_A=(\omega_Y|_A)\otimes(\det Q)^{-1}.
\]
The degree of $j^*\omega_Y$ is zero because $b_2(Y)=0$.
Also $p_a(A)=1$, so $\deg j^*\omega_A=0$. The degree of a
line bundle is preserved by finite birational normalization:
the normalization sequences for that line bundle and for
$\mathcal O_A$ have the same finite-length correction.
Consequently
\[
                              \deg E_0=0.                 \tag{4.55c}
\]

The actual reduced equation has generic order three along
$A$. Powers of a complete-intersection curve ideal are
unmixed, so $P_A^{(3)}=P_A^3$. Its third normal term is an
actual global section
\[
 q_3\in H^0\!\left(A,
          \operatorname{Sym}^3 Q\otimes\mathcal O_Y(S)|_A\right).
\]
The line twist has degree zero and becomes trivial on
$\mathbb P^1$. On the quotient projective bundle
$\mathbb P(E_0)$, with tautological class $\xi$ and fibre
class $F$, the pulled-back cubic has zero divisor of class
$3\xi$. Remove its common zero fibres, of total degree
$b\ge0$. The residual spectral curve $\Gamma$ has class
$3\xi-bF$. Its horizontal components are reduced, because
the generic $D_4$ cubic has three distinct roots.

We specify the spectral morphisms. At a generic point of
$H$, the first normal derivative of $T\to S\subset Y$
gives the direction of its smooth transverse branch and
hence a point of $\mathbb P(E_0)$. On the smooth compact
curve $H$, its two local coordinates extend after cancelling
their common power. This gives a holomorphic map over the
actual map $H\to\mathbb P^1$. Distinct generic preimages
on $H$ give distinct cubic roots. Hence its map onto the
spectral component $\Gamma_H$ is finite and birational,
and $H$ is that component's normalization. In degree three,
$\Gamma=\Gamma_H$. In degree two, precisely one remaining
rational vertical conductor component supplies the third
root, with degree one over $A$. Its direction curve
$\Gamma_R$ is a section and is isomorphic to $\mathbb P^1$.
Thus $\Gamma=\Gamma_H\cup\Gamma_R$. These exhaust the
three generic roots.

For a rank-two bundle of degree $e$ on $\mathbb P^1$,
\[
 \xi^2=e,\qquad \xi F=1,\qquad F^2=0,\qquad
 K_{\mathbb P(E_0)}=-2\xi+(e-2)F.
\]
Adjunction gives the exact arithmetic genus
\[
                    p_a(3\xi-bF)=3e-2-2b.               \tag{4.55d}
\]
With (4.55c) this is $-2-2b$. In degree three it is at
least $g\ge1$. In degree two the exact union sequence gives
\[
 p_a(\Gamma)=p_a(\Gamma_H)+p_a(\Gamma_R)
           +\operatorname{length}(\Gamma_H\cap\Gamma_R)-1
       \ge g-1\ge0.
\]
Both alternatives are impossible. Therefore $A$ is smooth.

\paragraph{The three-branch image point in degree three.}
Suppose now that $H\to A$ has degree three. Then $H$ is
the only conductor component mapping onto $A$. Since $A$
is smooth, a single three-branch image point consists of
one $A$ branch and two $B$ branches, denoted $+$ and $-$.

Let $c_B$ be the number of normalized conductor components
mapping to $B$. The entire conductor has $1+c_B$
normalized components and a tree graph, so its total
point-identification contribution is $c_B$. At the proposed
point write $r_A\ge1$, $r_+,r_-\ge c_B$ for the numbers
of conductor branches, and $n$ for the number of normalization
points. Each normalization point has at most one $A$
conductor branch, because $H$ is embedded.

A singleton support $\{A\}$ forces every point to carry
$A$. Then $n=r_A$, and the graph contribution is
$r_++r_-\ge2c_B>c_B$. A singleton $\{+\}$ forces
$r_+\ge n$, giving contribution at least
$r_A+r_-\ge1+c_B$; the same applies to $\{-\}$.
Thus every support has at least two elements. Consequently
$r_A+r_++r_-\ge2n$, and its graph contribution is at least
\[
       \frac{r_A+r_++r_-}{2}
                         \ge c_B+\frac{r_A}{2}>c_B.
\]
This is again impossible. Only two distinct $A$--$B$
junctions remain in degree three.

\paragraph{Both degree-three junctions are common cubic zeros.}
Let these junctions be $p,q$. If a local normalization
branch over $p$ has conductor only over $A$, pairwise
hypersurface intersections force every normalization point
over $p$ to carry $A$. Thus $n=r_A$, and the graph
contribution at $p$ is $r_B\ge c_B$, using the entire
budget. At $q$ some normalization point must carry both
$A$ and $B$: otherwise an $A$-only point and a $B$-only
point have disjoint supports. That point contributes
positively to the graph, contradicting the budget. There
is therefore no $A$-only local branch at either junction.

We retain all factors in the following local calculation.
Write the smooth curve as $A=\{x=y=0\}$. If its nonzero
cubic normal form has a simple root, use the normal blowup
chart $x=yu$. The implicit-function theorem on the strict
equation gives that root as $u=\alpha(y,z)$. The actual
function has the analytic factor
\[
                         f_1=x-y\alpha(y,z).
\]
This branch is smooth. Exact division writes $f=f_1f_2$,
and restriction of $f_2$ to this branch is $y^2$ times a
unit: the other cubic directions differ from the simple
root at the point. Its intersections with all other
branches are therefore supported only on $A$. It is the
forbidden $A$-only normalization branch.

Every nonzero binary cubic has a simple root unless it is
a cube. P4R.19 excludes a nonzero cube by its exact finite
weighted maps and compact intersection class. Hence the
whole cubic tensor vanishes at both $p$ and $q$.

\paragraph{The degree-three genus contradiction.}
For the now smooth curve $A=\mathbb P^1$, adjunction gives
$\deg N_{A/Y}=-2$, so $E=N_{A/Y}^*$ has degree two.
The actual cubic has at least two common zeros. Its residual
spectral divisor is $3\xi-bF$ with $b\ge2$, and is
irreducible with normalization $H$. Formula (4.55d) gives
\[
                    p_a(\Gamma)=4-2b\le0,
\]
contradicting $p_a(\Gamma)\ge g\ge1$. Degree three is
excluded. The degree-two alternative and its unique rational
third sheet are precisely the possibilities already identified
by the spectral morphism above.

The exact intersection polynomial and the graph inequalities
are recomputed in
\path{code/part4/d4_graph_spectral.py}; its saved output
reports PASS. The inequalities in this proof establish the
unrestricted graph counts.
\end{proof}


\subsection{The other image curve with at most two generic branches}

\begin{statement}{P4.56}{Proved}
In the remaining case of P4.55, the other conductor image
component $B$ has at least three generic normalization
branches. This conclusion imposes no preliminary bound on
its transverse multiplicity, conductor exponent, or tangency
order.
\end{statement}
\begin{proof}
We retain the smooth curve $A$, its degree-two horizontal
cover $H$, and its degree-one vertical cover $R$. All other
normalized conductor components map to $B$. Let $r_B$ be
the generic normalization branch number along $B$, and
$c_B$ the number of global conductor components over it.
Their finite degrees sum to $r_B$. We exclude $r_B=1,2$.
The normalized conductor has $2+c_B$ components and a tree
graph, with total identification contribution $1+c_B$.
Distinct points on one normalized component cannot be
identified, since that would create a graph loop.

\paragraph{A length-two intersection calculation.}
Let $F,G$ be two local hypersurface branches whose
scheme-theoretic intersection is supported on the smooth
curve $A$ and has generic length $m\in\{1,2\}$ along it.
The ring
\[
                 Q=\mathbb C\{x,y,z\}/(F,G)
\]
is a one-dimensional complete intersection. Its only
associated prime is the ideal of $A$. Choose $z$ as a
coordinate on $A$. Then $Q$ is finite and torsion-free
over $\mathbb C\{z\}$, hence free of rank $m$. The
reduction map to $\mathbb C\{z\}$ splits through the
chosen coordinate.

For $m=1$, the ring is $\mathbb C\{z\}$, and the linear
parts of $F,G$ have rank two. Both branches are smooth.
For $m=2$, the nilradical is a free rank-one module over
$\mathbb C\{z\}$. Its square vanishes generically and
therefore vanishes, by torsion-freeness. A generator
$\epsilon$ gives the exact analytic algebra isomorphism
\[
                    Q\simeq\mathbb C\{z,\epsilon\}/(\epsilon^2).
                                                            \tag{4.56a}
\]
Thus $Q$ has embedding dimension two, so one of $dF,dG$
is nonzero. If $G$ has generic normal order two along $A$,
it belongs to the ordinary square of its smooth ideal and
$dG$ vanishes on all of $A$. Then $dF\ne0$ and $F$ is
smooth. We use this calculation on the original intersection
scheme, retaining its multiplicity.

\paragraph{Two distinct image junctions.}
At a two-branch image point, let $r_A,r'_B$ count the
normalized conductor points over its two branches, and
let $n$ count normalization points of $S$. Pairwise
intersection of the hypersurface branches gives a common
element in their supports in $\{A,B\}$. Hence
$n\le\max(r_A,r'_B)$, and the graph contribution is at
least $\min(r_A,r'_B)$. Here $r_A=2$ or three, from the
one or two $H$ points and the unique $R$ point.

If $c_B=2$, both degrees are one, and $r'_B=2$ at either
junction. The two contributions total at least four,
exceeding the budget three. If $c_B=1$ and its degree is
two, an unramified junction costs at least two and the
other at least one, exceeding the budget two. Thus one
component $V$ remains, with exactly one normalized point
over each junction.

If $B$ belongs to every branch support at a junction,
there is only one normalization point there and its
contribution is $r_A\ge2$. The other junction makes this
impossible. Therefore $A$ belongs to every support. The
budget forces contribution exactly one at each junction:
every normalization point has one $A$ branch, and the
sole $V$ point is paired with one of them. Connectedness
forces the graph path $H$--$V$--$R$, with its edges at
the two different junctions.

At the $R$--$V$ junction, let $F$ be the corresponding
hypersurface branch. The other branches come from $H$
and carry only $A$. If its two preimages are distinct,
intersection with either such branch has support $A$ and
generic length one, because the full generic $A$ germ
is an ordinary triple point. The preceding calculation
makes $F$ smooth. If the two $H$ preimages coalesce,
their one hypersurface branch $G$ has generic normal
order two along $A$, and the intersection with $F$ has
length two. Formula (4.56a) again makes $F$ smooth.

This is impossible. The branch $F$ carries all the nearby
generic $B$ preimages: either the single branch of a
singular unibranch transverse germ, or both branches
when $V\to B^\nu$ has degree two and ramifies at this
point. In either case $F$ is intrinsically nonnormal
along the nearby $B$ germ. It cannot be smooth in a
neighbourhood of the junction.

\paragraph{A three-branch image point with one $B$ component.}
The image branches are $A,B_+,B_-$, since $A$ is smooth.
If $c_B=1$, the normalized points on $V$ over the two
$B$ signs are distinct and cannot be identified in $N$.
No normalization point can therefore carry both signs.
Pairwise hypersurface intersections force every point
carrying either sign to carry $A$ as well.

The graph budget two now forces exactly one point over
each sign, each paired with exactly one $A$ branch, and
no other identification. The pairs are $H$--$V$ and
$R$--$V$. An extra $H$ preimage is $A$-only. At the
$R$--$V$ point, the other $H$ branch or branches intersect
its hypersurface branch only along $A$, because their
$B$ supports have the opposite sign or are empty. The
same length-one or length-two calculation makes this
branch smooth, contradicting its intrinsic nonnormality
along the $B$ germ whose generic preimages it contains.

\paragraph{A three-branch point with two $B$ components.}
Write the two degree-one components as $V,W$. If $h=1$
or two is the number of $H$ points over the $A$ value,
the total normalized branch count is
\[
                         T=(h+1)+2+2=h+5.
\]
The graph budget three gives $n\ge h+2$ normalization
points. No singleton image support is possible: a
singleton $\{A\}$ forces $n\le h+1$, and a singleton
$\{B_+\}$ or $\{B_-\}$ forces $n\le2$. Hence every
point has at least two image labels and $n\le T/2$.
This excludes $h=2$. For $h=1$ there are precisely
three points, each with exactly two conductor branches
of different image labels. Up to signs and interchange
of $V,W$, their supports are
\[
 \begin{array}{c|c}
 \text{hypersurface branch}&\text{conductor branches}\\ \hline
 F_R&R,\ B_-\text{ on }V\\
 F_H&H,\ B_+\text{ on }W\\
 F_0&B_+\text{ on }V,\ B_-\text{ on }W.
 \end{array}
                                                            \tag{4.56b}
\]
The tree is the path $H$--$W$--$V$--$R$, and $H\to A$
is ramified with local degree two at its point.

The intersection of $F_R$ and $F_H$ is supported on $A$
and has generic length two. The equation of $F_H$ has
normal order two along $A$, so (4.56a) makes $F_R$
smooth. Choose its equation as $X$, a coordinate $Y$
on that surface cutting $A$, and a coordinate $z$ on
$A$. Preparation and depression in $Y$, preserving the
ideal $(X,Y)$, give the actual equation
\[
                       F_H=\text{unit}\,[Y^2+X^2k(X,z)].
                                                            \tag{4.56c}
\]
The second $X$ follows from normal multiplicity two
along every nearby point of $A$.

The monic finite partial normalization is
\[
      N'=\{u^2+k(X,z)=0\},\qquad Y=Xu.
                                                            \tag{4.56d}
\]
It is an isomorphism away from $A$. The branch $F_H$
is generically smooth along $B_+$: the two generic
$B$ sheets belong separately to $F_H$ and $F_0$.
Thus $N'$ has at most an isolated singularity here.
Its hypersurface Cohen--Macaulay property and regularity
in codimension one make it normal, so it is the actual
normalization.

The reduced divisor $X=0$ is supported on $H$ and has
equation $u^2+k(0,z)=0$. The actual $H$ is smoothly
embedded and has degree two with ramification over $A$.
An even zero of $k(0,z)$ gives two branches; an odd
zero of order at least three gives a singular curve.
Therefore its zero is simple, $k_z(0,0)\ne0$, and
$N'$ is smooth. The minimal normalization resolution
is the identity in this neighbourhood.

Let $m\ge1$ be the generic intersection multiplicity of
the two branches along $B$. It is the same on both
local signs, because they belong to the same irreducible
image component. Let $G$ define $F_0$, and set
\[
                      b=\operatorname{ord}_{p,A}(G|_A)>0.
\]
On the smooth branch $F_R$, its divisor is $mB_-$.
Restricting to $A$ gives
\[
                 b=m I_{F_R}(A,B_-)\ge m.                \tag{4.56e}
\]
On the smooth normalization (4.56d), the divisor of $G$
is $mW$. The actual pole support is normal crossing,
so $H$ and $W$ meet transversely. Its restriction to
$H$ consequently has order $m$. But it is the pullback
of the same ambient restriction $G|_A$ under the local
degree-two map. Hence
\[
                              2b=m,
\]
contradicting (4.56e). Every case with $r_B\le2$ is
excluded.

The exact incidence partitions and the normalization
substitution are checked in
\texttt{code/part4/d4\_other\_two\_branches.py}.
Its corrected run reports PASS; the initial structural
symbolic-equality assertion failure is preserved separately.
\end{proof}


\subsection{Exclusion of the two-junction spectral configuration}

\begin{statement}{P4.57}{Proved}
The remaining degree-two configuration of P4.55--P4.56
cannot have its image-conductor graph cycle realized at
two distinct intersections of its two components.
\end{statement}
\begin{proof}
The component $B$ has at least three generic normalization
branches, hence transverse multiplicity at least three.
Along the smooth $A$, choose coordinates $A=\{x=y=0\}$,
with $z$ on $A$. The actual equation belongs to $(x,y)^3$;
write $q_3(x,y;z)$ for its normal cubic.

A local branch of $B$ has first nonzero normal displacement
$(a,b)t^q$, with $(a,b)\ne0$. All ambient Hessian entries
of the actual equation vanish identically along $B$, by
its generic multiplicity. The first normal-order term of
its normal Hessian is
\[
       t^q\operatorname{Hess}_{x,y}
                    (q_3(\,\cdot\,,\,\cdot\,;0))(a,b).
\]
Terms of higher normal order and terms from varying $z$
have larger order in $t$. This matrix is therefore zero.
For a homogeneous binary cubic $Q$, a nonzero vector
annihilating this Hessian makes the directional derivative
of $Q$ the zero quadratic. In coordinates with that vector
as first axis, $Q$ is a cube of the other variable, or zero.
P4R.19 excludes the nonzero cube. Thus the cubic vanishes
at every $A$--$B$ junction, including tangent junctions.

\paragraph{The global spectral equality.}
Let $p,q$ be the two distinct junctions. The common cubic
zero divisor has degree $b\ge2$. On $\mathbb P(N_{A/Y}^*)$
its residual spectral curve is $\Gamma_H\cup\Gamma_R$,
of class $3\xi-bF$. The bundle has degree two, $\Gamma_R$
is a section, and $H$ normalizes $\Gamma_H$. Adjunction
and the exact union sequence give
\[
       g(H)+\delta(\Gamma_H)+\Gamma_H\cdot\Gamma_R=5-2b.
\]
More precisely, the displayed formula is an equality of
integers with $\delta(\Gamma_H)$ the total normalization
length. Since $g(H)\ge1$, it forces
\[
 b=2,\qquad g(H)=1,\qquad \delta(\Gamma_H)=0,
                  \qquad\Gamma_H\cdot\Gamma_R=0.          \tag{4.57a}
\]
The two common cubic zeros are simple, the double spectral
cover is smooth, and it is disjoint from the third section.

\paragraph{The graph support at a junction.}
Let $c_B$ count the global conductor components over $B$.
The graph budget is $c_B+1$. At either junction the
contribution is at least $\min(r_A,r'_B)$, where
$r_A\ge2$ and $r'_B\ge c_B$ count normalized points
over its two image branches.

If $c_B=2$, the two contributions total at least four,
exceeding the budget three. If $c_B=1$, a junction with
at least two $B$ preimages costs at least two, and the
other at least one, exceeding the budget two. Thus both
have one $B$ preimage. The argument of P4.56 gives the
path $H$--$V$--$R$, with one edge at each junction.
At the $R$--$V$ junction the local branch carries all
nearby generic $B$ sheets of $V$, regardless of its
finite degree. Its intersection with the $H$ branch or
branches is a length-one or length-two thickening of
smooth $A$. P4.56 makes it smooth, contradicting its
intrinsic nonnormality from the at least three $B$ sheets.

Therefore $c_B\ge3$. If $A$ belongs to every branch
support at a junction, that point costs at least
$r'_B\ge c_B$, and the other costs at least two.
This exceeds $c_B+1$. The two-element support argument
consequently forces $B$ to belong to every local
hypersurface branch at each junction. Branches carrying
only $B$ have been retained in this conclusion.

\paragraph{Choosing distinct normalization points.}
The ruling descends to $r_N:N\to H$. Its restriction
to $H$ is the identity, and $R$ maps to a fixed point
$b_0\in H$. If a point of $H$ and a point of $R$ have
the same image in $N$, that $H$ point is $b_0$.
This includes identifications through a contracted
exceptional chain. All such coincidences therefore lie
over the single value $a_0\in A$ which is the image of
$b_0$. Choose a junction different from $a_0$, and call
it $p$. Its $H$ preimages and its $R$ preimage are
distinct normalization points. Different $H$ points
are already distinct because $H$ is embedded.

The actual local equation factors into the corresponding
irreducible hypersurface factors and any factors carrying
only $B$. For an $A$-containing factor, let $b_i\ge0$
be the common $z$-order of its first nonzero normal term.
For a factor not containing $A$, let $c_j\ge1$ be the
order of its nonzero restriction to $A$ at $p$.
Gauss's lemma over $\mathbb C\{z\}$ gives
\[
                         \sum_i b_i+\sum_j c_j=1.         \tag{4.57b}
\]
Indeed primitive homogeneous polynomials have nonzero
product modulo $z$, so their common coefficient orders
add exactly. The right side is the simple cubic zero
of (4.57a).

If $H$ is unramified here, there are three
$A$-containing factors of normal order one. At least
two have nonzero linear normal terms, by (4.57b).
Their directions differ by (4.57a). The implicit-function
theorem makes their intersection exactly $A$, whereas
both must contain $B$. This is a contradiction.
Thus $H$ is ramified at $p$.

There are now two $A$-containing factors, denoted by
$F_H$ and $F_R$, of normal orders two and one. Their
residual spectral directions differ. Equality (4.57b)
has precisely the following three possibilities.

\paragraph{A factor carrying only $B$ uses the zero.}
Here $b_H=b_R=0$ and one such factor restricts to $A$
with order one. The branch $F_R$ is smooth. In coordinates
$X=F_R$, $A=\{X=Y=0\}$, restriction of $F_H$ to $X=0$
is $Y^2$ times a unit, since its residual direction
differs from the $R$ direction. Thus these two branches
intersect only along $A$. They both contain $B$, which
is impossible.

\paragraph{The quadratic factor uses the zero.}
Here $b_H=1$, $b_R=0$, and there is no other factor.
The branch $F_R$ is again smooth, and in its coordinates
exact division gives
\[
                         F_H|_{X=0}=Y^2h(Y,z).
\]
The value $h(0,z)$ has a simple zero: the residual
quadratic is nonzero on the $R$ direction, so its
common scalar factor survives with order one on this
restriction. Thus $h_z(0,0)\ne0$. Its zero curve is the
extra image branch $B$, with intersection multiplicity
one between $F_H$ and $F_R$.

The smooth $F_R$ branch supplies one generic $B$
normalization branch. Each branch supplied by $F_H$
has positive intersection multiplicity with it. There
can therefore be at most one such branch. The full
$B$ germ would have at most two branches, contradicting
P4.56.

\paragraph{The linear factor uses the zero.}
Finally $b_R=1$, $b_H=0$, with no other factor.
The nonzero quadratic form of $F_H$ has a double root.
Preparation and depression preserving $A$ give
\[
                   F_H=\text{unit}\,[Y^2+X^2k(X,z)].
\]
Its spectral curve is $u^2+k(0,z)=0$. By (4.57a) this
is smooth and its double cover is ramified at $p$.
Hence $k_z(0,0)\ne0$. The exact finite map
\[
 \{u^2+k(X,z)=0\}\longrightarrow\{F_H=0\},
                         \qquad Y=Xu
\]
has smooth source, is birational, and is an isomorphism
off $A$. It is the normalization. Thus $F_H$ supplies
one generic $B$ branch.

For $F_R$, the simple common zero of its linear normal
term makes its quadratic Taylor part
\[
                       zL(X,Y)+Q_2(X,Y),\qquad L\ne0.
\]
There is no $z^2$ term because $F_R$ vanishes on $A$.
Its Hessian has rank at least two. Rank three gives an
isolated critical point. In rank two the holomorphic
Morse lemma with a parameter gives the exact equation
$U^2+V^2+h(W)$. A nonisolated critical curve forces
$h'\equiv0$, hence $h=0$ because the factor vanishes
at the point. The resulting factorization
$(U+iV)(U-iV)$ contradicts local irreducibility.
Therefore $F_R$ has at most isolated singularities and
also supplies one generic $B$ branch. Again the full
branch number is two, contradicting P4.56.

All allocations in (4.57b) have been excluded, proving
the assertion. The exact genus budget and nonzero
Hessian minor are recomputed in
\texttt{code/part4/d4\_two\_junction.py}; its saved
output reports PASS.
\end{proof}


\subsection{The completed irreducible classification and sign theorem}

\begin{statement}{P4.58}{Proved}
Let $f:Y\to B$ be the fibration of Part~4.2, with
$b_1(Y)=b_2(Y)=0$, connected fibres and general fibre a
complex two-torus. Suppose every reduced fibre support
is irreducible or has normal crossings. Then
\[
 B=\mathbb P^1,\qquad c_3(Y)\in\{0,2\},
                       \qquad b_3(Y)=2-c_3(Y).
\]
Every reduced support has nonnegative Euler number.
An irreducible singular support whose normalization has
a rational resolution has Euler number two. Every other
irreducible support has Euler number zero. Such a support
is either smooth with a finite unramified torus cover,
or singular with smooth minimal elliptic ruled
normalization, or singular with smooth minimal Hopf
normalization. In particular the class-VII normalization
alternative is realized only by a smooth minimal Hopf
surface.

For a normal-crossing support the Euler number equals
its number of triple points, and is even. Its zero case
is exactly P4.8: in the reducible case the normalization
is a disjoint union of smooth minimal Hopf or non-K\"ahler
elliptic surfaces, with smooth elliptic conductor curves
and finite unramified branch maps. The equality
$c_3(Y)=0$ holds exactly when every support has one of
the corresponding zero-Euler forms.
\end{statement}
\begin{proof}
P4.26 proves the sign for every irreducible support,
including singular normalizations, and P4.6--P4.8 give
the full normal-crossing calculation. P4.35 proves the
exact value two for the rationally resolved singular
case. P4.36--P4.37 show that every class-VII case has
smooth minimal Hopf normalization and Euler number
zero. P4.32 gives the complete irreducible zero-Euler
classification and its converses.

We finish the positive-genus ruled alternative without
a transverse-type assumption on its vertical conductor
branches. P4R.13 and the preceding reductions show that
a positive-Euler ruled support would have a doubled
horizontal pole section and an image conductor with two
rational normalized components and one graph cycle.
P4.54 excludes contraction of that section, using the
actual weighted normalization and blowdown maps.
P4R.18 reduces its generic image to $D_4$ with degree
two or three. P4.55 makes that image smooth and excludes
degree three, leaving the degree-two section and one
rational vertical third sheet. P4.56 forces at least
three generic branches on the other image component.

The graph alternatives are exhaustive by P4.55. P4.57
excludes two distinct junctions. P4R.20 excludes a sole
three-branch image point, with the exact conductor-path
and endpoint normalization calculation. Thus no
positive-Euler ruled support remains. This proves the
asserted value zero for every nonrational irreducible
support, and P4.32 supplies its actual normalization.

Finally Euler additivity gives $c_3(Y)$ as the sum of
these nonnegative support Euler numbers. The base,
Poincar\'e-duality calculation $e(Y)=2-b_3(Y)$, and
parity are proved in P4.26. They yield the displayed
global values. A sum of nonnegative integers vanishes
exactly when every summand vanishes, giving the equality
classification through P4.32 and P4.8.

The hypothesis checks are the explicit constructions
P4.10 and P4.11. Removing $b_2=0$ permits the proved
Euler value $24$. Removing the torus-fibre hypothesis
permits the Hopf-fibre example, which fails the stated
normalization equality classification. The numerical
normal-crossing sign itself follows from the stronger
calculation in P4.7 and does not require that hypothesis.
\end{proof}

